\input _template

\let\angle\sphericalangle

\setlanguage{CS}

\Title{Přidávání bodů do obrázku}

\MathcompsLink{pridavani-bodu}

\Author{Patrik Bak}

\sec Úvod

Všichni známe základní metody řešení olympiádní geometrie, jako jsou

\begitems
\i počítání úhlů a tětivové čtyřúhelníky,
\i shodnost a podobnost trojúhelníků,
\i mocnost bodu a chordály.
\enditems

Jen kombinací těchto metod dokážeme vyřešit spoustu úloh. Ne vždy je však snadné přijít na to, kterou techniku použít. Obecné techniky, které nám pomáhají, jsou

\begitems
\i přeformulování zadání,
\i rozpoznání známých dílčích tvrzení,
\i eliminace bodů,
\i přidávání bodů do obrázku.
\enditems

V tomto materiálu se zaměříme na to poslední a procvičíme si některé konkrétní situace, které se opakují v úlohách všech obtížností.

\sec Obecné tipy

\secc Rýsování obrázků

Jednou z věcí, které mohou rozhodnout o tom, zda úlohu vyřešíte, nebo ne, je dobrý obrázek. Pár tipů:

\begitems
\i rýsujte velké obrázky,
\i nerýsujte různostranné trojúhelníky jako rovnostranné,
\i používejte rýsovací pomůcky,
\i používejte dobře ořezané tužky a gumu,
\i používejte barevné tužky (např. k označení úseček stejné délky),
\i nerýsujte kružnice, dokud nemusíte; pokud čtyři body leží na kružnici, stačí si to poznamenat stranou, nebo obarvit hrany tětivového čtyřúhelníku,
\i při rýsování kolmice použijte místo trojúhelníku s ryskou dvě kružnice (mně to vždy dávalo přesnější výsledky),
\i nerýsujte přímky a oblouky delší, než je nutné (víc rušivých prvků znamená, že vidíte méně); dvakrát si rozmyslete, než do obrázku přidáte přímku (nebo než ji zvýrazníte či obarvíte),
\i pomocné čáry průběžně mažte,
\i při rýsování přesného obrázku můžete použít tvrzení, které se snažíte dokázat, nebo dílčí tvrzení, která jste už objevili,
\i překreslujte často (např. po přeformulování) a klidně rýsujte tentýž obrázek vícekrát (např. k ověřování hypotéz),
\i pokud přesný obrázek nakreslit nejde (např. úloha má záludné implicitně definované body), narýsujte přesně tak velkou jeho část, jak jen dokážete,
\i při rýsování trojúhelníku~$ABC$ umístěte horní vrchol tak, aby situace vypadala vizuálně symetricky (v dobře postavených úlohách je to vrchol~$A$),
\i někdy je užitečné umístit neobvyklou přímku svisle nebo vodorovně (např. úhlopříčku čtyřúhelníku),
\i nejprve načrtněte hrubý náčrt od ruky a pomůcky začněte používat až po krátkém seznámení se situací (pokud úlohu nevyřešíte už z náčrtu); to vám dá cit pro to, jak velký by obrázek měl být a kde začít rýsovat.
\enditems

\secc Jak přemýšlet při řešení úloh

Jakmile dostaneme úlohu a nakreslíme obrázek, začneme analyzovat:

\begitems
\i Předpokládejme, že dokazované tvrzení platí. Můžeme z toho odvodit něco zajímavého?
\i Které předpoklady o definovaných bodech jsme použili a které nedokážeme uchopit? Jsou to úhlové, nebo délkové podmínky?
\i Co je na úloze nejhůř uchopitelné? Které body se zdají nejvíc libovolné?
\i Představte si řešení. Z jakého dílčího tvrzení by tvrzení mohlo plynout?
\i Které úhly určují konfiguraci? Které úhly dokážeme pomocí nich vyjádřit?
\i Můžeme přeformulováním situace nějaký bod eliminovat?
\enditems

V praxi ne vždy přinese úspěch samotná analýza, jinak by to bylo příliš snadné~\SlightSmile{} Stále si můžeme se situací pohrát:
\begitems
\i Naslepo počítejte nějaké úhly a uvidíte, co se stane. Všímejte si, jestli jsme náhodou nenašli tětivový čtyřúhelník nebo jinou zajímavou konfiguraci.
\i Naslepo si zapište nějaké poměry (např. z podobných trojúhelníků) a zkuste je zkombinovat se známými rovnostmi délek. Všímejte si, jestli lze nově získané vztahy interpretovat jako podobnost dvou jiných trojúhelníků.
\i Definujte nějaké pěkně vypadající body, které by mohly situaci vyjasnit nebo o ní odhalit něco zajímavého (to je právě to, co budeme v tomto materiálu procvičovat~\SlightSmile{}).
\i Vyzkoušejte, jak situace vypadá ve zvláštních případech, nebo zkoumejte volné body v mezních polohách.
\enditems

Můžeme také tvořit hypotézy. Dobrá hypotéza je v obrovském množství případů klíčem k vyřešení těžkých úloh. Autoři úloh se ze všech sil snaží klíčová tvrzení skrýt a člověk je často prostě musí uhodnout. Dobře narýsovaný obrázek hádání usnadní. Metody hádání:

\begitems
\i Hádejte pěkné geometrické vlastnosti, nejčastěji tětivové čtyřúhelníky.
\i Jakmile máme nějaký odhad, ověřte, zda v narýsovaném obrázku platí alespoň přibližně.
\i Poté ověřte, zda ho dokážeme dokázat, klidně i s použitím dokazovaného tvrzení; být si jistý je důležité.
\i Předpokládejme, že uhodnuté tvrzení platí. Co bychom v naší situaci mohli odvodit? Možná bychom došli k něčemu, co obecně platit nemůže? Možná lze úlohu s jeho pomocí vyřešit v pár krocích?
\i Používejte symetrickou úvahu: pokud tohle platí, pak ze symetrie musí platit i něco dalšího\dots
\i Pokud hypotézu neumíme dokázat ani vyvrátit, zapište ji do čistopisu slovy. Pokud platí a je podstatná, může přinést body.
\enditems

\sec Známá dílčí tvrzení

V dnešních úlohách budeme potřebovat několik tvrzení, která jsou obecně velmi užitečná a stojí za to je umět použít přímo ve složitém obrázku. Zde je krátký seznam některých z nejdůležitějších, jež by nám mohly pomoci i při řešení úloh z tohoto materiálu:

\EnvId{4G1lSZGp6WOW_d5PrWMRp-zakladni-uhly}
\Theorem{Základní úhly}{Nechť $ABC$ je ostroúhlý trojúhelník. Pak:

\begitems
\i je-li $H$ ortocentrum, pak $|\angle BHC| = 180^\circ - \alpha$,
\i je-li $O$ střed opsané kružnice, pak $|\angle BOC| = 2\alpha$,
\i je-li $I$ střed vepsané kružnice, pak $|\angle BIC| = 90^\circ + \frac\alpha2$,
\i je-li $I_a$ střed kružnice připsané proti vrcholu~$A$, pak $|\angle BI_aC|=90^\circ-\frac\alpha2$.
\enditems
}{}

\EnvId{bvirmYCiJc6UXSLVpp95M-svrckuv-bod}
\Theorem{Švrčkův bod}{Nechť $ABC$ je trojúhelník se středem vepsané kružnice~$I$ a $|AB|\neq|AC|$. Nechť~$I_a$ je střed kružnice připsané proti vrcholu~$A$. Pak se osa úhlu~$\angle BAC$ a osa úsečky~$BC$ protínají na kružnici opsané trojúhelníku~$ABC$ v bodě, který je středem úsečky~$II_a$ a zároveň středem kružnice opsané čtyřúhelníku $BICI_a$.}{}

\EnvId{e3oWC42hI0-yUVRoB0_8O-anti-svrckuv-bod}
\Theorem{Anti-Švrčkův bod}{Nechť $ABC$ je trojúhelník s $|AB|\neq|AC|$. Nechť $I_b$ a~$I_c$ jsou po řadě střed kružnice připsané proti vrcholu~$B$ a proti vrcholu~$C$. Pak se vnější osa úhlu~$\angle BAC$ a osa úsečky~$BC$ protínají na kružnici opsané trojúhelníku~$ABC$ v bodě~$M$, který je středem úsečky~$I_bI_c$ a zároveň středem kružnice opsané čtyřúhelníku $BCI_bI_c$.}{}

\EnvId{hXmQS3PNqrvP6cXDIywLj-spiralni-podobnost-chodi-po-dvou}
\Theorem{Spirální podobnost chodí po dvou}{Jsou-li trojúhelníky~$OBC$ a~$OB_0C_0$ podobné a souhlasně orientované, pak jsou podobné i trojúhelníky $OBB_0$ a $OCC_0$.}{}

\EnvId{xHpsKi1FxuaOE4TsQDYXt-zrcadleni-ortocentra}
\Theorem{Zrcadlení ortocentra}{Nechť $ABC$ je trojúhelník s ortocentrem~$H$. Pak:
\begitems
\i zrcadlový obraz~$H$ podle přímky~$BC$ leží na kružnici opsané trojúhelníku~$ABC$,
\i obraz~$H$ ve středové souměrnosti se středem ve středu úsečky~$BC$ leží na kružnici opsané trojúhelníku $ABC$ v bodě diametrálně protilehlém k~$A$.
\enditems}{}

\sec Přidávání bodů do obrázku

Konečně a jsme u věci~\SlightSmile{} Při řešení úlohy se často stává, že nedokážeme pořádně uchopit nějakou hypotézu nebo dokazované tvrzení. Důvodem může být, že v obrázku chybí všechny body potřebné k vyřešení úlohy. Autoři úloh mají sklon skrývat klíčové body a tím dělat úlohy zajímavými a často mnohem těžšími. Umět přidat ty správné body je pravé umění. Pár obecných rad:

\begitems
\i Body přidáváme, abychom lépe pochopili nějakou hypotézu nebo dokazované tvrzení.
\i Přidáváme body, které vytvářejí pěkné vlastnosti, jako je podobnost, tětivové čtyřúhelníky, rovnoběžnost, \dots
\i Přidáváme body proto, abychom mohli odstranit jiné, hůř uchopitelné body.
\enditems

Existují opakující se situace, kde se vždycky vyplatí zvážit přidání nového bodu, když nic jiného nezabírá. V tomto materiálu procvičujeme zejména tyto:

\begitems \style n

\i Pokud jsou v úloze středy úseček, vyplatí se přidat další středy a vytvořit tak střední příčky.
\i Pokud je v úloze střed úsečky, středová souměrnost podle tohoto středu vytvoří rovnoběžník.
\i Pokud je v úloze pravoúhlý trojúhelník, vyplatí se doplnit ho na rovnoramenný trojúhelník.
\i Pokud je v úloze ortocentrum, vyplatí se přemýšlet o jeho obrazech v zrcadlení podle strany nebo podle středu strany daného trojúhelníku.
\i Pokud jsou v úloze rovnoběžné úsečky, vyplatí se zvážit stejnolehlosti, které zobrazí jednu úsečku na druhou, a zrcadlit skrz ně další body.
\i Pokud dvě kružnice procházejí společným bodem, zajímavým bodem může být jejich druhý průsečík. Vytvoří přinejmenším chordálu.
\i Někdy stačí protnout vhodné přímky, nebo vhodnou přímku a kružnici. Důvodů k tomu může být několik, například se může objevit rovnoramenný trojúhelník nebo tětivový čtyřúhelník.

\enditems

Tento seznam má k úplnosti daleko, ale na dalších 24 úloh by měl stačit. Obecně platí, že při přidávání bodů do obrázku jsou možnosti téměř neomezené.

\sec Úlohy

Následující úlohy jsou zhruba čtyř typů:

\begitems \style n
\i Klíčovým krokem je přidání správného bodu; zbytek je snadný.
\i Přidání je poměrně přímočaré, ale dotažení úlohy vyžaduje trochu práce.
\i Před přidáním bodu je třeba situaci prozkoumat, což teprve přidání vnukne.
\i Totéž jako výše, ale ani po takovém přidání to není snadné.
\enditems

\secc Standardní úlohy

Následující úlohy jsou zhruba seřazeny podle obtížnosti a používají zhruba řečeno \uv{standardní} triky.

\EnvId{yt0KVCE0mByjG0W5AhbYA-petiuhelnik-se-stejnymi-uhly}
\Problem{0}{}{
    Pětiúhelník~$ABCDE$ má všechny vnitřní úhly stejné. Dokažte, že osa úsečky~$AB$, osa úsečky~$CD$ a osa úhlu~$DEA$ se protínají v jednom bodě.
}{
    Úloha o pěti bodech. Co víc si přát? Paradoxně to nejde bez přidávání bodů. Stejné úhly by měly něco napovědět.
}{
    Stejné úhly lze využít protnutím vhodných dvojic prodloužených stran, čímž vzniknou rovnoramenné trojúhelníky. Osy stran mají najednou úplně jiný význam.
}{
    Klíčem je protnout prodloužené strany. Stejné vnitřní úhly dají rovnoramenné trojúhelníky, v~nichž se osy úseček $AB$ a~$CD$ změní na osy úhlů.

    Součet vnitřních úhlů pětiúhelníku je $540^\circ$, takže každý z~nich má velikost $108^\circ$. Všechny jsou menší než $180^\circ$, pětiúhelník $ABCDE$ je tedy konvexní, a jelikož žádný z~nich není přímý, pro každou stranu leží zbylé tři vrcholy striktně ve stejné otevřené polorovině určené přímkou této strany.

    Prodlužme stranu $EA$ za bod $A$ a~stranu $CB$ za bod $B$. Body $E$ a~$C$ leží ve stejné polorovině určené přímkou $AB$, takže obě tyto polopřímky míří do opačné poloroviny a s~úsečkou $AB$ svírají úhly $180^\circ - 108^\circ = 72^\circ$. Jejich součet je menší než $180^\circ$, a tak se polopřímky protnou v~bodě $P$. V trojúhelníku $ABP$ je $|\angle PAB| = |\angle PBA| = 72^\circ$, tedy $|PA| = |PB|$ a~osa úsečky $AB$ je zároveň osou úhlu $\angle APB$.

    Stejně prodlužme stranu $BC$ za bod $C$ a~stranu $ED$ za bod $D$. Body $B$ a~$E$ leží ve stejné polorovině určené přímkou $CD$, obě polopřímky tedy míří do opačné a s~úsečkou $CD$ svírají úhly $72^\circ$. Protnou se proto v~bodě $Q$, pro který je $|\angle QCD| = |\angle QDC| = 72^\circ$, takže $|QC| = |QD|$ a~osa úsečky $CD$ je osou úhlu $\angle CQD$.

    Bod $P$ leží na polopřímce opačné k~polopřímce $BC$ a~bod $Q$ na polopřímce opačné k~polopřímce $CB$, takže na přímce $BC$ jdou body v~pořadí $P$, $B$, $C$, $Q$. Bod $E$ na přímce $BC$ neleží, a tedy $EPQ$ je trojúhelník. Dále $A$ leží mezi $P$ a~$E$ a bod $D$ mezi $Q$ a~$E$.

    Z~tohoto pořadí plyne, že polopřímka $PA$ je polopřímkou $PE$ a~polopřímka $PB$ polopřímkou $PQ$, takže osa úsečky $AB$ je vnitřní osou úhlu trojúhelníku $EPQ$ při vrcholu $P$. Podobně jsou polopřímky $QD$, $QC$ polopřímkami $QE$, $QP$, a tak je osa úsečky $CD$ vnitřní osou úhlu při vrcholu $Q$. Nakonec $EA$ je polopřímkou $EP$ a~$ED$ polopřímkou $EQ$, takže osa úhlu $\angle DEA$ je vnitřní osou úhlu při vrcholu $E$. Všechny tři přímky jsou tedy vnitřními osami úhlů trojúhelníku $EPQ$ a~protínají se ve středu jeho vepsané kružnice.
}

\EnvId{rIhQMrZeKZ2MH-PmkjT9D-bod-na-teznici-rovnoramenny-trojuhelnik}
\Problem{0}{}{
    Uvnitř trojúhelníku~$ABC$, na jeho těžnici~$AM$, leží bod~$K$ takový, že $|CK|=|AB|$. Nechť~$L$ je průsečík přímek~$CK$ a~$AB$. Dokažte, že trojúhelník~$AKL$ je rovnoramenný.
}{
    Bod~$K$ na těžnici vypadá podivně. Každopádně musíme využít, že jde o těžnici, tedy že~$M$ je střed, ideálně tak, aby nám to pomohlo lépe pochopit rovnost $|AB|=|CK|$.
}{
    Klíčem je zrcadlit~$A$ podle~$M$ do~$A'$, což dá bodu~$M$ silný význam, neboť rovnost~$|AB|=|CK|$ se najednou stane~$|A'C|=|CK|$.
}{
    Rovnost $|CK| = |AB|$ spojuje dvě úsečky, které nemají společný ani bod, ani směr, a jediné, co o těžnici víme, je to, že $M$ je střed. Střed úsečky přitom volá po středové souměrnosti podle něj: nechť $A'$ je obraz~$A$ v souměrnosti se středem~$M$. Úhlopříčky $AA'$ a~$BC$ čtyřúhelníku $ABA'C$ se navzájem půlí, takže jde o rovnoběžník, a proto
    $$
    |CA'| = |AB| = |CK|, \qquad CA' \parallel AB.
    $$
    Podmínka ze zadání se tím přesunula do jednoho trojúhelníku.

    Bod~$K$ leží uvnitř trojúhelníku $ABC$, tedy uvnitř úsečky $AM$, a jelikož $M$ je střed $AA'$, leží i uvnitř úsečky $AA'$. Polopřímka $CK$ vychází z vrcholu $C$ a prochází vnitřním bodem trojúhelníku, takže protíná protější stranu~$AB$ v jejím vnitřním bodě; tím je právě $L$ a bod~$K$ leží mezi $C$ a~$L$.

    Bod $C$ neleží na přímce $AM$, takže trojúhelník $CKA'$ není degenerovaný; z $|CA'| = |CK|$ je rovnoramenný se základnou $KA'$. Jelikož $K$ leží uvnitř $AA'$, je polopřímka $A'K$ totožná s~$A'A$, a tedy
    $$
    |\angle CKA'| = |\angle CA'K| = |\angle CA'A|.
    $$

    Přímka $AA'$ protíná úsečku $BC$ v jejím vnitřním bodě $M$, takže odděluje $B$ od $C$. Úhly $\angle CA'A$ a~$\angle BAA'$ jsou proto střídavými úhly u rovnoběžek $CA'$, $AB$ a příčky $AA'$. Polopřímka $AB$ obsahuje $L$ a polopřímka $AA'$ obsahuje $K$, čili
    $$
    |\angle CA'A| = |\angle BAA'| = |\angle LAK|.
    $$

    Nakonec jsou úhly $\angle AKL$ a~$\angle A'KC$ vrcholové: polopřímky $KA$, $KA'$ jsou opačné (bod $K$ leží uvnitř $AA'$) a opačné jsou i polopřímky $KL$, $KC$ (bod $K$ leží mezi $C$ a~$L$). Dáme-li vše dohromady, dostáváme
    $$
    |\angle AKL| = |\angle A'KC| = |\angle LAK|,
    $$
    takže trojúhelník $AKL$ je rovnoramenný se základnou $AK$ a platí $|LA| = |LK|$.
}

\EnvId{ull-eZWewTAOKLQHeu11Z-kolmice-k-mh-pres-ortocentrum}
\Problem{0}{}{
    Nechť $ABC$ je ostroúhlý trojúhelník s ortocentrem~$H$ a nechť~$M$ je střed~$BC$. Body~$P$ a~$Q$ leží po řadě na stranách~$AB$ a~$AC$ tak, že $P$, $H$, $Q$ leží na jedné přímce kolmé k~$MH$. Dokažte, že $|HP| = |HQ|$.
}{
    V úloze máme~$H$, střed~$BC$ a nějak hraje roli i přímka~$MH$. No, co si počneme s tím ortocentrem~\SlightSmile?
}{
    Je-li~$H'$ obraz~$H$ ve středové souměrnosti se středem~$M$, pak víme, že padne na kružnici opsanou trojúhelníku~$ABC$ do bodu protilehlého k~$A$. Objevily se nějaké tětivové čtyřúhelníky?
}{
    Označme $\alpha = |\angle BAC|$. V úloze vystupuje ortocentrum i střed strany, zrcadlíme tedy ortocentrum podle tohoto středu: nechť $H'$ je obraz bodu~$H$ ve středové souměrnosti se středem~$M$. Podle tvrzení \uv{Zrcadlení ortocentra} leží $H'$ na kružnici opsané trojúhelníku $ABC$, a to v bodě diametrálně protilehlém k~$A$.

    Trojúhelník je ostroúhlý, takže $H$ leží uvnitř něj; neleží tedy na přímce $BC$, čili $H \neq M$ a $H' \neq H$. Přímka $MH$ ze zadání je tak přímkou $HH'$.

    Úsečky $BC$ a $HH'$ mají společný střed~$M$, takže $BHCH'$ je rovnoběžník. Podle tvrzení \uv{Základní úhly} je $|\angle BHC| = 180^\circ - \alpha$ a součet sousedních úhlů rovnoběžníku je $180^\circ$, takže
    $$
    |\angle HBH'| = |\angle HCH'| = \alpha. \eqno(*)
    $$

    Bod~$H$ neleží na přímce $AB$, a tedy $P \neq H$. Ani $H'$ na přímce $AB$ neleží: ta protíná opsanou kružnici jen v bodech $A$, $B$, přičemž $H' \neq A$, a rovnost $H' = B$ by znamenala, že $AB$ je průměr, čili $|\angle ACB| = 90^\circ$. Bod~$P$ je tedy různý od~$H$ i od~$H'$ a můžeme uvažovat kružnici $\omega_1$ s průměrem $PH'$.

    Přímka $PQ$ prochází bodem~$H$ a je kolmá k $MH = HH'$, takže $|\angle PHH'| = 90^\circ$ a bod~$H$ leží na~$\omega_1$. Leží na ní i bod~$B$: pro $P = B$ je to zřejmé, a pro $P \neq B$ je přímka $PB$ přímkou $AB$ a z Thaletovy věty nad průměrem $AH'$ je $|\angle PBH'| = |\angle ABH'| = 90^\circ$. Vznikl tak tětivový čtyřúhelník $PBHH'$.

    Úhly $\angle HPH'$ a $\angle HBH'$ přísluší v~$\omega_1$ téže tětivě $HH'$, takže jsou shodné, nebo je jejich součet $180^\circ$. První z nich je ostrý, neboť trojúhelník $PHH'$ má pravý úhel při~$H$, a druhý se podle $(*)$ rovná $\alpha < 90^\circ$; jsou tedy shodné, čili $|\angle HPH'| = \alpha$. (Pro $P = B$ jde rovnou o týž úhel.)

    Stejná úvaha platí pro~$Q$. Bod $H'$ neleží na přímce $AC$, neboť $H' = C$ by znamenalo, že $AC$ je průměr, čili $|\angle ABC| = 90^\circ$. Kružnice $\omega_2$ s průměrem $QH'$ prochází bodem~$H$ díky $|\angle QHH'| = 90^\circ$ a bodem~$C$ díky $|\angle ACH'| = 90^\circ$, takže z tětivy $HH'$ a $(*)$ dostáváme $|\angle HQH'| = |\angle HCH'| = \alpha$.

    Trojúhelníky $PHH'$ a $QHH'$ mají pravý úhel při~$H$ a úhel $\alpha$ při vrcholu $P$, resp. $Q$, takže ze součtu úhlů je $|\angle PH'H| = |\angle QH'H| = 90^\circ - \alpha$. Shodují se tedy ve straně $HH'$ a v obou úhlech při jejích krajních bodech, čili jsou podle věty $usu$ shodné. Odtud $|HP| = |HQ|$.
}

\EnvId{n2g_94VghWgVCIT8C-mse-spojnice-stredu-stran-a-uhlopricky}
\Problem{0}{}{
    Ve čtyřúhelníku~$ABCD$ svírá přímka spojující středy stran~$BC$ a~$AD$ stejné úhly s oběma úhlopříčkami. Dokažte, že úhlopříčky jsou stejně dlouhé.
}{
    Úhel mezi střední příčkou a úhlopříčkou čtyřúhelníku není něco, co bych vídal každý den. Klíčem bude zaměřit se na samotné středy. Co takhle přidat jeden nebo více dalších středů?
}{
    Dobrý střed, který přidat, je střed~$AB$ nebo~$CD$. Díky rovnoběžnosti plynoucí ze středních příček dokážeme podmínku ze zadání pěkně přeformulovat. Pak už to dorazíme.
}{
    Zadání dává dva středy stran, a to je přesně situace, kdy se vyplatí přidat další střed a nechat vzniknout střední příčky. Označme $M$, $N$ středy stran~$BC$, $AD$ a přidejme střed~$P$ strany~$AB$. Body $P$, $M$ jsou středy stran~$AB$, $BC$ trojúhelníku~$ABC$ a body $P$, $N$ jsou středy stran~$AB$, $AD$ trojúhelníku~$ABD$, takže se obě úhlopříčky objeví naráz jako střední příčky vycházející z~bodu $P$:
    $$
    \gather
    PM \parallel AC, \qquad |PM| = \tfrac12|AC|, \\
    PN \parallel BD, \qquad |PN| = \tfrac12|BD|.
    \endgather
    $$
    Dokazovaná rovnost $|AC| = |BD|$ se tím mění na $|PM| = |PN|$, čili na to, že trojúhelník~$PMN$ je rovnoramenný, a podmínka ze zadání mluví přesně o úhlech, které přímka~$MN$ svírá s jeho rameny.

    Ověřme, že $PMN$ je opravdu trojúhelník. Přímky $AC$, $BD$ nejsou rovnoběžné: konvexní čtyřúhelník má úhlopříčky, které se protínají uvnitř, a v nekonvexním leží jeden vrchol uvnitř trojúhelníku určeného zbylými třemi, přičemž přímka spojující tento vrchol s protilehlým vrcholem (tedy jedna z úhlopříček) protíná protilehlou stranu tohoto trojúhelníku, čili druhou úhlopříčku. Kdyby bod~$P$ ležel na přímce~$MN$, splynuly by přímky $PM$ i $PN$ s přímkou~$MN$ a vyšlo by $AC \parallel MN \parallel BD$, spor.

    Nechť $\varphi \in [0^\circ, 90^\circ]$ je společná velikost úhlu, který přímka~$MN$ svírá s oběma úhlopříčkami; díky rovnoběžnostem výše jej svírá i s přímkami $PM$ a~$PN$. Vnitřní úhel $|\angle PMN|$ trojúhelníku~$PMN$ je proto $\varphi$ nebo $180^\circ - \varphi$ a totéž platí pro $|\angle PNM|$. Součet dvou vnitřních úhlů trojúhelníku je menší než $180^\circ$, což vylučuje dvojici $\varphi$, $180^\circ - \varphi$ se součtem rovným $180^\circ$, i dvojici $180^\circ - \varphi$, $180^\circ - \varphi$ se součtem $360^\circ - 2\varphi \geq 180^\circ$. Zbývá jediná možnost
    $$
    |\angle PMN| = |\angle PNM| = \varphi,
    $$
    takže trojúhelník~$PMN$ je rovnoramenný, $|PM| = |PN|$ a odtud $|AC| = 2|PM| = 2|PN| = |BD|$.
}

\EnvId{Z8MdAzTmbIxxpEyc-UnMy-ortocentra-tetivoveho-ctyruhelniku}
\Problem{0}{Česko-slovenská olympiáda 2009}{
    Nechť $ABCD$ je tětivový čtyřúhelník. Dokažte, že přímka spojující ortocentrum trojúhelníku~$ABC$ s ortocentrem trojúhelníku~$ABD$ je rovnoběžná s přímkou~$CD$.
}{
    Ortocentra jsou složité body, pokud narýsujeme všechny výšky. Takové kreslení hned zavrhněte. Jejich obrazy ve vhodných zrcadleních dají situaci lépe pochopit. Pomůže také přeformulovat, co vlastně dokazujeme, aby po zrcadlení ortocenter bylo tvrzení stále jasně vyslovitelné.
}{
    Nejprve stačí ukázat, že úsečka spojující ortocentra má stejnou délku jako~$CD$ (pak máme rovnoběžník). Toto zjištění nám dodá jistotu, že když tato ortocentra zrcadlíme podle osy, máme stále dobré tvrzení k dokázání, totiž rovnost dvou tětiv kružnice.
}{
    Označme $\omega$ kružnici opsanou čtyřúhelníku $ABCD$, $O$ její střed, $M$ střed strany $AB$ a~$H_1$, $H_2$ po řadě ortocentra trojúhelníků $ABC$ a~$ABD$. Body $A$, $B$, $C$, $D$ leží na $\omega$, takže žádné tři z~nich nejsou kolineární a~oba trojúhelníky jsou nedegenerované. Podstatné je, že oběma je opsaná táž kružnice $\omega$ a~že mají společnou stranu $AB$.

    Výšky nekreslíme, z~definice ortocentra nám stačí, že $CH_1 \perp AB$ a~$DH_2 \perp AB$, tedy $CH_1 \parallel DH_2$. Tvrzení by tak bylo hotové, kdyby byl čtyřúhelník $CH_1H_2D$ rovnoběžník. Rovnoběžnost jedné dvojice protilehlých stran už máme, takže jde jen o~délky: potřebujeme $|H_1H_2| = |CD|$ se souhlasnou orientací. V~tom je celý přínos přeformulování, protože délky se zrcadlením nemění. Ortocentra teď smíme vyměnit za jejich obrazy a~dokazované tvrzení půjde i~potom jasně vyslovit.

    Zrcadlíme podle středu společné strany, tedy podle bodu $M$; ten je pro oba trojúhelníky týž, a~to je jediný důvod, proč se výsledné obrazy dají porovnat. Podle věty o~zrcadlení ortocentra je obrazem $H_1$ ve středové souměrnosti podle $M$ bod $C_1$ kružnice $\omega$ diametrálně protilehlý k~$C$; táž věta pro trojúhelník $ABD$ dává, že obrazem $H_2$ je bod $D_1 \in \omega$ diametrálně protilehlý k~$D$. Z~$C \neq D$ plyne $C_1 \neq D_1$, a~tedy $H_1 \neq H_2$, takže přímka spojující obě ortocentra vůbec existuje.

    Středová souměrnost podle $M$ zobrazuje přímku $H_1H_2$ na přímku $C_1D_1$ a~středová souměrnost podle $O$ zobrazuje přímku $CD$ na tutéž přímku $C_1D_1$, jelikož $C_1$, $D_1$ jsou obrazy $C$, $D$. Zbývající tvrzení je tak rovnost dvou tětiv kružnice $\omega$ a~ta je po tomto pozorování zadarmo:
    $$
    |H_1H_2| = |C_1D_1| = |CD|.
    $$
    Středová souměrnost navíc zobrazuje každou přímku na přímku s~ní rovnoběžnou, takže $H_1H_2 \parallel C_1D_1$ i~$CD \parallel C_1D_1$, a~proto $H_1H_2 \parallel CD$. (Při některých polohách vrcholů obě přímky splynou, například když je $AB$ průměr $\omega$, protože pak $H_1 = C$ a~$H_2 = D$; totožné přímky považujeme za rovnoběžné.)
}

\EnvId{alkX6MavufXBEbJ3d6M4z-ortocentrum-v-rovnoramennem-lichobezniku}
\Problem{0}{}{
    Nechť $ABCD$ je rovnoramenný lichoběžník s $AB \parallel CD$ vepsaný do kružnice se středem~$O$. Nechť~$H$ je ortocentrum trojúhelníku $ABD$ a nechť~$P$ je průsečík~$DH$ a~$AB$. Dokažte, že $CH \parallel OP$.
}{
    Máme ortocentrum. Nemusí být hned zřejmé, co s ním, ale co takhle zrcadlit ho podle něčeho?
}{
    Nechť~$H'$ je obraz~$H$ v zrcadlení podle~$AB$. Padne na kružnici opsanou čtyřúhelníku~$ABCD$. Teď obrázek obsahuje vše, co potřebujeme, a zbývá jen spojit pár faktů.
}{
    V úloze vystupuje ortocentrum, podívejme se tedy na jeho obrazy v zrcadlení podle strany a podle středu strany trojúhelníku $ABD$. Označme $\omega$ kružnici opsanou čtyřúhelníku $ABCD$ a~$H'$ obraz bodu $H$ v osové souměrnosti podle přímky $AB$. Ukážeme, že $H'$ je bod diametrálně protilehlý k~$C$; zbytek už bude střední příčka.

    Přímka $DH$ je výška z~$D$ v trojúhelníku $ABD$, takže je kolmá na $AB$ a $P$ je její pata. Osová souměrnost podle $AB$ zobrazuje přímku $DH$ samu na sebe, neboť je na $AB$ kolmá; bod $H'$ proto leží na přímce $DH$ a $P$ je střed úsečky $HH'$.

    Nechť $M$ je střed úsečky $AB$ a $s$ její osa. Bod $O$ leží na $s$, neboť je stejně vzdálený od $A$ a~$B$. Přímka $s$ je kolmá na $AB$, tedy i na $CD$, a kolmice k tětivě vedená středem kružnice prochází středem této tětivy, takže $s$ je zároveň osou úsečky $CD$. Souměrnost podle $s$ proto zobrazuje $\omega$ na sebe, zachovává její střed $O$, zaměňuje $A$ s~$B$ a $C$ s~$D$.

    Podle druhé části věty \uv{Zrcadlení ortocentra}, použité na trojúhelník $ABD$ a jeho stranu $AB$, je obrazem bodu $H$ ve středové souměrnosti se středem~$M$ bod $E$ diametrálně protilehlý k~$D$. Středová souměrnost se středem~$M$ je přitom složením osových souměrností podle dvou navzájem kolmých přímek procházejících bodem $M$, tedy podle $AB$ a podle $s$: zrcadlíme-li bod $H$ nejprve podle $AB$ a potom podle $s$, dostaneme $E$. První krok dá $H'$, takže $H'$ je obrazem $E$ v souměrnosti podle $s$. Ta zobrazuje $D$ na $C$ a $O$ na $O$, čili bod diametrálně protilehlý k~$D$ zobrazuje na bod diametrálně protilehlý k~$C$. Bod $H'$ je tedy diametrálně protilehlý k~$C$; speciálně leží na $\omega$ a $O$ je střed úsečky $CH'$.

    Bod $C$ neleží na přímce $DH$, neboť ta je kolmá na $CD$, a tedy protíná přímku $CD$ jedině v~$D$. Různé jsou i body $H$ a~$H'$: kdyby splynuly, ležel by $H$ na $AB$, čili $H = P$, a zároveň na $\omega$, čili $P$ by byl jeden z bodů $A$, $B$. Souměrnost podle $s$ však zobrazuje patu kolmice z~$D$ na $AB$ na patu kolmice z~$C$ na $AB$, takže by těmito patami byly právě body $A$ a~$B$. Z rovnoběžnosti $CD \parallel AB$ by pak vyšlo $|CD| = |AB|$ a $ABCD$ by byl rovnoběžník, ne lichoběžník.

    Body $H$, $H'$, $C$ tedy tvoří trojúhelník, přičemž $P$ je střed jeho strany $HH'$ a $O$ střed strany $H'C$. Úsečka $OP$ je jeho střední příčkou, a proto $OP \parallel CH$.
}

\EnvId{WL5FFRrQH7pTblc3sQzBH-pravy-uhel-u-stredu-usecky}
\Problem{0}{Polsko 2002}{
    Nechť $ABCD$ je konvexní čtyřúhelník s $|AB|<|BC|$ a $|AD|<|CD|$. Body $P$ a~$Q$ leží po řadě na~$BC$ a~$CD$ tak, že $|BP|=|BA|$ a $|DQ|=|DA|$. Nechť~$M$ je střed~$PQ$. Dokažte, že pokud $|\angle BMD|=90^\circ$, pak je $ABCD$ tětivový čtyřúhelník.
}{
    Potřebujeme využít zvláštní podmínku $|\angle BMD|=90^\circ$. Bod~$M$ je střed úsečky, což už něco napovídá, a pravý úhel nápovědu jen zesiluje. Každopádně se vyplatí použít barevné tužky k označení úseček stejné délky.
}{
    Klíčem je zrcadlit řekněme~$D$ podle~$M$. Po označení všech stejných úseček vznikajících z výsledného rovnoběžníku a rovnoramenného trojúhelníku by vám mělo něco padnout do oka.
}{
    Bod $M$ je středem úsečky, takže se nabízí zrcadlit podle něj a vyrobit rovnoběžník; předpoklad $|\angle BMD|=90^\circ$ napovídá, který bod si to zaslouží. Nechť tedy $D'$ je obraz bodu $D$ ve středové souměrnosti se středem $M$ (úhel $\angle BMD$ je definovaný, takže $M \neq D$, a tedy i $D' \neq D$).

    Ze zadání je $|BP|=|BA|<|BC|$ a $|DQ|=|DA|<|DC|$, takže $P$ je vnitřním bodem strany $BC$ a~$Q$ vnitřním bodem strany $CD$. Čtyřúhelník $ABCD$ je konvexní, takže žádné tři jeho vrcholy neleží na jedné přímce; přímky $BC$ a~$CD$ se proto protínají jedině v~$C$ a $P \neq Q$.

    Úsečky $PQ$ a~$DD'$ mají společný střed $M$, takže $DPD'Q$ je rovnoběžník, a tedy
    $$
    |PD'| = |QD| = |DA|, \qquad \overline{PD'} = \overline{DQ}.
    $$
    Přímka $BM$ je kolmá na $MD$, čili na $DD'$, a prochází středem $M$ této úsečky. Je to tedy osa úsečky $DD'$, trojúhelník $BDD'$ je rovnoramenný a
    $$
    |BD'| = |BD|.
    $$

    Trojúhelníky $BPD'$ a~$BAD$ se tak shodují ve všech třech stranách: $|BP|=|BA|$ ze zadání, $|PD'|=|AD|$ z rovnoběžníku a $|BD'|=|BD|$ z rovnoramennosti. Podle $sss$ jsou shodné, přičemž stranám $BP$, $PD'$ odpovídají strany $BA$, $AD$, takže úhly u vrcholů $P$ a~$A$ se rovnají:
    $$
    |\angle BPD'| = |\angle BAD|.
    $$

    Zbývá vyčíst úhel $\angle BPD'$ z~původního obrázku. Bod $P$ leží uvnitř úsečky $BC$, takže $\overline{PB}$ je kladným násobkem $\overline{CB}$; bod $Q$ leží uvnitř úsečky $DC$, takže $\overline{PD'} = \overline{DQ}$ je kladným násobkem $\overline{DC}$, čili záporným násobkem $\overline{CD}$. Ramena úhlů $\angle BPD'$ a~$\angle BCD$ jsou tedy rovnoběžná, jedna dvojice souhlasně a druhá opačně orientovaná, odkud
    $$
    |\angle BAD| + |\angle BCD| = |\angle BPD'| + |\angle BCD| = 180^\circ.
    $$
    Konvexní čtyřúhelník, v němž je součet protilehlých úhlů $180^\circ$, je tětivový, takže $ABCD$ je tětivový čtyřúhelník.
}

\EnvId{Tax9-rYDqJvAbCiA0VqMC-paty-kolmic-a-stred-strany}
\Problem{0}{USA TST 2000}{
    Nechť $ABCD$ je tětivový čtyřúhelník, jehož úhlopříčky se protínají v~$P$. Nechť~$K$ a~$L$ jsou po řadě paty kolmic z~$P$ na~$AB$ a~$CD$ a nechť~$M$ je střed strany~$AD$. Dokažte, že $|MK|=|ML|$.
}{
    Vyplatí se narýsovat úsečku~$AD$ vodorovně (nezvykle), což dá nejlepší obrázek (překreslete hned teď, pokud jste tak neudělali). I v novém obrázku nemusí být jasné, co se má stát. Střed~$AD$ nevypadá, že by se s ním dalo snadno pracovat. Možná potřebuje ke spolupráci nějaké další středy. Ideálně takové, které také souvisejí s body~$K$ a~$L$.
}{
    Klíčem jsou zbývající středy stran trojúhelníku~$ADP$. Pak se stane kouzlo a jsme jen pár snadných argumentů od řešení.
}{
    Bod~$M$ je střed strany~$AD$ a sám o sobě nedává ani pěkný úhel, ani použitelnou délku. Když v úloze takto osaměle visí střed úsečky, přidáváme další středy: doplňme~$M$ na trojici středů stran trojúhelníku~$ADP$, tedy přidejme střed~$X$ úsečky~$AP$ a střed~$Y$ úsečky~$DP$. Odměna přijde okamžitě. Bod~$K$ leží na přímce~$AB$ a~$PK \perp AB$, takže $|\angle AKP| = 90^\circ$, případně $K = A$; tak či onak $K$ leží na Thaletově kružnici nad průměrem~$AP$, jejímž středem je právě~$X$. Stejně tak~$L$ leží na Thaletově kružnici nad průměrem~$DP$ se středem~$Y$. Jediné přidání tak svázalo $K$, $L$ i~$M$.

    Žádné tři z~bodů $A$, $B$, $C$, $D$ nejsou kolineární, neboť leží na kružnici. Bod~$P$ leží na přímkách $AC$ i~$BD$, takže $P \neq A$, $P \neq D$ a~$P$ neleží na přímce~$AB$ ani na přímce~$CD$ (jinak by přímka~$AB$ splynula s~$AC$, resp. $CD$ s~$BD$). Speciálně $K \neq P$ a~$L \neq P$.

    Z~Thaletových kružnic je
    $$
    |XK| = |XP| = \tfrac12|AP|, \qquad |YL| = |YP| = \tfrac12|DP|.
    $$
    Stejnolehlost se středem~$A$ a koeficientem~$\tfrac12$ zobrazí $P$ na~$X$ a~$D$ na~$M$, takže $|MX| = \tfrac12|DP|$ a polopřímka~$XM$ míří stejným směrem jako~$PY$. Stejnolehlost se středem~$D$ a koeficientem~$\tfrac12$ zobrazí $P$ na~$Y$ a~$A$ na~$M$, takže $|MY| = \tfrac12|AP|$ a polopřímka~$YM$ míří stejným směrem jako~$PX$. Délky si tedy odpovídají křížem:
    $$
    |MX| = |YL| = \tfrac12|DP|, \qquad |XK| = |MY| = \tfrac12|AP|.
    $$
    Trojúhelníky $MXK$ a~$LYM$ se tak shodují ve dvou dvojicích stran, a shodnou-li se i v úhlech, které tyto strany svírají, budou podle věty $sus$ shodné a vyjde $|MK| = |LM|$. Úloha se tím zredukovala na rovnost $|\angle MXK| = |\angle MYL|$.

    Tady se už orientovaným úhlům nevyhneme: paty $K$, $L$ mohou podle konfigurace padnout i mimo úsečky $AB$, $CD$ a polopřímka~$XP$ nemusí ležet uvnitř úhlu~$MXK$, zatímco znaménka ohlídají všechny polohy najednou. Při pevně zvolené orientaci roviny označme pro polopřímky $p$, $q$ se společným počátkem $\angle(p \to q)$ úhel otočení, které převede $p$ na~$q$, braný modulo $360^\circ$; pro přímky $p$, $q$ nechť $\angle(p,q)$ je orientovaný úhel modulo $180^\circ$, takže jeho dvojnásobek je dobře určen modulo $360^\circ$.

    Položme $\omega = \angle(XP \to XK)$ a~$\omega' = \angle(YP \to YL)$; jsou to úhly otočení se středy $X$, $Y$, která zobrazí $P$ na~$K$, resp. na~$L$. Klíčové je, že $\omega' = -\omega$. Body $A$, $K$, $P$ leží na kružnici se středem~$X$, takže středový úhel je dvojnásobkem obvodového a platí
    $$
    \omega = 2\angle(AP,AK) = 2\angle(AC,AB),
    $$
    neboť $P$ leží na přímce~$AC$ a~$K$ na přímce~$AB$. (Je-li $K = A$, což nastane právě při $AB \perp AC$, je~$X$ střed úsečky~$KP$, čili $\omega = 180^\circ$, a rovnost platí také.) Stejně tak $\omega' = 2\angle(DB,DC)$. Ze čtyř bodů na jedné kružnici je $\angle(AB,AC) = \angle(DB,DC)$, a proto
    $$
    \omega = -2\angle(AB,AC) = -2\angle(DB,DC) = -\omega'.
    $$

    Zbývá to poskládat. Polopřímka~$XM$ míří stejným směrem jako~$PY$, tedy opačně než~$YP$, a polopřímka~$XP$ míří stejným směrem jako~$MY$, tedy opačně než~$YM$; otočení obou ramen o~$180^\circ$ úhel nemění, takže
    $$
    \angle(XM \to XP) = \angle(YP \to YM) = -\angle(YM \to YP).
    $$
    Dohromady s~$\omega' = -\omega$ dostáváme
    $$
    \eqalign{
        \angle(XM \to XK) &= \angle(XM \to XP) + \omega \cr
        &= -\angle(YM \to YP) - \omega' = -\angle(YM \to YL). \cr
    }
    $$
    Orientované úhly $\angle(XM \to XK)$ a~$\angle(YM \to YL)$ jsou tedy opačné, čili neorientované úhly $|\angle MXK|$ a~$|\angle MYL|$ jsou stejné. Podle věty $sus$ jsou tedy trojúhelníky $MXK$ a~$LYM$ shodné, a proto $|MK| = |ML|$.
}

\EnvId{0TlY8HJdJ_fVKGlkNgzpO-bod-uvnitr-rovnobezniku}
\Problem{0}{}{
    Uvnitř rovnoběžníku $ABCD$ leží bod $P$, pro který platí $|\angle APB| + |\angle CPD| = 180^\circ$. Dokažte, že $|\angle PBC|=|\angle PDC|$.
}{
    Zkuste posunout bod~$P$ tak, aby podmínka na součet úhlů dávala smysl.
}{
    Klíčovým bodem je obraz bodu~$P$ v posunutí ve směru~$BC$ (tedy~$AD$); označme ho~$P'$. Pak máme tětivový čtyřúhelník $CP'DP$ i spoustu rovnoběžníků.
}{
    Podmínka $|\angle APB| + |\angle CPD| = 180^\circ$ váže dva úhly u téhož vrcholu~$P$, a tak se s ní nedá nic dělat. Rovnoběžník má ale dvě dvojice rovnoběžných shodných stran, což volá po posunutí zobrazujícím jednu dvojici na druhou. Posuneme-li jím i bod~$P$, přesune se jeden z těch dvou úhlů k novému vrcholu a ze součtu $180^\circ$ se stane podmínka tětivovosti.

    Nechť je tedy $\tau$ posunutí zobrazující $B$ na~$C$, tedy posunutí o vektor $BC = AD$, a nechť $P' = \tau(P)$. Jelikož $\tau(A) = D$ a~$\tau(B) = C$, přejde trojúhelník $ABP$ na trojúhelník $DCP'$, takže
    $$
    |\angle DP'C| = |\angle APB| = 180^\circ - |\angle CPD|.
    $$

    Nejprve určeme, kde bod~$P'$ leží. Bod~$P$ je vnitřním bodem rovnoběžníku, takže rovnoběžka s~$BC$ vedená bodem~$P$ protíná stranu $AB$ ve vnitřním bodě~$X$ a stranu $DC$ ve vnitřním bodě~$Y$, přičemž $P$ leží uvnitř úsečky $XY$. Bod $\tau(X)$ leží na přímce $\tau(AB) = DC$ i na přímce $XY$ (posunutí je ve směru $XY$), a ty se protínají jedině v~$Y$, čili $\tau(X) = Y$. Proto $|XP'| = |XP| + |XY| > |XY|$, bod~$P'$ leží na polopřímce $XY$ za bodem~$Y$, a body $P$, $P'$ tak leží v opačných polorovinách určených přímkou $CD$.

    Body $C$, $D$, $P$ neleží na jedné přímce; označme $\omega$ kružnici opsanou trojúhelníku $CDP$. Množina bodů poloroviny neobsahující $P$, z nichž je úsečka $CD$ vidět pod úhlem $180^\circ - |\angle CPD|$, je právě oblouk kružnice $\omega$ v této polorovině. Bod~$P'$ do ní podle předchozích dvou odstavců patří, takže leží na~$\omega$ a $CP'DP$ je tětivový čtyřúhelník. Tětiva $CD$ odděluje na~$\omega$ body $P'$ a~$P$, čili v cyklickém pořadí jdou body $C$, $P'$, $D$, $P$; body $P'$ a~$D$ proto leží na témž oblouku určeném tětivou $CP$.

    Zbývá přenést úhel $\angle PBC$ ke kružnici $\omega$. Posunutí $\tau$ zobrazuje $B$ na~$C$ a~$P$ na~$P'$, takže úsečky $BP$ a~$CP'$ jsou rovnoběžné, stejně dlouhé a stejně orientované; jelikož $P$ neleží na přímce $BC$, je $BPP'C$ nedegenerovaný rovnoběžník. Jeho protilehlé úhly u vrcholů $B$ a~$P'$ se rovnají:
    $$
    |\angle PBC| = |\angle PP'C|.
    $$
    Úhly $\angle PP'C$ a~$\angle PDC$ jsou obvodové úhly kružnice $\omega$ nad tětivou $CP$ z téhož oblouku, a tedy shodné. Dohromady
    $$
    |\angle PBC| = |\angle PP'C| = |\angle PDC|.
    $$
}

\EnvId{RV54g2hn6AMdIasT4IODZ-konvexni-sestiuhelnik-se-shodnymi-stranami}
\Problem{0}{Česko-slovenská olympiáda 2026}{
    Je dán konvexní šestiúhelník $ABCDEF$ takový, že $|AB| = |BC|$, $|CD| = |DE| \neq |AD|$, $|EF| = |FA|$, $|\angle BDC| + |\angle EDF| = |\angle FDB|$. Dokažte, že $|\angle CBA| + |\angle EDC| + |\angle AFE| = 360^\circ$.
}{
    Podmínka na součet úhlů je klíčová. Potřebujeme zavést bod, který podmínku na součet úhlů využije a udělá ji o něco smysluplnější.
}{
    Klíčovým bodem je zrcadlový obraz $C$ podle $BD$, nebo analogicky zrcadlový obraz $E$ podle $DF$. Je to týž bod. A je to vše, co potřebujeme.  Nezapomeňte využít nerovnostní část podmínky $|CD|=|DE| \neq |AD|$. Bez ní úloha neplatí.
}{
    Podmínka $|\angle BDC| + |\angle EDF| = |\angle FDB|$ říká, že úhly $BDC$ a~$EDF$ přesně vyplní úhel $BDF$. To nás vede k~tomu překlopit $C$ podle přímky $BD$ a~$E$ podle přímky $DF$: oba obrazy padnou dovnitř úhlu $BDF$ a~prvním krokem bude ukázat, že splynou.

    Nejprve si ujasněme konfiguraci. Úhlopříčky z~vrcholu $D$ dělí vnitřní úhel konvexního šestiúhelníku při $D$ v~tom pořadí, v~jakém jdou vrcholy po obvodu, takže polopřímky $DC$, $DB$, $DA$, $DF$, $DE$ jdou v~tomto úhlovém pořadí a~platí
    $$
    |\angle CDB| + |\angle BDF| + |\angle FDE| = |\angle CDE| < 180^\circ.
    $$
    Přímka $BD$ proto odděluje $C$ od bodů $A$, $F$, $E$ a~přímka $DF$ odděluje $E$ od bodů $C$, $B$, $A$. Jelikož oba sčítance v~podmínce ze zadání jsou kladné, je také $|\angle BDC| < |\angle FDB|$ a~$|\angle FDE| < |\angle FDB|$.

    Označme $P$ obraz bodu $C$ v~osové souměrnosti podle přímky $BD$ a~$Q$ obraz bodu $E$ v~osové souměrnosti podle přímky $DF$. Ukážeme, že $P = Q$. Bod $P$ leží na opačné straně přímky $BD$ než $C$, tedy na téže jako $F$, a~$|\angle BDP| = |\angle BDC| < |\angle FDB|$, takže polopřímka $DP$ leží uvnitř úhlu $BDF$. Symetricky $|\angle FDQ| = |\angle FDE| < |\angle FDB|$ a~uvnitř úhlu $BDF$ leží i~polopřímka $DQ$. Ze zadání je $|\angle BDP| = |\angle FDB| - |\angle EDF|$, a~tedy
    $$
    |\angle FDP| = |\angle FDB| - |\angle BDP| = |\angle EDF| = |\angle FDQ|.
    $$
    Polopřímky $DP$ a~$DQ$ svírají s~přímkou $DF$ stejný úhel a~leží u~ní na téže straně (obě uvnitř úhlu $BDF$), proto splývají. Spolu s~$|DP| = |DC| = |DE| = |DQ|$ dostáváme $P = Q$.

    Jediný bod $P$ tak splňuje naráz všechny tři délkové podmínky ze zadání:
    $$
    \gather
    |BP| = |BC| = |BA|, \\
    |DP| = |DC| = |DE|, \\
    |FP| = |FE| = |FA|.
    \endgather
    $$

    První a~třetí z~nich říkají, že $P$ leží na kružnici $k_B$ se středem $B$ a~poloměrem $|BA|$ i~na kružnici $k_F$ se středem $F$ a~poloměrem $|FA|$. Ty mají různé středy, takže se protínají nejvýše ve dvou bodech. Osová souměrnost podle přímky $BF$ zobrazí každou z~nich samu na sebe, a~jelikož vnitřní úhel šestiúhelníku při vrcholu $A$ je menší než $180^\circ$, bod $A$ na přímce $BF$ neleží a~jeho obraz $A'$ je od něj různý. Společnými body kružnic $k_B$, $k_F$ jsou tedy právě $A$ a~$A'$ a~bod $P$ je jedním z~nich. Možnost $P = A$ však odpadá: dávala by $|AD| = |PD| = |CD|$, což je ve sporu s~předpokladem $|CD| \neq |AD|$. Proto $P = A'$. Toto je jediné místo, kde se nerovnostní část podmínky použije, a~je nenahraditelné: kdyby $P$ splynul s~$A$, ležel by mimo trojúhelník $BDF$ a~závěrečný součet úhlů by se rozpadl.

    V~konvexním šestiúhelníku odřezává úhlopříčka $BF$ trojúhelník $ABF$, takže přímka $BF$ odděluje $A$ od $D$. Bod $P = A'$ proto leží v~otevřené polorovině s~hraniční přímkou $BF$ obsahující $D$. Zároveň polopřímka $DP$ leží uvnitř úhlu $BDF$ a~$P \neq D$. Průnikem vnitřku úhlu $BDF$ s~touto polorovinou je právě vnitřek trojúhelníku $BDF$, takže $P$ je vnitřním bodem trojúhelníku $BDF$.

    Osová souměrnost podle $BD$ zobrazí trojúhelník $BCD$ na $BPD$, souměrnost podle $DF$ zobrazí $DEF$ na $DPF$ a~souměrnost podle $BF$ zobrazí $BAF$ na $BPF$, čili
    $$
    \gather
    |\angle BPD| = |\angle BCD|, \\
    |\angle DPF| = |\angle DEF|, \\
    |\angle FPB| = |\angle FAB|.
    \endgather
    $$
    Bod $P$ leží uvnitř trojúhelníku $BDF$, takže úhly u~něj dávají plný úhel, a~tedy
    $$
    |\angle BCD| + |\angle DEF| + |\angle FAB| = 360^\circ.
    $$
    Součet vnitřních úhlů šestiúhelníku je $720^\circ$, odkud
    $$
    |\angle CBA| + |\angle EDC| + |\angle AFE| = 720^\circ - 360^\circ = 360^\circ.
    $$
}

\EnvId{L-IkptXbj4-UJ_vpMWLKS-bod-s-peti-shodnymi-uhly}
\Problem{0}{Česko-slovenská olympiáda 2024}{
    Bod $P$ ležící uvnitř konvexního čtyřúhelníku $ABCD$ splňuje rovnosti
    $$|\angle PAD| = |\angle ADP| = |\angle CBP| = |\angle PCB| = |\angle CPD|.$$
    Nechť $O$ je střed kružnice opsané trojúhelníku $CPD$. Dokažte, že $|OA| = |OB|$.
}{
    Úhly implikují rovnoběžné přímky. Napovídají správný tah vpřed.
}{
    Bodem, který zavést, je průsečík $X$ přímek $BC$ a $AD$, neboť $PCXD$ je teď rovnoběžník. Úlohu lze nyní vyřešit správným pohledem na to, co máme.
}{
    Označme $\varphi$ společnou velikost pěti úhlů ze zadání. Trojúhelník $APD$ (bod $P$ leží uvnitř čtyřúhelníku, na přímce $AD$ tedy neleží) má při vrcholech $A$ a~$D$ úhly velikosti $\varphi$, je proto rovnoramenný s~$|PA| = |PD|$ a~navíc $\varphi < 90^\circ$; stejně je $|PB| = |PC|$.

    Zbylé dvě rovnosti mluví o~úhlech u~příčky. Jelikož $P$ leží uvnitř konvexního čtyřúhelníku $ABCD$, jdou polopřímky $PA$, $PB$, $PC$, $PD$ kolem bodu $P$ v~tomto cyklickém pořadí a~součet úhlů $\angle APB$, $\angle BPC$, $\angle CPD$, $\angle DPA$ je $360^\circ$. Měříme-li tedy úhly od polopřímky $PD$ ve směru $D \to A \to B \to C$, leží polopřímka $PA$ pod úhlem $|\angle DPA| < 180^\circ$, kdežto polopřímka $PC$ pod úhlem $360^\circ - |\angle CPD| > 180^\circ$, takže přímka $PD$ odděluje $A$ od $C$. Úhly $\angle ADP$ a~$\angle DPC$ jsou tak střídavé úhly přímek $AD$ a~$PC$ u~příčky $DP$, a~jelikož jsou shodné, je $AD \parallel PC$. Analogicky přímka $PC$ odděluje $B$ od $D$ a~ze shodnosti úhlů $\angle PCB$ a~$\angle CPD$ dostaneme $BC \parallel PD$.

    Dvě dvojice rovnoběžek si říkají o~průsečík. Přímky $AD$ a~$BC$ rovnoběžné nejsou (jinak by byly rovnoběžné i~$PC$ a~$PD$, tedy by $C$, $P$, $D$ ležely na jedné přímce), takže se protínají v~bodě $X$. Z~$PC \parallel XD$ a~$CX \parallel PD$ je $PCXD$ rovnoběžník. Dokážeme dokonce víc, než se žádá:
    $$
    |OA| = |OX| = |OB|.
    $$

    Klíčové je pozorování, že \textbf{osa úsečky $PC$ je zároveň osou úsečky $AX$}.

    Nechť $M$ je střed úsečky $PC$, $N$ střed úsečky $AD$ a~$Y$ střed úsečky $AX$. Posunutí, které zobrazí $P$ na $C$, převede díky rovnoběžníku $PCXD$ bod $D$ na $X$, takže úsečka $AX$ vznikne z~úsečky $AD$ posunutím jednoho konce; její střed se přitom posune o~polovinu tohoto posunutí, čili posunutí zobrazující $P$ na $M$ zobrazuje $N$ na $Y$. Úsečka $MY$ je tedy obrazem úsečky $PN$ a~má stejný směr. Trojúhelník $APD$ je rovnoramenný, takže jeho těžnice $PN$ je zároveň výškou, čili $PN \perp AD$, a~proto i~$MY \perp AD$.

    Přímka $MY$ tak prochází středem $M$ úsečky $PC$ kolmo na $PC$ (neboť $PC \parallel AD$): je to osa úsečky $PC$. Zároveň je kolmá na přímku $AD$ v~bodě $Y$, což je střed úsečky $AX$, jejíž oba krajní body leží na~$AD$: je to tedy i~osa úsečky $AX$. (Body $A$ a~$X$ jsou různé: $|\angle ADP| = \varphi$, kdežto $|\angle XDP| = 180^\circ - \varphi$ jako úhel rovnoběžníku sousední k~$\angle DPC$, a~$\varphi \neq 90^\circ$.)

    Bod $O$ je střed kružnice opsané trojúhelníku $CPD$, leží tedy na ose úsečky $PC$, a~proto $|OA| = |OX|$. Záměna $A \leftrightarrow B$, $D \leftrightarrow C$ zachovává všechny předpoklady (čtyřúhelník $BADC$ je týž konvexní čtyřúhelník s~vrcholy vypsanými v~opačném pořadí a~pětice úhlů ze zadání přejde sama na sebe) i~body $P$, $X$, $O$, takže stejnou úvahou je osa úsečky $PD$ osou úsečky $BX$ a~$|OB| = |OX|$.
}

\EnvId{aI08awqy1yrLlCuEOFHJ7-pata-vysky-a-stred-oh}
\Problem{0}{Česko-slovenská olympiáda 2011}{
    Nechť $ABC$ je ostroúhlý trojúhelník, který není rovnostranný. Nechť~$P$ je pata výšky z~$A$ na stranu~$BC$, $H$ ortocentrum, $O$ střed opsané kružnice, $D$ průsečík polopřímky~$AO$ se stranou~$BC$ a~$E$ střed úsečky~$AD$. Dokažte, že přímka~$EP$ prochází středem úsečky~$OH$.
}{
    Situace je poměrně zamotaná a na první pohled toho zajímavého moc neobjevíme. Je v ní však bod~$H$, se kterým můžeme dělat všelijaké zajímavé věci. Která z nich se sem hodí?
}{
    Zrcadlete~$H$ podle~$BC$ a získáte~$H'$ ležící na kružnici opsané trojúhelníku~$ABC$. To znamená $|OA| = |OH'|$. Střed~$OH$ se také stane lépe uchopitelným.
}{
    Označme $\omega$ kružnici opsanou trojúhelníku $ABC$, $R$ její poloměr a~$N$ střed úsečky $OH$; trojúhelník není rovnostranný, takže $O \neq H$.

    Pokud $|AB| = |AC|$, leží na ose strany $BC$ kromě $O$ i body $A$, $H$ a~$P$, takže $D = P$ a přímka $EP$ je právě tato osa. Ta obsahuje $O$ i $H$, a tedy i střed úsečky $OH$. Dále proto předpokládejme $|AB| \neq |AC|$; pak přímka $AO$ není kolmá na $BC$.

    Jediný bod konfigurace, se kterým lze volně nakládat, je ortocentrum, a jeho obraz $H'$ v osové souměrnosti podle $BC$ udělá naráz dvě věci: podle tvrzení o zrcadlení ortocentra leží $H'$ na $\omega$, čili
    $$
    |OH'| = |OA| = R,
    $$
    a zároveň je $P$ středem úsečky $HH'$, čímž se do hry dostane i střed úsečky $OH$. Přímka $AP$ je totiž kolmá na $BC$, v této souměrnosti se tedy zobrazí sama na sebe a bod $P$ zůstane pevný; bod $H'$ proto leží na přímce $AP$ a $P$ je středem $HH'$.

    Trojúhelník je ostroúhlý, takže $H$ leží uvnitř něj, a tedy uvnitř úsečky $AP$. Bod $H'$ pak leží na polopřímce $AP$ za bodem $P$ a na výšce z~$A$ jdou body v pořadí $A$, $H$, $P$, $H'$. Speciálně $P$ je vnitřním bodem úsečky $AH'$ a $H'$ leží v opačné polorovině s hraniční přímkou $BC$ než $A$, takže $H' \neq A$.

    Nechť $A'$ je bod diametrálně protilehlý k~$A$. Kdyby bylo $H' = A'$, přímka $AA'$ by splývala s výškou z~$A$ a byla by kolmá na $BC$, což jsme vyloučili. Je tedy $H' \neq A'$ a z Thaletovy věty nad průměrem $AA'$ dostáváme $|\angle AH'A'| = 90^\circ$. Přímka $H'A'$ je proto kolmá na $AH'$, a jelikož $AH' \perp BC$, je
    $$
    H'A' \parallel BC.
    $$

    Uvažujme stejnolehlost $h$ se středem $A$, která zobrazuje $P$ na $H'$; její koeficient $|AH'| \, / \, |AP|$ je kladný, neboť $P$ leží uvnitř úsečky $AH'$. Přímku $BC$ zobrazí $h$ na rovnoběžku s~$BC$ procházející bodem $H'$, což je právě přímka $H'A'$, a přímku $AD$ zobrazí samu na sebe. Přímka $AD$ není rovnoběžná s~$BC$ (protíná ji v bodě $D$), takže s přímkou $H'A'$ má jediný společný bod, a tím je $A'$. Proto $h(D) = A'$, a jelikož stejnolehlost zachovává středy úseček, zobrazí se střed $E$ úsečky $AD$ na střed úsečky $AA'$, čili na $O$:
    $$
    h(P) = H', \qquad h(D) = A', \qquad h(E) = O.
    $$

    Bod $E$ je vnitřním bodem úsečky $AD$, neleží tedy na $BC$ a $E \neq P$. Obrazem přímky $EP$ v~$h$ je přímka $OH'$ a stejnolehlost zobrazuje přímku na přímku s ní rovnoběžnou (případně na ni samu). Přímka $EP$ je tedy rovnoběžná s~$OH'$, nebo s ní totožná.

    S bodem $N$ naložíme stejně. Stejnolehlost se středem $H$ a koeficientem $\frac12$ zobrazí $H'$ na střed úsečky $HH'$, čili na $P$, a bod $O$ na střed úsečky $HO$, čili na $N$. Z~$|OH'| = R > 0$ plyne $O \neq H'$, a tedy i $P \neq N$; přímka $PN$ je obrazem přímky $OH'$, takže je s ní rovnoběžná, nebo s ní totožná.

    Přímky $EP$ a~$PN$ procházejí obě bodem $P$ a obě jsou rovnoběžné s přímkou $OH'$, nebo s ní totožné. Taková přímka je bodem $P$ určena jednoznačně, takže $EP$ a~$PN$ splývají a bod $N$, čili střed úsečky $OH$, leží na přímce $EP$.
}

\EnvId{M3FyEsmtSGx30V-INIVXw-tecna-a-stred-usecky-ap}
\Problem{0}{}{
    Nechť $ABC$ je ostroúhlý trojúhelník vepsaný do kružnice~$k$. Tečna ke~$k$ v bodě~$A$ protíná přímku~$BC$ v~$P$. Nechť~$M$ je střed~$AP$ a nechť~$Q$ je druhý průsečík přímky~$MB$ s~$k$. Dokažte, že $|\angle PQA| = |\angle AQC|$.
}{
    Střed~$AP$ je zvláštní bod, který tak úplně nesedí. Ledaže bychom ho použili ke středové souměrnosti a přenesli tak jeden z úhlů, jejichž rovnost dokazujeme.
}{
    Dává smysl zrcadlit~$Q$ podle~$M$ a přenést úhel~$PQA$; ten se najednou rovná~$PQ'A$. Dokazované tvrzení nám pak po krátkém počítání úhlů řekne, co stačí ukázat.
}{
    Bod $M$ vstupuje do zadání jedině jako střed úsečky $AP$, a~střed úsečky volá po středové souměrnosti podle něj: ta vyrobí rovnoběžník. Označme $Q'$ obraz bodu $Q$ v~této souměrnosti; úhel $PQA$ se jí přenese do vrcholu $Q'$.

    Označme $t$ tečnu kružnice $k$ v~bodě $A$, na které leží $P$ i~$M$. Přitom $P \neq A$, neboť $A$ neleží na přímce $BC$; proto je i~$M \neq A$ a~$M \neq P$. Přímka $t$ protíná přímku $AB$ jedině v~$A$ a~přímku $BC$ jedině v~$P$, takže $M$ neleží ani na jedné z~nich; jelikož $Q$ leží na přímce $MB$, dostáváme $Q \neq A$ i~$Q \neq C$. Bod $Q$ proto neleží na $t$ (ta má s~$k$ společný jediný bod $A$), a~jelikož $M$ je středem úsečky $AP$ i~úsečky $QQ'$, je $AQPQ'$ rovnoběžník. Jeho protilehlé úhly dávají
    $$
    |\angle PQA| = |\angle AQ'P|.
    $$

    Jelikož $M$ leží na $t$ a~$M \neq A$, leží vně $k$ a~jeho mocnost vzhledem ke $k$ je $|MA|^2 > 0$. Body $B$, $Q$ tedy leží na téže polopřímce z~$M$ a~$|MB| \cdot |MQ| = |MA|^2$. Souměrnost podle $M$ zachovává $|MQ'| = |MQ|$ a~přenáší $Q'$ na opačnou polopřímku, takže $M$ leží uvnitř úsečky $BQ'$; uvnitř úsečky $AP$ leží také a~$|MA| = |MP|$, čili
    $$
    |MA| \cdot |MP| = |MB| \cdot |MQ'|.
    $$
    Podle kritéria tětivovosti leží $A$, $B$, $P$, $Q'$ na jedné kružnici. (Jsou to čtyři různé body: z~$Q' = A$ by vyšlo $Q = P$, který na $k$ neleží, z~$Q' = P$ zase $Q = A$, a~$Q' \neq B$, neboť $M$ leží mezi nimi. Body $A$, $B$, $P$ nejsou kolineární, neboť $B$ neleží na $t$.) Úhlopříčky čtyřúhelníku $ABPQ'$ se protínají v~jeho vnitřním bodě $M$, takže je konvexní a~jeho protilehlé úhly při vrcholech $B$ a~$Q'$ dávají
    $$
    |\angle AQ'P| = 180^\circ - |\angle ABP|.
    $$
    Zbývá nám tedy dokázat
    $$
    |\angle AQC| + |\angle ABP| = 180^\circ. \eqno(*)
    $$

    Bod $P$ leží vně $k$, čili mimo tětivu $BC$: buď leží $B$ mezi $P$ a~$C$, nebo $C$ mezi $P$ a~$B$. Na přímce $AC$ bod $P$ neleží, neboť ta protíná přímku $BC$ jedině v~$C$ a~$P \neq C$; body $A$, $P$, $C$ tedy tvoří trojúhelník, $M$ je vnitřním bodem jeho strany $AP$ a~přímka $BQ$ neprochází žádným jeho vrcholem (přes $A$ nebo $P$ by to byla přímka $t$, na které $B$ neleží, přes $C$ zase přímka $BC$, na které neleží $M$). Přímka, která vejde do trojúhelníku vnitřním bodem jedné strany a~vyhne se vrcholům, protne vnitřek právě jedné ze zbývajících dvou stran. S~přímkou $PC$ má přímka $BQ$ společný jediný bod $B$: pokud tedy $B$ leží uvnitř úsečky $PC$, je to právě ta protnutá strana a~přímka $BQ$ nechá $A$ i~$C$ v~téže polorovině; pokud $B$ uvnitř úsečky $PC$ neleží, tedy pokud $C$ leží mezi $P$ a~$B$, protne přímka $BQ$ vnitřek úsečky $AC$ a~body $A$, $C$ oddělí.

    Dvě tětivy kružnice se protínají v~jejím vnitřním bodě právě tehdy, když každá z~nich odděluje krajní body té druhé. V~prvním případě proto leží $Q$ a~$B$ v~téže polorovině s~hraniční přímkou $AC$, čili $|\angle AQC| = |\angle ABC|$, a~polopřímky $BP$, $BC$ jsou opačné, čili $|\angle ABP| = 180^\circ - |\angle ABC|$. Ve druhém případě přímka $AC$ body $Q$ a~$B$ odděluje, čili $|\angle AQC| = 180^\circ - |\angle ABC|$, a~polopřímky $BP$, $BC$ splývají, čili $|\angle ABP| = |\angle ABC|$. Tak či onak platí $(*)$, a~tedy
    $$
    |\angle PQA| = |\angle AQ'P| = 180^\circ - |\angle ABP| = |\angle AQC|.
    $$
}

\EnvId{JDPS5QjvQROqZ2N6JEwq6-lichobeznik-a-body-na-kl}
\Problem{0}{}{
    Nechť $ABCD$ je lichoběžník s $AB \parallel CD$. Body $K$ a~$L$ leží po řadě na~$AB$ a~$CD$ tak, že $|AK| \cdot |CL| = |BK| \cdot |DL|$. Body $P$ a~$Q$ leží na úsečce~$KL$ tak, že $|\angle APB| = |\angle DCB|$ a $|\angle CQD|=|\angle CBA|$. Dokažte, že $B$, $C$, $P$, $Q$ leží na jedné kružnici.
}{
    Podmínka na polohy~$K$ a~$L$ je zajímavá. Její přepis do tvaru poměru napovídá, že v tom hraje roli nějaká stejnolehlost. Co takhle v ní něco dalšího zrcadlit?
}{
    Uvažte stejnolehlost zobrazující~$AB$ na $DC$; při ní se~$K$ zobrazí na~$L$. Zobrazíme-li při ní~$P$ na~$P'$, pak body $K$, $L$, $P$, $Q$, $P'$ všechny leží na jedné přímce. Úloha se promění v pěkné cvičení na počítání úhlů.
}{
    Body $K$ a~$L$ bereme jako vnitřní body stran $AB$ a~$CD$. Kdyby byl $K$ krajním bodem strany $AB$, zadaná rovnost by vynutila, že $L$ je příslušný krajní bod strany $CD$, a~úsečka $KL$ by splynula s~ramenem lichoběžníku. Na rameni $AD$ však pro každý jeho bod $P$ platí $|\angle APB| < 180^\circ - |\angle DAB| = |\angle ADC|$ a~pro každý jeho bod $Q$ platí $|\angle CQD| < |\angle DAB|$, takže předepsané rovnosti by naráz dávaly $|\angle DCB| < |\angle ADC|$ i~$|\angle ADC| < |\angle DCB|$. Na rameni $BC$ nevyhovuje dokonce ani jeden z~těch dvou bodů, neboť tam je $|\angle APB| < |\angle DCB|$ a~$|\angle CQD| < |\angle CBA|$. Taková poloha tedy do zadání nespadá. Zadanou podmínku pak umíme přepsat na poměr
    $$
    \frac{|AK|}{|KB|} = \frac{|DL|}{|LC|},
    $$
    který říká, že $K$ a~$L$ dělí rovnoběžné strany $AB$ a~$DC$ ve~stejném poměru. Přesně tak si odpovídají body ve~stejnolehlosti zobrazující $AB$ na $DC$, a~když nám rovnoběžné úsečky takovou stejnolehlost nabízejí, vyplatí se jí zobrazit další bod. Zobrazíme jí bod $P$.

    Jelikož $ABCD$ je lichoběžník, ramena $AD$ a~$BC$ nejsou rovnoběžná a~jejich přímky se protínají v~bodě $S$. Kdyby $S$ ležel uvnitř úsečky $AD$, ležel by striktně mezi rovnoběžkami $AB$ a~$CD$, a~tedy i~uvnitř úsečky $BC$, čili protilehlé strany lichoběžníku by se protínaly. Bod $S$ proto na úsečce $AD$ neleží: body $A$, $D$ leží na téže polopřímce vycházející z~bodu $S$ a~úsečka $AD$ neprotíná přímku $BC$, takže $A$ a~$D$ leží v~téže polorovině určené přímkou $BC$.

    Nechť $h$ je stejnolehlost se středem $S$ a~koeficientem $k = |SD| \, / \, |SA| > 0$, takže $h(A) = D$. Přímka $AB$ se při ní zobrazí na rovnoběžku s~$AB$ procházející bodem $D$, čili na přímku $CD$, a~přímka $SB = BC$ je samodružná; obraz $h(B)$ je tedy průsečík přímek $CD$ a~$BC$, čili $h(B) = C$. Lichoběžník není rovnoběžník, takže $|DC| \neq |AB|$ a~$k \neq 1$. Bod $S$ navíc neleží ani na jedné z~rovnoběžek $AB$, $CD$: je samodružný a~$h$ zobrazuje $AB$ na $CD$, takže by ležel na obou.

    Stejnolehlost zachovává dělicí poměry, takže $h(K)$ je ten bod úsečky $DC$, který ji dělí v~poměru $|AK| : |KB|$. Podle přepsané podmínky je to právě $L$, čili
    $$
    h(K) = L.
    $$
    Bod $L$ tak leží na přímce $SK$: přímka $\ell = KL$ prochází středem $S$, je při $h$ samodružná a~leží na ní i~obraz $P' = h(P)$.

    Koeficient $k$ je kladný, takže $h$ zobrazuje každou polopřímku z~bodu $S$ samu na sebe. Body $K$, $L = h(K)$ a~$h(L)$ proto leží na téže polopřímce z~$S$ a~jejich vzdálenosti od $S$ jsou po řadě $|SK|$, $k|SK|$ a~$k^2|SK|$. Pro $k > 1$ tato trojice roste, pro $k < 1$ klesá; v~obou případech je prostřední hodnota striktně mezi krajními, takže $L$ leží mezi $K$ a~$h(L)$. Jelikož $K$ i~$L$ leží na téže polopřímce z~$S$ a~jsou od $S$ různé, bod $S$ na úsečce $KL$ neleží.

    Jelikož $0 < |\angle APB| = |\angle DCB| < 180^\circ$, bod $P$ neleží na přímce $AB$, speciálně $P \neq K$. Obrazem úsečky $KL$ je úsečka $L\,h(L)$, takže $P'$ leží na ní a~$P' \neq L$. Bod $L$ tedy odděluje $P'$ od $K$, a~jelikož $\ell$ protíná přímku $CD$ jedině v~$L$, leží $P'$ a~$K$ (a~s~ním celá přímka $AB$) v~opačných polorovinách určených přímkou $CD$.

    Stejnolehlost $h$ zobrazuje trojúhelník $APB$ na trojúhelník $DP'C$, takže
    $$
    |\angle DP'C| = |\angle APB| = |\angle DCB|.
    $$
    Bod $P'$ leží na $\ell$ a~$P' \neq L$, takže neleží na přímce $CD$; nechť $\gamma$ je kružnice opsaná trojúhelníku $DP'C$. Podle věty o~úsekovém úhlu svírá tečna ke $\gamma$ v~bodě $C$ s~polopřímkou $CD$ v~polorovině neobsahující $P'$ úhel velikosti $|\angle DP'C|$. Polopřímka $CB$ leží v~téže polorovině a~s~polopřímkou $CD$ svírá úhel $|\angle DCB| = |\angle DP'C|$; polopřímka z~bodu $C$ je však v~dané polorovině velikostí úhlu s~$CD$ určena jednoznačně, takže přímka $BC$ se dotýká $\gamma$ v~bodě $C$.

    Přímka $BC$ je příčkou rovnoběžek $AB$ a~$CD$ a~body $A$, $D$ leží u~ní na téže straně, takže $|\angle CBA| + |\angle BCD| = 180^\circ$, čili
    $$
    |\angle DQC| = |\angle CBA| = 180^\circ - |\angle DCB| = 180^\circ - |\angle DP'C|.
    $$
    Jelikož $0 < |\angle DQC| < 180^\circ$, bod $Q$ neleží na přímce $CD$, čili $Q \neq L$, a~jelikož úsečka $KL$ protíná přímku $CD$ jedině v~$L$, leží $Q$ v~téže polorovině jako $K$, tedy na opačné straně přímky $CD$ než $P'$. Oblouk kružnice $\gamma$ ležící v~této polorovině je přitom množinou právě těch jejích bodů, z~nichž je $DC$ vidět pod úhlem $180^\circ - |\angle DP'C|$ (obvodové úhly z~opačných oblouků se doplňují do $180^\circ$), a~taková množina je jediná. Proto $Q$ leží na $\gamma$.

    Přímky $\ell$ a~$BC$ jsou různé, neboť jinak by $K$ byl průsečíkem přímek $AB$ a~$BC$, čili vrcholem $B$; obě procházejí bodem $S$, takže $\ell \cap BC = \{S\}$, a~$S$ na úsečce $KL$ neleží. Body $P$, $Q$ proto neleží na přímce $BC$ a~body $B$, $C$, které jsou od $S$ různé, neleží na $\ell$. Na $\ell$ neleží ani $D$, jelikož $\ell$ protíná přímku $CD$ jedině ve vnitřním bodě $L$ strany $CD$. Je-li $P = Q$, tvrzení platí (body $B$, $C$, $P$ nejsou kolineární, stačí vzít kružnici opsanou trojúhelníku $BCP$); nechť tedy $P \neq Q$, takže $B$, $C$, $P$, $Q$ jsou navzájem různé.

    Dále počítáme s~orientovanými úhly přímek modulo $180^\circ$, čímž se vyhneme rozboru vzájemných poloh bodů $P$, $Q$, $P'$ na přímce $\ell$; konfigurační fakta, na nichž úloha doopravdy stojí, jsme už použili. Stejnolehlost $h$ zobrazuje přímku $PB$ na přímku $P'C$, čili $PB \parallel P'C$, a~body $P$, $Q$, $P'$ leží na $\ell$. Proto
    $$
    \angle(PB, PQ) = \angle(P'C, P'Q) = \angle(DC, DQ) = \angle(CB, CQ),
    $$
    kde druhá rovnost jsou obvodové úhly nad tětivou $CQ$ kružnice $\gamma$ z~bodů $P'$ a~$D$ a~třetí je úsekový úhel pro tečnu $CB$ a~tutéž tětivu. Rovnost $\angle(PB,PQ) = \angle(CB,CQ)$ znamená, že body $B$, $C$, $P$, $Q$ leží na jedné kružnici nebo na jedné přímce; přímka je vyloučena, neboť $B$ neleží na $\ell$.
}

\EnvId{Kr9IKyGuib12QPGum8wz1-ctyri-body-na-stranach-ctyruhelniku}
\Problem{0}{CPSJ 2019, soutěž družstev}{
    Nechť $ABCD$ je tětivový čtyřúhelník. Body $K$, $L$, $M$, $N$ leží po řadě na stranách~$AB$, $BC$, $CD$, $DA$ tak, že $|\angle ADK| = |\angle BCK|$, $|\angle BAL| = |\angle CDL|$, $|\angle CBM| = |\angle DAM|$ a $|\angle DCN|=|\angle ABN|$. Dokažte, že $KM \perp LN$.
}{
    Narýsujte jen jeden bod, řekněme~$K$, a zkuste lépe pochopit zvláštní úhlovou podmínku, kterou je definován. Přidávání dalších bodů začne dávat smysl teprve tehdy, až něco objevíte.
}{
    V obrázku jen s bodem~$K$ lze počítáním úhlů dostat $|\angle BKC|=|\angle KDC|$. Přímka~$AB$ je tedy tečnou ke kružnici opsané trojúhelníku~$KCD$. To připomíná mocnost bodu: stačí jen přidat bod, ze kterého ji počítáme. Jakmile je tohle jasné, můžeme dopojit zbývající body~$L$, $M$, $N$.
}{
    Body $A$, $B$, $C$, $D$ leží v tomto pořadí na kružnici $\omega$, takže čtyřúhelník $ABCD$ je konvexní. Všechny čtyři podmínky ze zadání jsou přitom jedna a táž: bod $X$ leží na straně $A_1A_2$ čtyřúhelníku $A_1A_2A_3A_4$ a platí $|\angle A_1A_4X| = |\angle A_2A_3X|$, postupně pro
    $$
    (A_1,A_2,A_3,A_4) = (A,B,C,D),\ (B,C,D,A),\ (C,D,A,B),\ (D,A,B,C).
    $$
    Cyklický posun značení $(A,B,C,D) \to (B,C,D,A)$ tedy zachovává zadání a čtveřici $(K,L,M,N)$ převede na $(L,M,N,K)$. Každé tvrzení o bodě $K$ tak stačí dokázat jen jednou.

    Když se $X$ pohybuje po straně $AB$ z~$A$ do~$B$, úhel $\angle ADX$ roste od nuly po $|\angle ADB|$ a úhel $\angle BCX$ klesá od $|\angle BCA|$ k nule. Bod $K$ proto existuje, je jediný a je vnitřním bodem strany $AB$; totéž platí pro $L$, $M$, $N$.

    Vyřiďme si nejprve případ, kdy jsou některé dvě protilehlé strany rovnoběžné; díky posunu můžeme předpokládat $AB \parallel CD$. Pak je $ABCD$ rovnoramenný lichoběžník a osa $s$ úsečky $AB$ je zároveň osou úsečky $CD$. Osová souměrnost podle~$s$ zamění $A \leftrightarrow B$ a~$C \leftrightarrow D$, čili zachová konfiguraci s přeznačením $(A,B,C,D) \to (B,A,D,C)$, při němž přejde $(K,L,M,N)$ na $(K,N,M,L)$. Z jednoznačnosti tedy tato souměrnost zobrazí $K$ na $K$, $M$ na $M$ a~$L$ na $N$. Body $K$, $M$ proto leží na~$s$ a jsou různé, takže $KM = s$, kdežto $LN$ je na~$s$ kolmá. Dále už předpokládejme, že žádné dvě protilehlé strany rovnoběžné nejsou.

    Klíčové pozorování lze udělat v obrázku s jediným novým bodem. Označme $\varphi = |\angle ADK| = |\angle BCK|$, $\beta = |\angle ABC|$ a~$\delta = |\angle CDA|$. Polopřímka $BK$ splývá s polopřímkou $BA$, takže ze součtu úhlů v trojúhelníku $BCK$ a z~$\beta + \delta = 180^\circ$ dostaneme
    $$
    |\angle BKC| = 180^\circ - \beta - \varphi = \delta - \varphi.
    $$
    Bod $K$ je vnitřním bodem strany $AB$ konvexního čtyřúhelníku, takže polopřímka $DK$ leží uvnitř úhlu $ADC$, a tedy
    $$
    |\angle KDC| = \delta - \varphi = |\angle BKC|.
    $$
    Úsečka $KC$ dělí čtyřúhelník na trojúhelník $KBC$ a čtyřúhelník $AKCD$, čili body $B$ a~$D$ leží v opačných polorovinách přímky $KC$. Podle obrácené věty o úsekovém úhlu je proto přímka $AB$ tečnou kružnice opsané trojúhelníku $KCD$ v bodě $K$.

    Tečna je nám k ničemu, dokud nemáme bod, ze kterého mocnost počítáme. Stačí protnout dvě vhodné přímky: nechť $P$ je průsečík přímek $AB$ a~$CD$. Jelikož je $ABCD$ konvexní, leží $P$ mimo úsečky $AB$ i~$CD$, tedy vně $\omega$, a mocnost bodu $P$ ke kružnici $(KCD)$ jednou přes tečnu $PK$ a jednou přes sečnu $CD$ dává
    $$
    |PK|^2 = |PC| \cdot |PD| = |PA| \cdot |PB|,
    $$
    kde druhá rovnost je mocnost bodu $P$ k~$\omega$. Dvojnásobný posun $(A,B,C,D) \to (C,D,A,B)$ převede $K$ na $M$ a bod $P$ na průsečík přímek $CD$, $AB$, čili zase na $P$. Proto $|PM|^2 = |PA| \cdot |PB|$, a tedy
    $$
    |PK| = |PM|.
    $$
    Po jednom, resp. třech posunech dostaneme stejně $|QL| = |QN|$, kde $Q$ je průsečík přímek $BC$ a~$DA$.

    Body $K$, $M$ jsou různé, protože leží na různých přímkách $AB$, $CD$, jejichž jediný společný bod $P$ není vnitřním bodem žádné strany; stejně $L \neq N$. Označme $p$ osu úsečky $KM$ a~$q$ osu úsečky $LN$. Z~$|PK| = |PM|$ leží $P$ na~$p$, takže osová souměrnost podle~$p$ nechá $P$ na místě a zamění $K$ s~$M$, čímž zobrazí přímku $AB$ na přímku $CD$. Stejně tak souměrnost podle~$q$ zobrazí $BC$ na $DA$.

    Zbývá porovnat směry os $p$ a~$q$. Zde přejdeme na orientované úhly přímek modulo $180^\circ$, zapisované $\angle(g,h)$; výsledek vyjde nulový, takže rozbor vzájemných poloh úplně odpadne. Pro osovou souměrnost podle přímky $p$ a libovolnou přímku $g$ s obrazem $g'$ platí $\angle(g,p) = \angle(p,g')$, což pro $g = AB$, $g' = CD$ dává $\angle(AB,CD) = 2\angle(AB,p)$, a stejně $\angle(BC,DA) = 2\angle(BC,q)$. Proto
    $$
    \eqalign{
        2\angle(p,q) &= 2\angle(p,AB) + 2\angle(AB,BC) + 2\angle(BC,q) \cr
        &= -\angle(AB,CD) + 2\angle(AB,BC) + \angle(BC,DA) \cr
        &= \angle(AB,DA) - \angle(BC,CD). \cr
    }
    $$
    Tětivovost čtyřúhelníku $ABCD$ znamená právě $\angle(AB,DA) = \angle(CB,CD)$, takže pravá strana je nula a~$\angle(p,q)$ je $0^\circ$ nebo $90^\circ$. Jelikož $KM \perp p$ a~$LN \perp q$, jsou přímky $KM$ a~$LN$ rovnoběžné nebo kolmé.

    Rovnoběžnost vyloučíme polohou bodů. Průnikem přímky $KM$ s konvexním čtyřúhelníkem $ABCD$ je úsečka $KM$ a její vnitřní body leží uvnitř čtyřúhelníku, protože $K$, $M$ jsou vnitřními body protilehlých stran. Bod $L$ leží na hranici čtyřúhelníku a je různý od $K$ i~$M$, takže na přímce $KM$ neleží; stejně $N$. Úsečka $KM$ přitom dělí čtyřúhelník na $KBCM$ a~$AKMD$, přičemž $L$ je vnitřním bodem strany $BC$ prvního z nich a~$N$ vnitřním bodem strany $DA$ druhého. Body $L$ a~$N$ tak leží v opačných otevřených polorovinách přímky $KM$, čili přímky $KM$ a~$LN$ rovnoběžné nejsou. Jsou tedy kolmé.
}

\EnvId{t8rjDz8xH7IC_FmnHz15z-ctyruhelnik-se-tremi-shodnymi-stranami}
\Problem{0}{CAPS 2024, P4}{
    Nechť $ABCD$ je čtyřúhelník takový, že $|AB| = |BC| = |CD|$. Na polopřímkách $CA$, $BD$ jsou po řadě body $X$, $Y$ takové, že $|BX| = |CY|$. Nechť $P$, $Q$, $R$, $S$ jsou po řadě středy úseček $BX$, $CY$, $XD$, $YA$. Dokažte, že body $P$, $Q$, $R$, $S$ leží na jedné kružnici.
}{
    Máme spoustu středů. Zrcadlíme, nebo zavádíme ještě další? Zde dává větší smysl to druhé.
}{
    Klíčovým bodem, který zavést, je střed~$XY$, neboť vytvoří spoustu středních příček. A ty dají jasnou cestu, co použít k důkazu toho, že body leží na kružnici.
}{
    Body $P$, $R$ jsou středy stran $XB$, $XD$ trojúhelníku $XBD$, takže $PR \parallel BD$; stejně tak $Q$, $S$ jsou středy stran $YC$, $YA$ trojúhelníku $YCA$, takže $QS \parallel CA$ (z~$B \neq D$ a~$A \neq C$ plyne $P \neq R$ a~$Q \neq S$, obě přímky jsou tedy dobře definované). Samotné rovnoběžnosti o kružnici neříkají nic, chybí nám průsečík přímek $PR$, $QS$ a délky odměřené z něj. Jelikož v úloze vystupují samé středy úseček, přidejme ještě jeden, který s každým z nich vytvoří střední příčku: střed~$M$ úsečky $XY$.

    Potom je $MP$ střední příčka trojúhelníku $XBY$, $MR$ trojúhelníku $XDY$, $MQ$ trojúhelníku $YCX$ a $MS$ trojúhelníku $YAX$. Všechny čtyři lze zapsat naráz: chápeme-li body jako polohové vektory, dostaneme z~$M = \tfrac12(X+Y)$ a z definic bodů $P$, $Q$, $R$, $S$
    $$
    \gather
    P - M = \tfrac12 (B - Y), \qquad R - M = \tfrac12 (D - Y), \\
    Q - M = \tfrac12 (C - X), \qquad S - M = \tfrac12 (A - X).
    \endgather
    $$

    Bod $Y$ leží na přímce $BD$, takže $P - M$ i $R - M$ jsou rovnoběžné s~$BD$ a $M$ leží na přímce $PR$; stejně tak $X$ leží na přímce $CA$, takže $M$ leží na přímce $QS$. Úhlopříčky $AC$ a~$BD$ čtyřúhelníku $ABCD$ nejsou rovnoběžné, proto jsou $PR$ a~$QS$ různé přímky a jejich jediným společným bodem je $M$.

    Dále počítejme s orientovanými délkami na přímkách: pro kolineární body součin $\overline{UV} \cdot \overline{UW}$ nezávisí na volbě orientace a ušetří nám rozbor toho, zda $X$ leží uvnitř úsečky $AC$, nebo mimo ni. Zvolíme-li na přímce $PR$ tutéž orientaci jako na $BD$ a na $QS$ tutéž jako na $CA$, dávají předchozí rovnosti
    $$
    \gather
    \overline{MP} \cdot \overline{MR} = \tfrac14\, \overline{YB} \cdot \overline{YD}, \\
    \overline{MQ} \cdot \overline{MS} = \tfrac14\, \overline{XC} \cdot \overline{XA}.
    \endgather
    $$
    Podle kritéria tětivovosti tak stačí dokázat rovnost
    $$
    \overline{YB} \cdot \overline{YD} = \overline{XC} \cdot \overline{XA};
    $$
    je-li společná hodnota nenulová, jsou body $P$, $Q$, $R$, $S$ všechny různé od $M$ a kritérium dává hledanou kružnici.

    Na obou stranách stojí mocnost bodu, stačí jen najít ty správné kružnice. Označme $s = |AB| = |BC| = |CD|$. Rovnost $|BA| = |BC| = s$ říká přesně to, že kružnice $\omega_B$ se středem $B$ a poloměrem $s$ prochází body $A$, $C$; stejně tak $|CB| = |CD| = s$ říká, že kružnice $\omega_C$ se středem $C$ a poloměrem $s$ prochází body $B$, $D$. Přímka $CA$ protíná $\omega_B$ právě v bodech $A$ a~$C$ a bod $X$ na ní leží, takže
    $$
    \overline{XC} \cdot \overline{XA} = |BX|^2 - s^2,
    $$
    a stejně tak $\overline{YB} \cdot \overline{YD} = |CY|^2 - s^2$. Podmínka $|BX| = |CY|$ tyto dvě mocnosti ztotožňuje, čímž je požadovaná rovnost dokázána.

    Zbývá případ, kdy je společná hodnota $|BX|^2 - s^2$ nulová, tedy $|BX| = |CY| = s$. Pak $X$ leží na $\omega_B$, čili $X \in \{A, C\}$, a $Y$ leží na $\omega_C$, čili $Y \in \{B, D\}$. V každé ze čtyř takových možností dva ze čtyř bodů splynou a zbylé tři jsou středy tří po sobě jdoucích stran čtyřúhelníku $ABCD$: například pro $X = A$, $Y = B$ je $S = P$ střed strany $AB$ a zbylé dva body jsou středy stran $DA$ a~$BC$. Jsou to tři po sobě jdoucí vrcholy Varignonova rovnoběžníku, jehož sousední strany jsou rovnoběžné s~$AC$ a~$BD$; ty rovnoběžné nejsou, takže naše tři body neleží na jedné přímce a~kružnice jimi procházející existuje.
}

\EnvId{unjF6AEffP5thl4GWrGD5-uhlopricky-lichobezniku-a-vnitrni-bod}
\Problem{0}{}{
    Úhlopříčky lichoběžníku~$ABCD$ s $AB \parallel CD$ se protínají v~$P$. Bod~$Q$ leží uvnitř trojúhelníku~$BCP$ tak, že $|\angle AQB|= |\angle CQD|$. Dokažte, že $|\angle DQP| = |\angle BAQ|$.
}{
    Bod~$Q$ je definován dost libovolně. Nezbývá než přidat nějaký další bod, který přinese smysluplnou strukturu. Sám lichoběžník dává docela dobrou nápovědu.
}{
    Zrcadlete~$Q$ ve stejnolehlosti se středem~$P$, která zobrazuje~$AB$ na~$CD$. Najednou se definice~$Q$ promění v něco smysluplného a odnikud máme vše, co potřebujeme; jen to poskládáme dohromady.
}{
    Rovnoběžné strany $AB$, $CD$ nabízejí stejnolehlost se středem $P$, která jednu z nich zobrazí na druhou. Podmínka $|\angle AQB| = |\angle CQD|$ svazuje dvojici $A$, $B$ s dvojicí $C$, $D$ přesně tak, jak to dělá tato stejnolehlost, zobrazme jí tedy i bod $Q$.

    Lichoběžník $ABCD$ je konvexní, takže $P$ je vnitřním bodem obou úseček $AC$, $BD$. Z rovnoběžnosti $AB \parallel CD$ jsou trojúhelníky $PAB$ a $PCD$ podobné, čili
    $$
    \frac{|PC|}{|PA|} = \frac{|PD|}{|PB|}.
    $$
    Stejnolehlost $h$ se středem $P$ a koeficientem $k = -|PC| \, / \, |PA| < 0$ proto zobrazuje $A \mapsto C$ a $B \mapsto D$ (záporné znaménko odpovídá tomu, že $P$ leží uvnitř obou úhlopříček). Označme $Q_1 = h(Q)$.

    Stejnolehlost zachovává velikosti úhlů, takže
    $$
    |\angle CQ_1D| = |\angle AQB| = |\angle CQD|.
    $$

    Abychom z toho dostali tětivový čtyřúhelník, potřebujeme ještě, že $Q$ a~$Q_1$ leží na téže straně přímky $CD$. Nechť $\pi$ je otevřená polorovina s hraniční přímkou $CD$ obsahující bod $A$. Bod $B$ leží na přímce rovnoběžné s~$CD$ a od ní různé, takže $B \in \pi$, a $P$ je vnitřním bodem úsečky $AC$ s $A \in \pi$, takže také $P \in \pi$. Vnitřek trojúhelníku $BCP$, jehož jediný vrchol na přímce $CD$ je $C$, tedy celý leží v~$\pi$; speciálně $Q \in \pi$. Jelikož $h(B) = D$, $h(C) = A$ a $h(P) = P$, zobrazí se trojúhelník $BCP$ na trojúhelník $DAP$ a $Q_1$ leží uvnitř něj. Stejnou úvahou (z jeho vrcholů leží na přímce $CD$ jedině $D$) je i $Q_1 \in \pi$.

    Body $Q$, $Q_1$ tedy leží na téže straně přímky $CD$ a vidí úsečku $CD$ pod stejným úhlem, takže podle obrácené věty o obvodovém úhlu leží $C$, $D$, $Q$, $Q_1$ na jedné kružnici; označme ji $\omega$.

    Koeficient $k$ je záporný a $Q \neq P$, takže $P$ je vnitřním bodem úsečky $QQ_1$. Polopřímky $QP$ a~$QQ_1$ proto splývají a platí $|\angle DQP| = |\angle DQQ_1|$.

    Bod $P$ je vnitřním bodem tětivy $QQ_1$, čili leží uvnitř $\omega$. Pro bod uvnitř kružnice se cyklické pořadí bodů na kružnici shoduje s cyklickým pořadím polopřímek, které k nim z něj vedou. Polopřímka $PC$ je opačná k~$PA$, polopřímka $PD$ je opačná k~$PB$ a přímka $AC$ odděluje $B$ od $D$, takže kolem $P$ jdou tyto čtyři polopřímky v pořadí $PA$, $PB$, $PC$, $PD$. Bod $Q$ leží uvnitř trojúhelníku $BCP$, takže polopřímka $PQ$ leží uvnitř úhlu $\angle BPC$ a opačná polopřímka $PQ_1$ uvnitř úhlu $\angle DPA$. Kolem $P$ tak polopřímky jdou v pořadí
    $$
    PA, \quad PB, \quad PQ, \quad PC, \quad PD, \quad PQ_1,
    $$
    a cyklické pořadí bodů na $\omega$ je $Q$, $C$, $D$, $Q_1$.

    Tětiva $DQ_1$ proto nechává body $C$ a~$Q$ na témže oblouku a z věty o obvodovém úhlu dostáváme $|\angle DQQ_1| = |\angle DCQ_1|$. Nakonec $h$ zobrazuje $A \mapsto C$, $B \mapsto D$, $Q \mapsto Q_1$, takže $|\angle BAQ| = |\angle DCQ_1|$. Spojením tří rovností
    $$
    |\angle DQP| = |\angle DQQ_1| = |\angle DCQ_1| = |\angle BAQ|.
    $$
}

\EnvId{cvVqE4nT1JdVqFMok0P_I-osy-usecek-bd-a-ce}
\Problem{0}{IMO 2018, P1}{
    Nechť $ABC$ je ostroúhlý trojúhelník s kružnicí opsanou $\Gamma$. Body $D$ a $E$ leží po řadě uvnitř stran $AB$ a $AC$, přičemž $|AD| = |AE|$. Osy úseček $BD$ a $CE$ protínají kratší oblouky $AB$ a $AC$ kružnice $\Gamma$ po řadě v bodech $F$ a $G$. Dokažte, že přímky $DE$ a $FG$ jsou rovnoběžné (nebo splývají).
}{
    Máme $|FB| = |FD|$, ale není jasné, jak to využít. Dává nám to stejné úhly $\angle FBD$ a $\angle FDB$. Zejména ten druhý vypadá osamoceně. Zkuste něco s něčím protnout a přenést tyto úhly někam jinam.
}{
    Klíčem je prodloužit úsečku $FD$ na tětivu $\Gamma$ a obdobně prodloužit $GE$ na druhé straně. Najednou dostaneme rovnoramenné trojúhelníky, které dobře souhrají s podmínkou $|AD| = |AE|$. Žádné další body nejsou potřeba; stane se z toho cvičení na počítání úhlů.
}{
    Označme $\alpha$, $\beta$, $\gamma$ vnitřní úhly trojúhelníku $ABC$. Trojúhelník je ostroúhlý, takže oblouk $AB$ neobsahující $C$ má velikost $2\gamma < 180^\circ$ a~je tedy ten kratší; bod $F$ proto leží na oblouku $AB$ neobsahujícím $C$ a~stejně tak $G$ leží na oblouku $AC$ neobsahujícím $B$.

    Z~$|FB| = |FD|$ máme $|\angle FDB| = |\angle FBD|$, jenže úhel při $D$ nemá na přímce $AB$ s~čím spolupracovat. Přeneseme ho na kružnici přesně v~duchu poslední ze situací v~seznamu výše: prodloužíme $FD$ na tětivu $\Gamma$. Nechť je tedy $T$ druhý průsečík přímky $FD$ s~$\Gamma$ a~nechť $S$ je druhý průsečík přímky $GE$ s~$\Gamma$.

    Bod $D$ je vnitřním bodem tětivy $AB$, leží tedy uvnitř $\Gamma$, a~proto na přímce $FT$ leží mezi $F$ a~$T$. Body $F$ a~$C$ leží v~opačných polorovinách určených přímkou $AB$; jelikož $D$ leží na přímce $AB$ mezi $F$ a~$T$, leží bod $T$ v~téže polorovině jako $C$, tedy na oblouku $AB$ obsahujícím $C$. Stejně tak $E$ leží mezi $G$ a~$S$ a~bod $S$ leží na oblouku $AC$ obsahujícím $B$.

    Označme $\varphi = |\angle ABF|$. Bod $D$ leží mezi $A$ a~$B$ i~mezi $F$ a~$T$, takže z~vrcholových úhlů
    $$
    |\angle ADT| = |\angle BDF| = |\angle DBF| = \varphi.
    $$
    Oblouk $AF$ neobsahující $B$ je částí oblouku $AB$ neobsahujícího $C$, takže body $B$ a~$T$ leží na témže z~oblouků určených tětivou $AF$ a~obvodové úhly nad ní dávají
    $$
    |\angle ATD| = |\angle ATF| = |\angle ABF| = \varphi.
    $$
    Trojúhelník $ADT$ je proto rovnoramenný a~$|AT| = |AD|$. Záměnou $B \leftrightarrow C$, $D \leftrightarrow E$, $F \leftrightarrow G$, $T \leftrightarrow S$ stejně tak vyjde $|AS| = |AE|$, a~jelikož $|AD| = |AE|$, dostáváme
    $$
    |AT| = |AS|.
    $$

    Polopřímky $AD$ a~$AB$ splývají, takže ze součtu úhlů v~trojúhelníku $ADT$ je $|\angle BAT| = 180^\circ - 2\varphi$. Bod $T$ leží na oblouku $AB$ obsahujícím $C$, tedy $|\angle ATB| = \gamma$, a~součet úhlů v~trojúhelníku $ABT$ dává
    $$
    |\angle ABT| = 180^\circ - (180^\circ - 2\varphi) - \gamma = 2\varphi - \gamma.
    $$
    Analogicky pro $\psi = |\angle ACG|$ vyjde $|\angle ACS| = 2\psi - \beta$.

    Úhly $\angle ABT$ a~$\angle ACS$ jsou obvodové úhly kružnice $\Gamma$ nad shodnými tětivami $AT$ a~$AS$. Shodné tětivy vytínají z~kružnice shodnou dvojici oblouků a~obvodový úhel nad tětivou je polovinou toho z~jejích dvou oblouků, který neobsahuje vrchol úhlu; úhly $\angle ABT$ a~$\angle ACS$ jsou proto shodné, nebo mají součet $180^\circ$.

    Ukážeme, že oba jsou ostré, čímž druhou možnost vyloučíme. Body $B$ a~$C$ leží na témže z~oblouků určených tětivou $AF$, takže $\varphi = |\angle ACF|$, a~jelikož úsečka $CF$ protíná tětivu $AB$, leží polopřímka $CF$ uvnitř úhlu $ACB$ a~$\varphi < \gamma$. Číslo $2\varphi - \gamma$ je úhel trojúhelníku $ABT$, tedy kladné, dohromady $0 < 2\varphi - \gamma < \gamma < 90^\circ$; stejně tak $0 < 2\psi - \beta < \beta < 90^\circ$. Součet dvou ostrých úhlů je menší než $180^\circ$, takže nastává první možnost:
    $$
    2\varphi - \gamma = 2\psi - \beta. \eqno(*)
    $$

    Nechť $M$ je střed oblouku $BC$ neobsahujícího $A$, tedy druhý průsečík osy úhlu $BAC$ s~$\Gamma$. Trojúhelník $ADE$ je rovnoramenný, osa úhlu $BAC$ je proto osou úsečky $DE$ a~platí $AM \perp DE$. Stačí tedy dokázat, že i~$AM \perp FG$.

    Na tom z~oblouků určených tětivou $FG$, který obsahuje $A$, leží pouze body oblouků $FA$ a~$AG$, zatímco $B$, $M$, $C$ leží na druhém. Tětiva $FG$ tak odděluje $A$ od $M$ a~úsečka $AM$ protíná $FG$ ve vnitřním bodě; označme ho $X$. Body $C$ a~$F$ leží na témže z~oblouků určených tětivou $AG$ (na tom, který obsahuje $B$), takže $|\angle AFG| = |\angle ACG| = \psi$. Dále $|\angle AFB| = 180^\circ - \gamma$, neboť $F$ a~$C$ leží na opačných obloucích určených tětivou $AB$, a~ze součtu úhlů v~trojúhelníku $ABF$ je $|\angle FAB| = \gamma - \varphi$. Polopřímka $AM$ jde vnitřkem úhlu $BAC$, zatímco $F$ leží v~opačné polorovině než $C$ vzhledem k~přímce $AB$; polopřímka $AB$ proto leží uvnitř úhlu $FAM$ a
    $$
    |\angle FAX| = |\angle FAB| + |\angle BAM| = \gamma - \varphi + \tfrac\alpha2.
    $$
    Ze vztahu $(*)$ je $\varphi - \psi = \frac{\gamma - \beta}2$, tedy $\gamma - \varphi + \psi = \frac{\beta + \gamma}2$, a~součet úhlů v~trojúhelníku $AFX$ dává
    $$
    \gather
    |\angle AXF| = 180^\circ - (\gamma - \varphi + \tfrac\alpha2) - \psi \\
    = 180^\circ - \tfrac{\beta + \gamma}2 - \tfrac\alpha2 = 90^\circ.
    \endgather
    $$
    Přímky $DE$ a~$FG$ jsou tak obě kolmé na přímku $AM$, jsou tedy rovnoběžné, nebo splývají.
}

\EnvId{Dk4xgZzPmhroFmRx3omVb-rovnobezne-tecny-k-vepsane-kruznici}
\Problem{0}{IMO 2024, P4}{
    Nechť $ABC$ je trojúhelník takový, že $|AB| < |AC| < |BC|$. Nechť $\omega$ je vepsaná kružnice trojúhelníku $ABC$ a $I$ její střed. Nechť $X$ je bod na přímce $BC$ různý od $C$ takový, že přímka procházející $X$ rovnoběžně s $AC$ se dotýká $\omega$. Podobně nechť $Y$ je bod na přímce $BC$ různý od $B$ takový, že přímka procházející $Y$ rovnoběžně s $AB$ se dotýká $\omega$. Přímka $AI$ protíná kružnici opsanou trojúhelníku $ABC$ podruhé v bodě $P \neq A$. Nechť $K$ a $L$ jsou po řadě středy stran $AC$ a $AB$. Dokažte, že $|\angle KIL| + |\angle YPX| = 180^\circ$.
}{
    S těmi rovnoběžnými přímkami máme skoro daný rovnoběžník.
}{
    Nechť~$Z$ je průsečík rovnoběžných přímek ze zadání; pak s využitím stejnolehlosti se středem~$A$ a koeficientem~$2$ máme $|\angle KIL| = |\angle BZC|$. Ty ošklivé středy už nemáme, neboť dokazujeme jen $|\angle BZC| + |\angle YPX| = 180^\circ$.
}{
    Klíčem k dotažení úlohy je najít tětivové čtyřúhelníky. Ale aspoň nemusíme zavádět další body.
}{
    Tečny ze zadání jsou rovnoběžné se stranami $AC$ a~$AB$, a~u rovnoběžek se vyplatí uvažovat o~zobrazení, které jednu z~nich pošle na druhou. Tady je takové po ruce: středová souměrnost podle $I$ zobrazuje $\omega$ samu na sebe, tedy tečnu na tečnu, a~každou přímku na rovnoběžku s~ní. Nechť $Z$ je obraz bodu $A$ v~této souměrnosti, tedy bod, pro který je $I$ středem úsečky $AZ$.

    Přímka vedená bodem $X$ rovnoběžně s~$AC$ je tečnou $\omega$ různou od přímky $AC$: kdyby s~$AC$ splývala, ležel by $X$ na přímkách $AC$ i~$BC$, čili by bylo $X = C$, což zadání vylučuje. Tečny $\omega$ daného směru jsou však jen dvě, takže tato přímka je obrazem přímky $AC$ v~souměrnosti podle $I$, a~proto prochází obrazem bodu $A$, čili bodem $Z$. Stejně přímka vedená bodem $Y$ rovnoběžně s~$AB$ je obrazem přímky $AB$ a~také prochází bodem $Z$. Bod $Z$ je tedy průsečíkem obou tečen ze zadání.

    Stejnolehlost se středem $A$ a~koeficientem $2$ zobrazí $K$ na $C$, $L$ na $B$ a~$I$ na $Z$, přičemž zachovává velikosti úhlů. Proto $|\angle KIL| = |\angle CZB|$ a~zbývá dokázat
    $$
    |\angle BZC| + |\angle YPX| = 180^\circ.
    $$
    Nepříjemné středy $K$, $L$ se tím z~úlohy ztratily.

    Označme $a = |BC|$, $b = |CA|$, $c = |AB|$, $\alpha = |\angle BAC|$, dále $s$ polovinu obvodu, $S$ obsah trojúhelníku $ABC$, $r$ poloměr $\omega$ a~$v_a$, $v_b$, $v_c$ jeho výšky. Ze vztahů $S = rs$ a~$2S = a v_a = b v_b = c v_c$ dostaneme
    $$
    \frac{2r}{v_a} = \frac{a}{s}, \qquad \frac{2r}{v_b} = \frac{b}{s}, \qquad \frac{2r}{v_c} = \frac{c}{s},
    $$
    přičemž všechny tři podíly jsou menší než $1$, neboť trojúhelníkové nerovnosti dávají $a, b, c < s$.

    Měřme vzdálenosti od přímky $BC$ se znaménkem, kladně na straně bodu $A$. Bod $I$ má vzdálenost $r$ a~je středem úsečky $AZ$, takže vzdálenost bodu $Z$ je $2r - v_a < 0$. Bod $Z$ tedy leží v~opačné polorovině než $A$ a~úsečka $AZ$ protíná přímku $BC$; ten průsečík označme $T$. Leží na ose úhlu $\angle BAC$, čili uvnitř úsečky $BC$.

    Druhá tečna $\omega$ rovnoběžná s~$AC$ má od přímky $AC$ vzdálenost $2r$, a~to na té straně, kde leží $I$, čili kde leží i~$B$. Z~nerovnosti $2r < v_b$ proto odděluje $B$ od přímky $AC$, speciálně od bodu $C$, a~tak protíná vnitřek úsečky $BC$: bod $X$ leží uvnitř $BC$. Stejně z~$2r < v_c$ leží uvnitř $BC$ i~bod $Y$.

    Nechť $h$ je stejnolehlost se středem $T$, která zobrazuje $A$ na $Z$. Její koeficient je záporný, neboť $T$ leží uvnitř úsečky $AZ$. Přímku $BC$ zobrazuje $h$ na sebe a~přímku $AC$ na rovnoběžku s~$AC$ procházející bodem $Z$, což je přímka $ZX$; proto $h(C) = X$ a~stejně $h(B) = Y$. Ze záporného koeficientu leží $X$ na polopřímce opačné k~$TC$ a~$Y$ na polopřímce opačné k~$TB$, takže na přímce $BC$ máme pořadí $B$, $X$, $T$, $Y$, $C$. Označíme-li $k$ koeficient stejnolehlosti $h$ v~absolutní hodnotě, pak
    $$
    |TZ| = k \, |TA|, \qquad |TX| = k \, |TC|, \qquad |TY| = k \, |TB|.
    $$

    Přímka $AI$ je osou úhlu $\angle BAC$, takže $P$ je středem oblouku $BC$ neobsahujícího $A$; leží proto v~opačné polorovině než $A$ a~na přímce $BC$ neleží. Ze zadání je $|AB| \neq |AC|$, takže podle tvrzení o~Švrčkově bodu je $P$ středem kružnice opsané čtyřúhelníku $BICI_a$, kde $I_a$ je střed kružnice připsané proti vrcholu $A$; speciálně $|PI| = |PB|$.

    Nechť $F$ je bod dotyku $\omega$ se stranou $AB$ a~nechť $M$ je střed strany $BC$. Úhel při $F$ v~trojúhelníku $AFI$ je pravý a~$|AF| = s - a$. Bod $P$ je středem oblouku $BC$, takže $|PB| = |PC|$ a~přímka $PM$ je osou úsečky $BC$; úhel při $M$ v~trojúhelníku $BMP$ je proto také pravý a~$|BM| = a/2$. Obvodové úhly nad tětivou $PC$ nakonec dávají $|\angle PBC| = |\angle PAC| = \alpha/2 = |\angle FAI|$. Pravoúhlé trojúhelníky $AFI$ a~$BMP$ se tak shodují v~ostrém úhlu, čili jsou podobné (s~$F \leftrightarrow M$, $I \leftrightarrow P$), a~porovnáním odvěsny s~přeponou dostaneme
    $$
    \frac{|AI|}{|AF|} = \frac{|BP|}{|BM|}, \qquad \text{tedy} \qquad |AI| = \frac{2(s-a)}{a}\,|PB|.
    $$
    Jelikož $|PI| = |PB|$, je proto
    $$
    \frac{|AI|}{|PI|} = \frac{2(s-a)}{a} = \frac{b+c-a}{a} < 1,
    $$
    neboť ze zadání je $b < a$ i~$c < a$; právě tady se využívá, že $BC$ je nejdelší strana.

    Bod $I$ leží uvnitř kružnice opsané, tedy uvnitř tětivy $AP$. Bod $Z$ leží na polopřímce $AI$ a~$|AZ| = 2|AI| < |AI| + |IP| = |AP|$, takže $Z$ leží uvnitř úsečky $IP$. Na přímce $AP$ tak máme pořadí $A$, $T$, $Z$, $P$.

    Tětivy $BC$ a~$AP$ kružnice opsané se protínají v~bodě $T$, takže z~mocnosti bodu $T$ je $|TB| \cdot |TC| = |TA| \cdot |TP|$; tuto společnou hodnotu označme $m$. Odtud
    $$
    \gather
    |TB| \cdot |TX| = k \, m = |TZ| \cdot |TP|, \\
    |TC| \cdot |TY| = k \, m = |TZ| \cdot |TP|.
    \endgather
    $$
    Bod $T$ leží mimo úsečku $BX$ (pořadí $B$, $X$, $T$), mimo úsečku $CY$ (pořadí $T$, $Y$, $C$) i~mimo úsečku $ZP$ (pořadí $T$, $Z$, $P$), takže podle kritéria tětivovosti leží na jedné kružnici body $B$, $X$, $Z$, $P$ a~na jedné kružnici i~body $C$, $Y$, $Z$, $P$.

    Body $Z$ a~$P$ leží v~téže polorovině určené přímkou $BC$, tedy na témž oblouku první kružnice určeném tětivou $BX$. Obvodové úhly nad touto tětivou se proto rovnají, čili $|\angle BZX| = |\angle BPX|$, a~z~druhé kružnice stejně $|\angle CZY| = |\angle CPY|$.

    Body $B$, $X$, $Y$, $C$ leží na přímce $BC$ v~tomto pořadí a~body $Z$, $P$ na ní neleží, takže polopřímky z~$Z$ i~polopřímky z~$P$ k~nim jdou v~témž pořadí a~platí
    $$
    \gather
    |\angle BZC| = |\angle BZX| + |\angle XZY| + |\angle YZC|, \\
    |\angle BPC| = |\angle BPX| + |\angle XPY| + |\angle YPC|.
    \endgather
    $$

    Stejnolehlost $h$ zobrazuje trojúhelník $ABC$ na trojúhelník $ZYX$, takže $|\angle XZY| = \alpha$. Body $A$ a~$P$ leží na kružnici opsané po opačných stranách tětivy $BC$, takže $|\angle BPC| = 180^\circ - \alpha$. Po dosazení dvou dokázaných rovností úhlů dostáváme
    $$
    \gather
    |\angle BZC| = |\angle BPX| + \alpha + |\angle YPC|, \\
    |\angle YPX| = 180^\circ - \alpha - |\angle BPX| - |\angle YPC|,
    \endgather
    $$
    a~součet těchto dvou úhlů je $180^\circ$. Jelikož $|\angle KIL| = |\angle BZC|$, je tvrzení dokázáno.
}

\EnvId{h_nSioAPSd9vwFsAInRyp-stredy-oblouku-dvou-kruznic}
\Problem{0}{}{
    Nechť $ABC$ je ostroúhlý trojúhelník se středem~$M$ strany~$BC$. Nechť~$P \neq M$ je bod na straně~$BC$. Označme~$Q$ střed oblouku~$AC$ neobsahujícího~$P$ kružnice opsané trojúhelníku~$APC$. Podobně~$R$ je střed oblouku~$AB$ neobsahujícího~$P$ kružnice opsané trojúhelníku~$APB$. Dokažte, že $Q$, $M$, $P$, $R$ leží na jedné kružnici.
}{
    Nejproblematičtější bod je~$M$: jako středu mu chybí pěkné úhly kolem něj. Je třeba mu dát lepší význam. Dokazované tvrzení také napovídá, co chceme udělat.
}{
    Klíčem je zrcadlit řekněme~$R$ podle~$M$. Pak se vyplatí obarvit úsečky stejné délky; cílem je obarvit další dvojici. Nejsme daleko od kýženého konce.
}{
    Aby kružnice ze zadání existovaly, musí být $P$ vnitřním bodem strany $BC$. Úhly trojúhelníku $ABC$ označme jako obvykle $\alpha$, $\beta$, $\gamma$; trojúhelník je ostroúhlý, takže všechny jsou menší než $90^\circ$. Položme $2\varphi = |\angle APB|$, čili $|\angle APC| = 180^\circ - 2\varphi$ a~$0^\circ < \varphi < 90^\circ$.

    Nejprve si ujasněme polohu bodů $Q$ a~$R$. Bod $R$ leží na oblouku $AB$ neobsahujícím $P$, takže na kružnici $(APB)$ leží body v cyklickém pořadí $A$, $P$, $B$, $R$. Tětiva $PB$ leží na přímce $BC$ a body $A$, $R$ leží na tomtéž oblouku, který tato tětiva vytíná, proto $A$ a~$R$ leží v téže polorovině určené přímkou $BC$. Tětiva $AB$ naopak odděluje $P$ od $R$, takže $R$ a~$C$ leží v opačných polorovinách určených přímkou $AB$ (bod $C$ je u přímky $AB$ na téže straně jako $P$). Stejná úvaha pro kružnici $(APC)$ s cyklickým pořadím $A$, $P$, $C$, $Q$ dává, že $Q$ leží v téže polorovině určené přímkou $BC$ jako $A$ a že $Q$, $B$ leží v opačných polorovinách určených přímkou $AC$. Speciálně $Q$ ani $R$ neleží na přímce $BC$.

    Čtyřúhelník $APBR$ je konvexní, takže polopřímka $PR$ leží uvnitř úhlu $APB$; z rovnosti oblouků $AR$ a~$RB$ je $|\angle APR| = |\angle RPB| = \varphi$ a~$|RA| = |RB|$. Analogicky polopřímka $PQ$ leží uvnitř úhlu $APC$, $|\angle QPC| = 90^\circ - \varphi$ a~$|QA| = |QC|$. Polopřímky $PB$ a~$PC$ jsou opačné a polopřímka $PA$ dělí polorovinu s hraniční přímkou $BC$ obsahující $A$ na úhly $APB$ a~$APC$, uvnitř nichž po řadě leží polopřímky $PR$ a~$PQ$. Proto
    $$
    |\angle RPQ| = 180^\circ - \varphi - (90^\circ - \varphi) = 90^\circ.
    $$
    Bod $P$ tak leží na kružnici s průměrem $QR$ a zbývá dokázat, že na ní leží i $M$, čili že $|\angle QMR| = 90^\circ$.

    Bod $M$ je zatím jen středem úsečky a kolem středu úsečky nejsou žádné úhly, které by šlo počítat. Dáme mu lepší význam tím, co si střed úsečky vždy žádá: nechť $R_0$ je obraz bodu $R$ ve středové souměrnosti se středem $M$, takže $BRCR_0$ je rovnoběžník. Jelikož $R$ neleží na přímce $BC$, je $R_0 \neq R$, a jelikož na ní neleží ani $Q$, je $Q \neq M$. Bod $M$ je středem úsečky $RR_0$, takže dokazovaná rovnost $|\angle QMR| = 90^\circ$ říká přesně to, že $Q$ leží na ose úsečky $RR_0$, čili
    $$
    |QR| = |QR_0|.
    $$

    Teď se vyplatí vyznačit úsečky stejné délky. Jednu dvojici dává $Q$, totiž $|QA| = |QC|$. Druhou dává $R$: středová souměrnost podle $M$ zobrazuje $B$ na $C$ a~$R$ na $R_0$, takže $|CR_0| = |BR| = |AR|$. Trojúhelníky $QAR$ a~$QCR_0$ tak mají shodné dvě dvojice stran a třetí dvojici $QR$, $QR_0$ dostaneme podle věty sus, jakmile ověříme
    $$
    |\angle RAQ| = |\angle R_0CQ|.
    $$

    Přitom budeme opakovaně používat následující pozorování. Leží-li body $X$, $Z$ v opačných polorovinách určených přímkou $VY$, leží uvnitř úhlu $XVZ$ právě jedna ze dvou polopřímek této přímky s počátkem $V$. V prvním případě je $|\angle XVZ| = |\angle XVY| + |\angle YVZ|$, ve druhém $|\angle XVZ| = 360^\circ - |\angle XVY| - |\angle YVZ|$, a jelikož $|\angle XVZ| \leq 180^\circ$, o tom, který případ nastal, rozhoduje právě součet $|\angle XVY| + |\angle YVZ|$.

    Obvodové úhly nad tětivou $RB$ z téhož oblouku dávají $|\angle RAB| = |\angle RPB| = \varphi$, analogicky $|\angle QAC| = |\angle QPC| = 90^\circ - \varphi$. Body $R$, $C$ leží v opačných polorovinách určených přímkou $AB$ a~$\varphi + \alpha < 180^\circ$, takže
    $$
    |\angle RAC| = |\angle RAB| + |\angle BAC| = \varphi + \alpha
    $$
    a polopřímka $AB$ leží uvnitř úhlu $RAC$. Body $R$, $B$ proto leží v téže polorovině určené přímkou $AC$, a jelikož $Q$, $B$ leží v opačných, odděluje přímka $AC$ body $R$ a~$Q$. Součet $(\varphi + \alpha) + (90^\circ - \varphi) = 90^\circ + \alpha$ je menší než $180^\circ$, čili
    $$
    |\angle RAQ| = |\angle RAC| + |\angle CAQ| = 90^\circ + \alpha.
    $$

    U vrcholu $C$ počítejme stejně. Z~$|QA| = |QC|$ je $|\angle QCA| = |\angle QAC| = 90^\circ - \varphi$, přímka $AC$ odděluje $Q$ od $B$ a~$(90^\circ - \varphi) + \gamma < 180^\circ$, takže $|\angle QCB| = 90^\circ - \varphi + \gamma$. Z~$|RA| = |RB|$ je $|\angle RBA| = |\angle RAB| = \varphi$, přímka $AB$ odděluje $R$ od $C$ a~$\varphi + \beta < 180^\circ$, takže $|\angle RBC| = \varphi + \beta$. Středová souměrnost podle $M$ zobrazuje polopřímku $BC$ na polopřímku $CB$ a polopřímku $BR$ na polopřímku $CR_0$, proto $|\angle R_0CB| = |\angle RBC| = \varphi + \beta$.

    Táž souměrnost zobrazuje polorovinu určenou přímkou $BC$ a obsahující $A$ na tu druhou, takže $R_0$ a~$Q$ leží v opačných polorovinách určených přímkou $BC$. Tentokrát je
    $$
    |\angle QCB| + |\angle BCR_0| = 90^\circ + \beta + \gamma = 270^\circ - \alpha,
    $$
    což je díky $\alpha < 90^\circ$ více než $180^\circ$. Uvnitř úhlu $QCR_0$ tedy leží polopřímka opačná k polopřímce $CB$ a~
    $$
    |\angle QCR_0| = 360^\circ - (270^\circ - \alpha) = 90^\circ + \alpha.
    $$

    Trojúhelníky $QAR$ a~$QCR_0$ jsou proto shodné: $|QA| = |QC|$, $|AR| = |CR_0|$ a úhly mezi těmito stranami jsou oba $90^\circ + \alpha$. Odtud $|QR| = |QR_0|$, čili $|\angle QMR| = 90^\circ$ a také bod $M$ leží na kružnici s průměrem $QR$.

    Zbývá ověřit, že jde o čtyři různé body. Z~$|\angle RPQ| = 90^\circ$ jsou $Q$, $P$, $R$ navzájem různé, bod $M$ je různý od $P$ podle zadání a od $Q$ i~$R$ proto, že ty neleží na přímce $BC$. Body $Q$, $M$, $P$, $R$ tak leží na jedné kružnici, a to na kružnici s průměrem $QR$.
}

\EnvId{aq53bc6NlKQWV07NwXRnF-tecna-a-obraz-bodu-r}
\Problem{0}{IMO 2017, P4}{
    Nechť~$R$ a~$S$ jsou různé body na kružnici~$k$ a nechť~$t$ značí tečnu ke~$k$ v bodě~$R$. Nechť~$R_0$ je obraz~$R$ ve středové souměrnosti se středem~$S$. Bod~$I$ je zvolen na kratším oblouku~$RS$ kružnice~$k$ tak, že kružnice~$k_0$ opsaná trojúhelníku~$ISR_0$ protíná~$t$ ve dvou různých bodech. Nechť~$A$ je společný bod~$k_0$ a~$t$ bližší k~$R$. Přímka~$AI$ protíná~$k$ znovu v~$J$. Dokažte, že $JR_0$ je tečnou ke~$k_0$.
}{
    Bod~$S$ je středem nějaké úsečky a to znemožňuje počítání úhlů. Přesto lze něco počítáním úhlů získat, což může napovědět, jak střed~$S$ použít k definici nového bodu.
}{
    Klíčem je objevit $AR_0 \parallel RJ$. To dává docela jasnou nápovědu: obraz~$J$ ve středové souměrnosti se středem~$S$ padne na~$AR_0$. Odtud to dorazíme, máme vše potřebné pro počítání úhlů.
}{
    V celém řešení počítáme s orientovanými úhly přímek modulo $180^\circ$: symbol $\angle(p, q)$ značí úhel, o který je třeba otočit přímku $p$, aby byla rovnoběžná s přímkou $q$. Konfigurace se totiž mění podle polohy $I$ na oblouku i podle toho, který z průsečíků přímky $t$ a kružnice $k_0$ je $A$, a orientované úhly nás těchto rozborů zbaví. Použijeme dvě jejich standardní vlastnosti: čtyři navzájem různé body $P$, $Q$, $X$, $Y$ leží na jedné kružnici nebo přímce právě tehdy, když $\angle(XP, XQ) = \angle(YP, YQ)$, a pro tečnu $\tau$ kružnice v jejím bodě $P$ a další její body $Q$, $X$ platí $\angle(\tau, PQ) = \angle(XP, XQ)$.

    Nejprve si odbudeme ověření, že všechny body, se kterými pracujeme, jsou navzájem různé. Kružnice $k_0$ existuje, takže $I$ neleží na přímce $SR_0$, tedy ani na přímce $RS$; speciálně $I \neq R$. Přímka $t$ má s~$k$ jediný společný bod $R$, takže na ní neleží $I$ ani $S$; přímka $RS$ je proto od $t$ různá a protínají se jedině v~$R$, takže na $t$ neleží ani $R_0$. Bod $A$ je tak různý od $I$, $S$ i~$R_0$. Dále $R$ neleží na $k_0$, neboť jinak by $k_0$ měla s přímkou $RS$ tři různé společné body $R$, $S$, $R_0$; proto $A \neq R$. Bod $R_0$ neleží na $k$, protože střed tětivy $RR_0$ by ležel uvnitř $k$, kdežto $S$ leží na $k$; proto $J \neq R_0$. Přímka $AI$ protíná $k_0$ právě v bodech $A$ a~$I$, takže na ní neleží $S$ a máme $J \neq S$; kdyby na ní ležel $R$, splynula by s~$t$ (obě by procházely různými body $A$ a~$R$) a bod $I$ by ležel na $t$, takže i $J \neq R$. Konečně podle zadání má přímka $AI$ s~$k$ druhý průsečík, tedy $J \neq I$.

    Bod $S$ je středem úsečky $RR_0$ a s takovou podmínkou úhly nic nezmohou. Přímka $R_0S$ je ale tatáž jako $RS$, a na jedno pozorování to stačí. Body $A$, $I$, $S$, $R_0$ leží na $k_0$, body $R$, $S$, $I$, $J$ leží na $k$ a body $A$, $I$, $J$ leží na jedné přímce, takže
    $$
    \gather
    \angle(R_0A, RS) = \angle(R_0A, R_0S) = \angle(IA, IS) \\
    = \angle(IJ, IS) = \angle(RJ, RS),
    \endgather
    $$
    a tedy
    $$
    AR_0 \parallel RJ.
    $$

    Rovnoběžnost rovnou říká, jak střed $S$ využít: středová souměrnost $\sigma$ se středem $S$ zobrazí $R$ na $R_0$, takže obraz $J' = \sigma(J)$ leží na přímce vedené bodem $R_0$ rovnoběžně s~$RJ$, což je právě přímka $AR_0$. Jde o známou situaci, kdy střed úsečky v zadání volá po středové souměrnosti, která vyrobí rovnoběžník $RJR_0J'$. Body $A$, $R_0$, $J'$ tedy leží na jedné přímce.

    Ukážeme, že $A$, $R$, $S$, $J'$ leží na jedné kružnici. Přímka $J'S$ je tatáž jako $JS$, neboť $S$ je středem úsečky $JJ'$ a~$J \neq S$. Spolu s~$AR_0 \parallel RJ$ a s větou o úsekovém úhlu pro tečnu $t$ a tětivu $RS$ kružnice $k$ dostaneme
    $$
    \angle(J'A, J'S) = \angle(JR, JS) = \angle(t, RS) = \angle(RA, RS).
    $$
    Naše čtyři body jsou navzájem různé, netriviální jsou jen $J' \neq R$ (protože $J \neq R_0$), $J' \neq S$ (protože $J \neq S$) a~$J' \neq A$ (jinak by $S$ jakožto střed úsečky $JJ'$ ležel na přímce $AJ$, tedy na přímce $AI$). Na jedné přímce neleží, protože jinak by $A$ ležel na přímce $RS$, a tedy by byl $A = R$. Leží tedy na jedné kružnici.

    Zbývá poznatek přenést zpět k~$R_0$. Souměrnost $\sigma$ zachovává orientované úhly, zobrazuje $J$ na $J'$, $R_0$ na $R$ a přímku $R_0S$, tedy přímku $RR_0$ procházející jejím středem, samu na sebe. Proto
    $$
    \gather
    \angle(JR_0, R_0S) = \angle(J'R, RS) \\
    = \angle(AJ', AS) = \angle(AR_0, AS),
    \endgather
    $$
    kde druhá rovnost plyne z tětivovosti bodů $A$, $R$, $S$, $J'$ a třetí z kolinearity bodů $A$, $R_0$, $J'$. Body $A$, $S$, $R_0$ jsou tři různé body kružnice $k_0$, takže pravá strana je úhel, který s tětivou $R_0S$ svírá tečna ke~$k_0$ v bodě $R_0$. Přímka procházející bodem $R_0$ je přitom svým orientovaným úhlem s~$R_0S$ určena jednoznačně, a proto je $JR_0$ touto tečnou.
}

\EnvId{ZuNLEoEwr-J_fVOQWtSbR-dva-prumery-a-ortocentrum}
\Problem{0}{IMO 2013, P4}{
    Nechť~$ABC$ je ostroúhlý trojúhelník s ortocentrem~$H$ a nechť~$W$ je bod na straně~$BC$. Označme~$M$ a~$N$ paty výšek z~$B$ a~$C$ v trojúhelníku~$ABC$. Nechť~$k_1$ je kružnice opsaná trojúhelníku~$BWN$ a nechť~$X$ je bod na~$k_1$ takový, že~$XW$ je průměrem~$k_1$. Analogicky nechť~$k_2$ je kružnice opsaná trojúhelníku~$CWM$ a nechť~$Y$ je bod na~$k_2$ takový, že~$YW$ je průměrem~$k_2$. Dokažte, že $X$, $Y$, $H$ leží na jedné přímce.
}{
    Úloha nás skoro vybízí, abychom definovali druhý průsečík kružnic~$k_1$ a~$k_2$. Ten je přirozeně propojí a můžeme také dobře využít jejich průměry.
}{
    Je-li~$T$ druhý průsečík~$k_1$ a~$k_2$, pak nám průměry dají pravé úhly, z nichž~$T$, $X$, $Y$ leží na jedné přímce. Důkaz toho, že~$H$ leží na téže přímce, lze přeformulovat jako kolmost. Nezeliminovali jsme náhodou~$X$ a~$Y$?
}{
    Jelikož je trojúhelník $ABC$ ostroúhlý, paty $N$, $M$ leží uvnitř stran $AB$, $AC$ a ortocentrum $H$ uvnitř trojúhelníku, konkrétně uvnitř úsečky $AA'$, kde $A'$ je pata výšky z~$A$. Aby kružnice $k_1$, $k_2$ existovaly, musí být bod $W$ různý od $B$ i od $C$; jsou to pak dvě různé kružnice, neboť tři různé kolineární body $B$, $W$, $C$ nemohou ležet na jedné kružnici.

    Vzájemné polohy bodů v obrázku závisí na tom, kde na straně $BC$ bod $W$ leží, na podstatě argumentu to však nic nemění. Počítáme proto v orientovaných úhlech modulo $180^\circ$: $\angle(p,q)$ značí úhel, o který je třeba otočit přímku $p$, aby byla rovnoběžná s~$q$.

    Kružnice $k_1$ a~$k_2$ procházejí společným bodem $W$, a to je přesně situace, kdy se vyplatí přidat jejich druhý průsečík; označme ho $T$; jeho existenci dokážeme za chvíli. Díky němu lze využít průměry ze zadání: z průměru $WX$ a Thaletovy věty je $|\angle WTX| = 90^\circ$ (pro $X = T$ triviálně), takže $X$ leží na přímce $\ell$ vedené bodem $T$ kolmo k $TW$, a stejně tak z průměru $WY$ leží na $\ell$ i~$Y$. Body $X$ a~$Y$ jsme tím z úlohy odstranili a zbývá dokázat, že na $\ell$ leží i~$H$, což pro $T \neq H$ znamená
    $$
    TH \perp TW.
    $$
    (Je-li $T = H$, jsme hotovi, neboť $T$ na $\ell$ leží.)

    Klíčovou kružnicí je $\omega$ nad průměrem $AH$. Bod $H$ leží na výšce $CN$ kolmé k $AB$ a na výšce $BM$ kolmé k $AC$, takže $|\angle ANH| = |\angle AMH| = 90^\circ$ a kromě $A$ a~$H$ leží na $\omega$ i paty $N$ a~$M$. Přímku $BC$ přitom $\omega$ neprotíná: body $A$ i~$H$ se promítají do $A'$, takže střed $\omega$, tedy střed úsečky $AH$, má od $BC$ vzdálenost $(|AA'| + |HA'|)/2$, kdežto poloměr $\omega$ je $|AH|/2 = (|AA'| - |HA'|)/2$, což je méně.

    Nyní existence bodu $T$. Kdyby se $k_1$ a~$k_2$ v bodě $W$ dotýkaly, měly by tam společnou tečnu $t$ a věta o úsekovém úhlu by v obou kružnicích dala
    $$
    \gather
    \angle(WN, WM) = \angle(WN, t) + \angle(t, WM) \\
    = \angle(BN, BW) + \angle(CW, CM) \\
    = \angle(AB, BC) + \angle(BC, AC) = \angle(AN, AM).
    \endgather
    $$
    Bod $W$ by tedy ležel na kružnici opsané trojúhelníku $ANM$, čili na $\omega$, která přímku $BC$ neprotíná. To je spor, takže kružnice $k_1$ a~$k_2$ mají kromě $W$ ještě druhý společný bod $T$.

    Týž výpočet, jen s obvodovými úhly namísto úsekového, ukáže, že $T$ leží na $\omega$. Nejprve $T \neq N$: jinak by na $k_2$ ležely body $C$, $M$, $N$, které leží i na kružnici $\gamma$ nad průměrem $BC$ (vždyť $|\angle BNC| = |\angle BMC| = 90^\circ$), takže by bylo $k_2 = \gamma$ a bod $W$ by jako společný bod $\gamma$ a přímky $BC$ splynul s~$B$ nebo $C$. Stejně $T \neq M$. Obvodové úhly v~$k_1$ a v~$k_2$ tak dávají
    $$
    \gather
    \angle(TN, TM) = \angle(TN, TW) + \angle(TW, TM) \\
    = \angle(BN, BW) + \angle(CW, CM) = \angle(AN, AM),
    \endgather
    $$
    čili body $A$, $N$, $M$, $T$ leží na jedné kružnici a~$T$ leží na $\omega$.

    Zbývá kolmost; nechť tedy $T \neq H$. Body $T$, $B$, $W$, $N$ leží na $k_1$ a body $T$, $A$, $N$, $H$ na $\omega$, takže
    $$
    \gather
    \angle(TW, TH) = \angle(TW, TN) + \angle(TN, TH) \\
    = \angle(BW, BN) + \angle(AN, AH).
    \endgather
    $$
    První sčítanec je $\angle(BC, AB)$ a druhý $\angle(AB, AH)$, což je kvůli $AH \perp BC$ přesně $\angle(AB, BC) + 90^\circ$. Součet je tedy $90^\circ$, čili $TH \perp TW$. Bod $H$ proto leží na přímce $\ell$, na níž leží i $X$ a~$Y$.
}

\EnvId{FhO0eO_oi-IndJmQ9N6Ql-dva-body-na-strane-bc}
\Problem{0}{IMO 2014, P4}{
    Na straně~$BC$ ostroúhlého trojúhelníku $ABC$ jsou zvoleny body~$P$ a~$Q$ tak, že $|\angle PAB| = |\angle ACB|$ a $|\angle QAC| = |\angle CBA|$. Body~$M$ a~$N$ leží po řadě na polopřímkách~$AP$ a~$AQ$ tak, že $|AP| = |PM|$ a $|AQ| = |QN|$. Dokažte, že přímky~$BM$ a~$CN$ se protínají na kružnici opsané trojúhelníku~$ABC$.
}{
    Body~$Q$ a~$P$ jsou po řadě středy~$AN$ a~$AM$. Tento fakt je sám o sobě těžké využít. Klíčem je definovat nějaké další středy. V nejlepším případě nám tyto středy pomohou eliminovat nepříliš přirozeně vypadající body~$M$ a~$N$.
}{
    Definujte středy~$J$ a~$K$ stran~$AC$ a~$AB$. Střední příčky nám dají $BM \parallel KP$ a $CN \parallel QJ$. Zbývá tedy ukázat, že úhel mezi úsečkami~$QJ$ a~$KP$ je $180^\circ - \alpha$. Zdá se, že jsme nepokročili, neboť úhly u těžnic jsou obecně ošklivé. V konfiguraci je však mnoho podobných trojúhelníků.
}{
    Označme $\alpha = |\angle BAC|$, $\beta = |\angle CBA|$, $\gamma = |\angle ACB|$. Rovnosti $|AP| = |PM|$ a $|AQ| = |QN|$ říkají, že $P$ je střed úsečky $AM$ a~$Q$ je střed úsečky $AN$.

    Nejprve převeďme úhlové podmínky na délky. Bod $P$ leží na straně $BC$, takže $|\angle ABP| = \beta = |\angle CBA|$, a~ze zadání $|\angle BAP| = \gamma = |\angle BCA|$; podle věty $uu$ je $\triangle BAP \sim \triangle BCA$ (v~pořadí vrcholů $B \leftrightarrow B$, $A \leftrightarrow C$, $P \leftrightarrow A$). Analogicky $\triangle CAQ \sim \triangle CBA$, a~z~poměrů stran
    $$
    \frac{|BP|}{|BA|} = \frac{|BA|}{|BC|}, \qquad \frac{|CQ|}{|CA|} = \frac{|CA|}{|CB|}. \eqno(*)
    $$

    Body $M$ a~$N$ vstupují do zadání jedině přes středy $P$ a~$Q$ a samy o sobě se uchopit nedají. Jsou-li v úloze středy úseček, vyplatí se doplnit další středy, aby vznikly střední příčky: nechť $K$, $J$, $L$ jsou po řadě středy stran $AB$, $AC$, $BC$. V~trojúhelníku $ABM$ jsou $K$ a~$P$ středy stran $AB$ a~$AM$, takže $KP \parallel BM$; v~trojúhelníku $ACN$ jsou $J$ a~$Q$ středy stran $AC$ a~$AN$, takže $JQ \parallel CN$. Body $M$, $N$ tím z~obrázku zmizely: úhel přímek $BM$ a~$CN$ je týž jako úhel přímek $KP$ a~$JQ$ a stačí ukázat, že je roven $180^\circ - \alpha$.

    Obě přímky protínají $BC$, změřme tedy jejich úhly vůči ní. Jelikož $|BK| = \tfrac12|AB|$ a $|BL| = \tfrac12|BC|$, dává $(*)$
    $$
    \frac{|BK|}{|BL|} = \frac{|AB|}{|BC|} = \frac{|BP|}{|BA|},
    $$
    přičemž $\angle KBP$ a~$\angle LBA$ je týž úhel $\beta$ (bod $K$ leží na úsečce $AB$ a~body $P$, $L$ na úsečce $BC$). Podle věty $sus$ je $\triangle BKP \sim \triangle BLA$ (v~pořadí vrcholů $B \leftrightarrow B$, $K \leftrightarrow L$, $P \leftrightarrow A$), a~tedy $|\angle BPK| = |\angle BAL|$. Stejně z~$|CJ| : |CL| = |AC| : |BC| = |CQ| : |CA|$ a~ze společného úhlu při $C$ je $\triangle CJQ \sim \triangle CLA$ a $|\angle CQJ| = |\angle CAL|$. Úhly u těžnice $AL$ jsou sice ošklivé, ale potřebujeme jen jejich součet: bod $L$ leží uvnitř strany $BC$, takže polopřímka $AL$ leží uvnitř úhlu $BAC$ a
    $$
    |\angle BPK| + |\angle CQJ| = |\angle BAL| + |\angle LAC| = \alpha.
    $$

    Zbývá přenést tyto dva úhly k~vrcholům $B$ a~$C$. Bod $M$ je obrazem bodu $A$ ve středové souměrnosti podle bodu $P$ ležícího na $BC$, leží tedy v~polorovině $\rho$ s~hraniční přímkou $BC$ neobsahující $A$; totéž platí pro $N$. Naopak $K$ a~$J$ leží uvnitř úseček $AB$ a~$AC$, tedy mimo $\rho$. Přímka $BC$ je příčkou rovnoběžek $KP$, $BM$, body $K$, $M$ leží u ní na opačných stranách a~polopřímky $PB$, $BC$ míří opačně (bod $P$ leží na $BC$ a~$P \neq B$), takže jde o~střídavé úhly; stejně pro rovnoběžky $JQ$, $CN$. Dostáváme
    $$
    |\angle MBC| = |\angle BPK|, \qquad |\angle NCB| = |\angle CQJ|.
    $$

    Polopřímky $BM$ a~$CN$ tak vycházejí z~konců úsečky $BC$ do téže poloroviny $\rho$ a~svírají s~$BC$ úhly, jejichž součet je $\alpha < 180^\circ$. Protnou se proto v~bodě $X$ ležícím v~$\rho$ a~ze součtu úhlů v~trojúhelníku $BXC$ je
    $$
    |\angle BXC| = 180^\circ - |\angle MBC| - |\angle NCB| = 180^\circ - \alpha.
    $$
    Body $A$ a~$X$ leží v~opačných polorovinách určených přímkou $BC$ a $|\angle BXC| + |\angle BAC| = 180^\circ$, takže $ABXC$ je tětivový čtyřúhelník. Průsečík $X$ přímek $BM$ a~$CN$ proto leží na kružnici opsané trojúhelníku $ABC$.
}

\EnvId{3AK1oHdV1pF38vTXfQc41-stredy-kruznic-u-pat-vysek}
\Problem{0}{ISL 2012, G3}{
    Nechť $ABC$ je ostroúhlý trojúhelník. Označme $D$, $E$, $F$ paty výšek po řadě z~$A$, $B$, $C$. Nechť~$I_1$ a~$I_2$ jsou po řadě středy vepsaných kružnic trojúhelníků~$BFD$ a~$CED$ a nechť~$O_1$ a~$O_2$ jsou po řadě středy kružnic opsaných trojúhelníkům~$AI_1B$ a~$AI_2C$. Dokažte, že $I_1I_2 \parallel O_1O_2$.
}{
    Body~$O_1$ a~$O_2$ vypadají opravdu uměle. A skutečně jsou. Oba lze eliminovat definováním nového bodu a přeformulováním rovnoběžnosti do něčeho jiného. I bez těchto podivných bodů v konfiguraci platí víc, než úloha uvádí, a není to úplně snadné vidět. Vyplatí se to uhodnout.
}{
    Nechť~$T$ je druhý průsečík našich dvou kružnic se středy~$O_1$ a~$O_2$. Chceme ukázat, že~$AT$ je kolmé na~$I_1I_2$. To stále není zadarmo. Pomůže uhodnout a dokázat, že $BCI_2I_1$ je tětivový čtyřúhelník.
}{
    Označme $\alpha$, $\beta$, $\gamma$ vnitřní úhly trojúhelníku $ABC$, dále $I$ střed jeho vepsané kružnice a~$r$ její poloměr. Kružnice opsané trojúhelníkům $AI_1B$ a~$AI_2C$, tedy kružnice se středy $O_1$ a~$O_2$, označme $\omega_1$ a~$\omega_2$.

    Body $O_1$, $O_2$ vstupují do tvrzení jedině skrze přímku $O_1O_2$, a ta je osou společné tětivy $AT$, kde $T$ je druhý společný bod kružnic $\omega_1$, $\omega_2$. Stačí tedy dokázat, že $AT \perp I_1I_2$; existenci bodu $T$ dokážeme za chvíli a ukáže se přitom, že přímka $AT$ není nic jiného než osa úhlu $\angle BAC$. Z přesného obrázku lze uhodnout, že v konfiguraci platí víc, než úloha tvrdí: $BCI_2I_1$ je tětivový čtyřúhelník. To je jádro řešení a začneme jím.

    Trojúhelník $ABC$ je ostroúhlý, takže $D$, $E$, $F$ leží po řadě uvnitř stran $BC$, $CA$, $AB$. Úhly $\angle AFC$ a~$\angle ADC$ jsou pravé, takže $A$, $F$, $D$, $C$ leží na kružnici s~průměrem $AC$ a mocnost bodu $B$ vzhledem k ní dává $|BF| \cdot |BA| = |BD| \cdot |BC|$, neboli
    $$
    \frac{|BF|}{|BC|} = \frac{|BD|}{|BA|}.
    $$
    Tento společný poměr označme $k$; jelikož $BD$ je odvěsna a~$BA$ přepona pravoúhlého trojúhelníku $ABD$, je $0 < k < 1$. Úhel $\angle FBD$ je přitom týž jako $\angle CBA$, takže trojúhelníky $BFD$ a~$BCA$ se shodují v úhlu při $B$ i v poměru jeho ramen, tedy jsou podobné s~koeficientem $k$, přičemž $F \leftrightarrow C$ a~$D \leftrightarrow A$.

    Bod $I_1$ je střed vepsané kružnice trojúhelníku $BFD$, takže přímka $BI_1$ je osou úhlu $\angle FBD$, tedy osou úhlu $\angle ABC$, a~$I_1$ leží na polopřímce $BI$. Podobnost zobrazuje střed vepsané kružnice na střed vepsané kružnice, takže $|BI_1| = k|BI|$, a z~$0 < k < 1$ plyne, že $I_1$ leží uvnitř úsečky $BI$.

    Součin $|IB| \cdot |II_1|$ změříme pomocí obrazu středu $I$ v osové souměrnosti. Nechť $X$ je bod dotyku vepsané kružnice se stranou $BC$ a nechť $I'$ je obraz bodu $I$ v souměrnosti podle přímky $BC$. Bod $X$ je patou kolmice z~$I$ na $BC$, tedy středem úsečky $II'$; proto $|IX| = r$, $|II'| = 2r$ a polopřímky $IX$, $II'$ splývají. Dále $|BI'| = |BI|$ a $|\angle IBI'| = 2|\angle IBC| = \beta$.

    Nechť $Q$ je pata kolmice z~$I'$ na přímku $BI$. Úhel $\angle IBI'$ je nenulový a ostrý, takže $I'$ neleží na přímce $BI$ a $Q$ leží na polopřímce $BI$. Pravoúhlé trojúhelníky $BQI'$ a~$BDA$ mají pravé úhly při $Q$ a~$D$ a shodují se v úhlu při $B$, neboť $|\angle QBI'| = \beta = |\angle DBA|$. Jsou tedy podobné a
    $$
    |BQ| = |BI'| \cdot \frac{|BD|}{|BA|} = k|BI| = |BI_1|,
    $$
    a jelikož $Q$ i~$I_1$ leží na polopřímce $BI$, je $Q = I_1$. Zvláště $|\angle II_1I'| = 90^\circ$.

    Bod $I_1$ leží uvnitř úsečky $BI$, takže polopřímky $II_1$ a~$IB$ splývají, a jak jsme viděli, splývají i polopřímky $II'$ a~$IX$. Pravoúhlé trojúhelníky $II_1I'$ a~$IXB$ (s~pravými úhly při $I_1$ a~$X$) se tak shodují v úhlu při $I$, jsou tedy podobné, přičemž $I_1 \leftrightarrow X$ a~$I' \leftrightarrow B$. Odtud
    $$
    |IB| \cdot |II_1| = |IX| \cdot |II'| = 2r^2.
    $$

    U vrcholu $C$ je situace stejná. Úhly $\angle AEB$ a~$\angle ADB$ jsou pravé, takže $A$, $E$, $D$, $B$ leží na kružnici s~průměrem $AB$ a mocnost bodu $C$ dává $|CE| \cdot |CA| = |CD| \cdot |CB|$, neboli
    $$
    \frac{|CE|}{|CB|} = \frac{|CD|}{|CA|} = k',
    $$
    kde $0 < k' < 1$, neboť $CD$ je odvěsna a~$CA$ přepona pravoúhlého trojúhelníku $ACD$. Trojúhelníky $CED$ a~$CBA$ jsou proto podobné s~koeficientem $k'$ (s~$E \leftrightarrow B$ a~$D \leftrightarrow A$), neboť úhel $\angle ECD$ je týž jako $\angle BCA$. Přímka $CI_2$ je proto osou úhlu $\angle ACB$, bod $I_2$ leží uvnitř úsečky $CI$ a $|CI_2| = k'|CI|$. Týž bod $I'$ splňuje $|CI'| = |CI|$ a $|\angle ICI'| = 2|\angle ICB| = \gamma$; pata kolmice z~$I'$ na přímku $CI$ tak leží na polopřímce $CI$ a z~podobnosti s~pravoúhlým trojúhelníkem $CDA$ je její vzdálenost od $C$ rovna $k'|CI|$, čili je to bod $I_2$. Táž dvojice podobných pravoúhlých trojúhelníků jako u $B$ pak dává $|IC| \cdot |II_2| = |IX| \cdot |II'| = 2r^2$.

    Přímky $BI$ a~$CI$ jsou různé a protínají se jedině v $I$, takže body $B$, $I_1$, $I_2$, $C$ (žádný z nich není $I$) jsou navzájem různé a neleží na jedné přímce. Bod $I$ leží mimo úsečku $BI_1$ i mimo úsečku $CI_2$, takže z~rovnosti $|IB| \cdot |II_1| = |IC| \cdot |II_2|$ podle kritéria tětivovosti leží body $B$, $I_1$, $I_2$, $C$ na jedné kružnici.

    Táž rovnost nám dá i bod $T$. Body $B$, $I_1$ leží na $\omega_1$ a bod $I$ leží na přímce $BI_1$ mimo úsečku $BI_1$, takže $I$ leží vně $\omega_1$ a jeho mocnost vzhledem k~$\omega_1$ je $|IB| \cdot |II_1| = 2r^2$; stejně mocnost $I$ vzhledem k~$\omega_2$ je $|IC| \cdot |II_2| = 2r^2$. Bod $I$ tedy leží na chordále kružnic $\omega_1$, $\omega_2$, na níž leží i jejich společný bod $A$.

    Nechť $Z$ je bod dotyku vepsané kružnice se stranou $AB$. Trojúhelník $AZI$ má pravý úhel při $Z$, úhel $\frac\alpha2$ při $A$ a přeponu $AI$. Z~$\alpha < 90^\circ$ je úhel při $A$ menší než $45^\circ$, takže úhel při $I$ je větší než $45^\circ$; proti většímu úhlu leží větší strana, čili $|AZ| > |ZI| = r$ a podle Pythagorovy věty
    $$
    |IA|^2 = r^2 + |AZ|^2 > 2r^2.
    $$
    Kdyby se přímka $AI$ dotýkala $\omega_1$, dotýkala by se jí v bodě $A$ a mocnost $I$ by byla $|IA|^2 \neq 2r^2$. Přímka $AI$ tedy protíná $\omega_1$ ještě v bodě $T_1 \neq A$, přičemž $I$ leží vně $\omega_1$, takže $A$ a~$T_1$ leží na téže polopřímce z~$I$ a $|IT_1| = 2r^2/|IA|$. Totéž s~$\omega_2$ dává bod $T_2$ na téže polopřímce s~$|IT_2| = 2r^2/|IA|$, čili $T_1 = T_2$. Kružnice $\omega_1$, $\omega_2$ tak mají kromě $A$ ještě druhý společný bod $T = T_1$ a ten leží na přímce $AI$.

    Bod $I_1$ leží uvnitř úsečky $BI$, tedy uvnitř trojúhelníku $ABC$. Kdyby $\omega_1 = \omega_2$, procházela by tato kružnice body $A$, $B$, $C$, byla by to kružnice opsaná trojúhelníku $ABC$ a bod $I_1$ by na ní ležel, ne uvnitř. Kružnice $\omega_1$, $\omega_2$ jsou tedy různé a mají společný bod, takže nejsou soustředné a $O_1 \neq O_2$. Oba středy jsou stejně vzdálené od $A$ i od $T$, takže přímka $O_1O_2$ je osou úsečky $AT$, a jelikož přímka $AT$ je přímka $AI$, dostáváme $O_1O_2 \perp AI$.

    Zbývá dokázat, že $I_1I_2 \perp AI$. Nechť $I_a$ je střed kružnice připsané proti vrcholu $A$. Leží na ose úhlu $\angle BAC$ na opačné straně přímky $BC$ než $A$, takže $I$ leží uvnitř úsečky $AI_a$. Přímky $BI$, $BI_a$ jsou vnitřní a vnější osa úhlu při $B$, takže $|\angle IBI_a| = 90^\circ$, a stejně $|\angle ICI_a| = 90^\circ$. Podle věty o základních úhlech (po přejmenování vrcholů) je
    $$
    |\angle AIB| = 90^\circ + \tfrac\gamma2, \qquad |\angle AIC| = 90^\circ + \tfrac\beta2.
    $$

    Nechť $J$ je pata kolmice z~$I_1$ na přímku $AI$. Bod $I_1$ leží na přímce $BI$ a je různý od $I$, takže na přímce $AI$ neleží a úhly trojúhelníku $IJI_1$ při $I$ a~$I_1$ jsou ostré (zvláště $J \neq I$). Jelikož $I_1$ leží na polopřímce $IB$, je $|\angle AII_1| = |\angle AIB| > 90^\circ$, takže $J$ neleží na polopřímce $IA$, nýbrž na polopřímce k ní opačné, čili na polopřímce $II_a$. Úhly $\angle JII_1$ a~$\angle I_aIB$ jsou tak týž úhel a pravoúhlé trojúhelníky $IJI_1$, $IBI_a$ jsou podobné (s~$J \leftrightarrow B$, $I_1 \leftrightarrow I_a$), odkud
    $$
    |IJ| \cdot |II_a| = |II_1| \cdot |IB| = 2r^2.
    $$
    Stejně pata $J'$ kolmice z~$I_2$ na $AI$ leží na polopřímce $II_a$, neboť $|\angle AII_2| = |\angle AIC| > 90^\circ$, a z~podobnosti trojúhelníků $IJ'I_2$, $ICI_a$ dostaneme
    $$
    |IJ'| \cdot |II_a| = |II_2| \cdot |IC| = 2r^2.
    $$

    Body $J$, $J'$ tedy leží na téže polopřímce z~$I$ a $|IJ| = |IJ'|$, čili $J = J'$. Přímky $I_1J$ a~$I_2J$ jsou pak obě kolmicí na $AI$ v bodě $J$, a proto splývají. Body $I_1$, $J$, $I_2$ tak leží na jedné přímce a $I_1I_2 \perp AI$. Dohromady s~$O_1O_2 \perp AI$ to dává $I_1I_2 \parallel O_1O_2$.
}

\EnvId{a0vidDqjikMyjuIYb5vVc-ctyruhelnik-z-os-uhlu}
\Problem{0}{MEMO 2012, T6}{
    Nechť~$ABCD$ je konvexní čtyřúhelník, jehož žádné dvě strany nejsou rovnoběžné a $|\angle ABC| = |\angle CDA|$. Předpokládejme, že po dvou vzaté průsečíky os úhlů u sousedních vrcholů~$ABCD$ tvoří čtyřúhelník~$EFGH$. Nechť~$K$ je průsečík úhlopříček~$EFGH$. Dokažte, že průsečík přímek~$AB$ a~$CD$ leží na kružnici opsané trojúhelníku~$BKD$.
}{
    Úloha se rýsuje trochu neohrabaně. Nenechte se zastrašit. Nestandardní úhlová podmínka vypadá, že by mohla něco pěkného znamenat. Znamená: stačí jen něco protnout. Co se týče os úhlů, nebojte se jich; dvě často implikují třetí a tady jich je spousta.
}{
    Nechť~$E$ leží na osách u~$A$ a~$B$, $F$ na osách u~$B$ a~$C$, $G$ na osách u~$C$ a~$D$ a~$H$ na osách u~$D$ a~$A$. Všechny body~$E$, $F$, $G$, $H$ jsou díky osám úhlů středy vepsaných nebo připsaných kružnic nějakých trojúhelníků. Každopádně body~$P$, $F$, $H$ a $Q$, $E$, $G$ leží na jedné přímce a tyto přímky jsou osy úhlů~$\angle CPB$ a~$\angle CQD$. Teď už jsme velmi blízko.
}{
    Stačí protnout vhodné přímky. Kromě průsečíku $P$ přímek $AB$ a~$CD$ ze zadání přidáme průsečík $Q$ přímek $BC$ a~$AD$; oba existují, protože žádné dvě strany nejsou rovnoběžné. Teprve tato dvojice dá podmínce $|\angle ABC| = |\angle CDA|$ tvar, se kterým se dá pracovat.

    Označme $\alpha = |\angle DAB|$, $\beta = |\angle ABC|$, $\gamma = |\angle BCD|$ a~$\delta = |\angle CDA|$, takže $\beta = \delta$ a~$\alpha + 2\beta + \gamma = 360^\circ$.

    Přímka každé strany konvexního čtyřúhelníku je opornou přímkou a zbylé dva vrcholy leží striktně v jedné polorovině. Body $A$, $B$ tedy leží striktně na téže straně přímky $CD$, čili $P$ neleží na úsečce $AB$ a polopřímky $PA$, $PB$ splývají; symetricky splývají $PC$ s~$PD$, $QB$ s~$QC$ a~$QA$ s~$QD$. Úhly $\angle BPC$, $\angle APD$ a~$\angle BPD$ jsou proto týž úhel, stejně jako $\angle AQB$, $\angle CQD$ a~$\angle BQD$.

    Přímka $BC$ je příčkou přímek $AB$ a~$CD$ a body $A$, $D$ leží na téže její straně, takže $AB \parallel CD$ právě tehdy, když $\beta + \gamma = 180^\circ$; to je vyloučeno. Záměna označení $(A,B,C,D) \to (C,D,A,B)$ je cyklickým posunem: zachová předpoklady i body $P$ a~$Q$, zamění $E$ s~$G$ a~$F$ s~$H$ (úhlopříčky $EFGH$, a tedy ani bod $K$, se nezmění) a zamění $B$ s~$D$, přičemž $\beta + \gamma$ přejde na $\alpha + \beta = 360^\circ - \beta - \gamma$. Dokazované tvrzení je vůči ní invariantní, takže smíme předpokládat $\beta + \gamma < 180^\circ$.

    Body $B$, $C$ leží striktně na téže straně přímky $AD$, zatímco $A$ a~$D$ leží na ní. Proto $A$ leží mezi $P$ a~$B$ právě tehdy, když jsou $P$ a~$B$ v opačných polorovinách, což nastane právě tehdy, když jsou v opačných polorovinách $P$ a~$C$, čili právě tehdy, když $D$ leží mezi $P$ a~$C$. Úhly trojúhelníku $PBC$ u vrcholů $B$ a~$C$ jsou tak buď $\beta$ a~$\gamma$, nebo $180^\circ - \beta$ a~$180^\circ - \gamma$; druhá možnost dává součet $360^\circ - \beta - \gamma > 180^\circ$, takže platí první a bod $A$ leží mezi $P$ a~$B$ a bod $D$ mezi $P$ a~$C$. Táž úvaha s přímkou $AB$ dává, že $B$ leží mezi $Q$ a~$C$ právě tehdy, když $A$ leží mezi $Q$ a~$D$, a úhly trojúhelníku $QCD$ u~$C$ a~$D$ jsou $\gamma$ a~$\delta$, neboť $360^\circ - \gamma - \delta = 360^\circ - \gamma - \beta > 180^\circ$. Odtud
    $$
    |\angle BPC| = 180^\circ - \beta - \gamma = 180^\circ - \gamma - \delta = |\angle CQD|.
    $$
    To je celý obsah podmínky $\beta = \delta$: úhly u~$P$ a~$Q$ jsou shodné.

    Teď osy úhlů. V trojúhelníku $PBC$ jsou osy čtyřúhelníku u~$B$ a~$C$ jeho vnitřními osami, takže $F$ je střed jeho vepsané kružnice. V trojúhelníku $PAD$ je $|\angle PAD| = 180^\circ - \alpha$ a~$|\angle PDA| = 180^\circ - \delta$, takže osy čtyřúhelníku u~$A$ a~$D$ jsou jeho vnějšími osami a $H$ je střed kružnice připsané proti vrcholu $P$. Oba body proto leží na vnitřní ose $p$ úhlu $\angle BPC$. Analogicky je $G$ střed vepsané kružnice trojúhelníku $QCD$ a~$E$ střed kružnice připsané trojúhelníku $QAB$ proti vrcholu $Q$, takže oba leží na vnitřní ose $q$ úhlu $\angle CQD$. Úhlopříčkami čtyřúhelníku $EFGH$ jsou tedy přímky $EG = q$ a~$FH = p$ a bod $K$ je průsečíkem přímek $p$ a~$q$.

    Ze zjištěných polopřímek: $BP$ je polopřímka $BA$ a~$BQ$ je opačná k~$BC$, takže $|\angle PBQ| = 180^\circ - \beta$; podobně $DP$ je opačná k~$DC$ a~$DQ$ je polopřímka $DA$, takže $|\angle PDQ| = 180^\circ - \delta$. Tyto dva úhly jsou shodné. Bod $A$ leží striktně mezi $P$ a~$B$ i striktně mezi $Q$ a~$D$, přičemž neleží na přímce $PQ$ (jinak by přímky $AB$ a~$AD$ splynuly). Bod $B$ leží na polopřímce z bodu $P$ přes $A$ za bodem $A$, a jelikož $P$ leží na přímce $PQ$, je $B$ striktně v téže polorovině jako $A$; stejně $D$ leží na polopřímce z bodu $Q$ přes $A$ za bodem $A$, čili také v této polorovině. Body $B$ a~$D$ tak vidí úsečku $PQ$ pod stejným úhlem a z téže strany, takže $P$, $B$, $D$, $Q$ leží na jedné kružnici $\omega$.

    Jelikož $B$ a~$D$ leží v téže polorovině s hraniční přímkou $PQ$, tětivy $BD$ a~$PQ$ se neprotínají, a tedy $P$ i $Q$ leží na témže z oblouků určených body $B$ a~$D$. Přímka $p$ je vnitřní osou úhlu $\angle BPD$, takže její druhý průsečík $M$ s~$\omega$ dělí oblouk $BD$ neobsahující $P$ na dva shodné oblouky, neboť obvodové úhly $\angle BPM$ a~$\angle MPD$ jsou shodné. Bod $Q$ leží na témže oblouku jako $P$, takže ani jeden z oblouků $BM$, $MD$ neobsahuje $Q$; obvodové úhly nad nimi dávají $|\angle BQM| = |\angle MQD|$ a díky poloze $M$ na oblouku $BD$ neobsahujícím $Q$ leží polopřímka $QM$ uvnitř úhlu $\angle BQD$. Přímka $QM$ je tedy vnitřní osou úhlu $\angle BQD$, čili $QM = q$.

    Bod $M$ leží na $p$ i na $q$, proto $M = K$. Body $B$, $K$, $D$ jsou tři různé body kružnice $\omega$, takže $\omega$ je kružnice opsaná trojúhelníku $BKD$, a leží na ní i bod $P$.
}

\EnvId{DzF1H7zPMIba3UVbwJoNr-bod-na-kruznici-apob}
\Problem{0}{MEMO 2016, I3}{
    Nechť~$ABC$ je ostroúhlý trojúhelník se středem opsané kružnice~$O$ a $|\angle BAC| > 45^\circ$. Bod~$P$ leží uvnitř něj tak, že $A$, $P$, $O$, $B$ leží na jedné kružnici a přímka~$BP$ je kolmá na~$CP$. Bod~$Q$ leží na úsečce~$BP$ tak, že~$AQ$ a~$PO$ jsou rovnoběžné. Dokažte, že $|\angle QCB| = |\angle PCO|$.
}{
    Je zřejmé, že v úloze něco chybí. Vyplatí se proto nejprve získat lepší obrázek situace přeformulováním tvrzení jako rovnosti nějakých jiných úhlů, ideálně takových, které lze vyjádřit pomocí základních úhlů~$ABC$. To může vnuknout přidání. Nezapomeňme, že úloha stále obsahuje známý indikátor potenciálně dobrých přidání: pravý úhel.
}{
    Tvrzení je ekvivalentní $|\angle PCQ|=|\angle OCB|$, což je ekvivalentní $|\angle BAC|=|\angle CQP|$. Takže zrcadlíme-li~$Q$ podle~$P$ do~$Q'$, stačí dokázat, že body leží na jedné kružnici. Ani to není zřejmé, ale v tomto bodě to lze dotáhnout.
}{
    Označme $\alpha = |\angle BAC|$ a~$\gamma = |\angle BCA|$, nechť $\omega$ je kružnice procházející body $A$, $P$, $O$, $B$ a~nechť $k$ je kružnice s~průměrem $BC$. Z~$|\angle BPC| = 90^\circ$ leží $P$ na $k$. Trojúhelník $ABC$ je ostroúhlý, takže $O$ leží uvnitř něj a~podle tvrzení o~základních úhlech je $|\angle BOC| = 2\alpha$ a~$|\angle AOB| = 2\gamma$. Z~rovnoramenných trojúhelníků $OBC$ a~$OAB$ tak dostáváme
    $$
    |\angle OCB| = 90^\circ - \alpha, \qquad |\angle OAB| = 90^\circ - \gamma.
    $$

    Předpoklad $\alpha > 45^\circ$ použijeme jen jednou, a~to hned teď: dává $|\angle BOC| = 2\alpha > 90^\circ$, čili $O$ leží uvnitř kružnice $k$, kdežto z~$\alpha < 90^\circ$ leží $A$ mimo $k$. Kružnice $\omega$ prochází vnitřním bodem $O$ kružnice $k$, takže $\omega \neq k$ a~tyto dvě kružnice mají společné právě body $B$ a~$P$ (ty jsou různé, neboť $P$ leží uvnitř trojúhelníku). Dělí $\omega$ na dva oblouky, z~nichž jeden leží uvnitř $k$ a~druhý mimo $k$, takže $A$ a~$O$ leží na různých obloucích: na $\omega$ jdou naše čtyři body v~pořadí $A$, $P$, $O$, $B$. Speciálně $P \neq O$, takže přímka $PO$ je vůbec dobře definovaná; a~jelikož $A$, $P$, $O$ jsou tři různé body kružnice, nejsou kolineární, takže z~$AQ \parallel PO$ máme $Q \neq P$.

    Dokážeme, že $|\angle PCQ| = 90^\circ - \alpha$, což už stačí. Bod $Q$ totiž leží na úsečce $BP$, takže polopřímka $CQ$ leží uvnitř úhlu $BCP$ a
    $$
    |\angle BCP| = |\angle BCQ| + |\angle QCP| \geq |\angle QCP|.
    $$
    Je-li tedy $|\angle QCP| = 90^\circ - \alpha$, je $|\angle BCO| = 90^\circ - \alpha \leq |\angle BCP|$, a~jelikož $O$ i~$P$ leží uvnitř trojúhelníku, tedy v~téže polorovině s~hraniční přímkou $BC$, leží uvnitř úhlu $BCP$ i~polopřímka $CO$ a
    $$
    |\angle BCP| = |\angle BCO| + |\angle OCP|.
    $$
    Porovnáním obou rozkladů úhlu $BCP$ a~vykrácením shodných úhlů $|\angle QCP| = |\angle BCO|$ dostaneme $|\angle BCQ| = |\angle OCP|$, což je dokazovaná rovnost.

    Jelikož $Q \neq P$ leží na úsečce $BP$, splývají polopřímky $PQ$ a~$PB$, a~tak $|\angle QPC| = |\angle BPC| = 90^\circ$. V~pravoúhlém trojúhelníku $QPC$ je proto rovnost $|\angle PCQ| = 90^\circ - \alpha$ ekvivalentní $|\angle CQP| = \alpha$; stačí tedy dokázat $|\angle CQP| = |\angle BAC|$.

    Pravý úhel u~$P$ si žádá doplnění pravoúhlého trojúhelníku $QPC$ na rovnoramenný: nechť $Q'$ je obraz $Q$ ve středové souměrnosti podle $P$. Přímka $CP$ je osou úsečky $QQ'$, takže $|CQ| = |CQ'|$ a~$|\angle CQP| = |\angle CQ'P|$. Bod $Q$ leží mezi $B$ a~$P$ a~bod $P$ mezi $Q$ a~$Q'$, čili na polopřímce $BP$ jdou body v~pořadí $Q$, $P$, $Q'$; polopřímky $Q'P$ a~$Q'B$ proto splývají a~$|\angle CQ'P| = |\angle CQ'B|$. Přímky $AB$ a~$BC$ procházejí bodem $B$, takže každý bod polopřímky $BP$ kromě $B$ leží vůči nim v~téže polorovině jako vnitřní bod $P$: bod $Q'$ tak leží v~téže polorovině jako $C$ vzhledem k~$AB$ a~v~téže jako $A$ vzhledem k~$BC$. Stačí proto dokázat, že $Q'$ leží na kružnici opsané trojúhelníku $ABC$: z~obvodových úhlů nad tětivou $BC$ pak $|\angle CQ'B| = |\angle CAB| = \alpha$, a~tedy $|\angle CQP| = \alpha$.

    Z~pořadí bodů na $\omega$ tětiva $AB$ neodděluje $P$ od $O$ a~tětiva $OB$ neodděluje $P$ od $A$, takže z~obvodových úhlů
    $$
    \gather
    |\angle APB| = |\angle AOB| = 2\gamma, \\
    |\angle OPB| = |\angle OAB| = 90^\circ - \gamma.
    \endgather
    $$
    Polopřímky $PQ$ a~$PB$ splývají, čili $|\angle APQ| = 2\gamma$. Přímka $BP$ je příčkou rovnoběžek $QA$ a~$PO$ a~body $A$, $O$ leží vůči ní v~opačných polorovinách, takže jde o~střídavé úhly a~$|\angle AQP| = |\angle OPB| = 90^\circ - \gamma$. Ze součtu úhlů v~trojúhelníku $APQ$ je
    $$
    |\angle QAP| = 180^\circ - 2\gamma - (90^\circ - \gamma) = 90^\circ - \gamma = |\angle AQP|,
    $$
    takže trojúhelník $APQ$ je rovnoramenný a~$|PA| = |PQ| = |PQ'|$.

    Rovnoramenný je tím pádem i~trojúhelník $APQ'$, přičemž $P$ leží mezi $Q$ a~$Q'$, takže jeho úhel u~$P$ je $|\angle APQ'| = 180^\circ - |\angle APQ| = 180^\circ - 2\gamma$. Zbylé dva úhly jsou shodné, čili $|\angle AQ'P| = \gamma$. Polopřímky $Q'P$ a~$Q'B$ splývají, a~tak $|\angle AQ'B| = \gamma = |\angle ACB|$. Body $Q'$ a~$C$ leží v~téže polorovině s~hraniční přímkou $AB$, proto $Q'$ leží na oblouku $AB$ kružnice opsané trojúhelníku $ABC$ obsahujícím $C$.
}

\EnvId{GK8F64_JP-rSXryIXL_wp-petiuhelnik-a-dva-stredy-kruznic}
\Problem{0}{MEMO 2017, I3}{
    Je dán konvexní pětiúhelník $ABCDE$. Označme $P$ průsečík přímek~$CE$ a~$BD$. Dokažte, že pokud $|\angle PAD| = |\angle ACB|$ a $|\angle CAP| = |\angle EDA|$, pak~$P$ leží na přímce určené středy kružnic opsaných trojúhelníkům~$ABC$ a~$ADE$.
}{
    Když nevíme, co dělat, chvíli počítejte úhly. To by nám mělo umožnit najít v obrázku jednoduchou pěknou vlastnost, jež může vnuknout, co přidat.
}{
    Klíčem je objevit $BC \parallel DE$. Obrázek obsahuje průsečík~$CE$ a~$BD$. To může vnuknout přidání obrazu~$A$ ve stejnolehlosti zobrazující~$BC$ na~$DE$. Najednou se vše stane krásným.
}{
    Označme $\omega_1$ a~$\omega_2$ po řadě kružnice opsané trojúhelníkům $ABC$ a~$ADE$ a jejich středy $O_1$, $O_2$.

    Nejprve si ujasněme polohu bodu $P$. Úhlopříčka $AC$ odděluje $B$ od bodů $D$, $E$ a úhlopříčka $AD$ odděluje $E$ od bodů $B$, $C$. Bod $P$ je vnitřním bodem úsečky $CE$, leží tedy v otevřené polorovině určené přímkou $AC$ a bodem $E$, a zároveň je vnitřním bodem úsečky $BD$, takže leží v otevřené polorovině určené přímkou $AD$ a bodem $B$. Leží proto uvnitř úhlu $\angle CAD$ a platí
    $$
    |\angle CAP| + |\angle PAD| = |\angle CAD|.
    $$

    Úhlopříčky $CA$ a~$DA$ leží uvnitř vnitřních úhlů pětiúhelníku při vrcholech $C$ a~$D$, takže
    $$
    \gather
    |\angle BCD| = |\angle BCA| + |\angle ACD|, \\
    |\angle CDE| = |\angle CDA| + |\angle ADE|.
    \endgather
    $$
    Podle zadání je $|\angle BCA| = |\angle PAD|$ a~$|\angle ADE| = |\angle CAP|$, ze součtu úhlů v trojúhelníku $ACD$ je $|\angle ACD| + |\angle CDA| = 180^\circ - |\angle CAD|$, a tedy
    $$
    |\angle BCD| + |\angle CDE| = |\angle CAD| + 180^\circ - |\angle CAD| = 180^\circ.
    $$
    Body $B$ a~$E$ leží v téže polorovině s hraniční přímkou $CD$, takže $\angle BCD$ a~$\angle CDE$ jsou vnitřní úhly u příčky $CD$ na téže straně. Jejich součet $180^\circ$ dává
    $$
    BC \parallel DE.
    $$

    Dvě rovnoběžné úsečky vybízejí k~úvaze o~stejnolehlosti, která jednu zobrazí na druhou a jíž přeneseme i~další bod. Její střed už v obrázku máme: přímky $BD$ a~$CE$ se protínají právě v~$P$. Čtyřúhelník $BCDE$ je konvexní, takže $P$ leží uvnitř obou úseček $BD$ i $CE$, a stejnolehlost $h$ se středem $P$ a koeficientem $k = -|PD| \, / \, |PB| < 0$ zobrazuje $B$ na bod $D$. Přímku $BC$ zobrazí na rovnoběžku s ní vedenou bodem $D$, čili na přímku $DE$, a přímku $CE$ na sebe samu, neboť ta prochází středem $P$. Proto $h(C) = E$.

    Přidaným bodem je obraz $A' = h(A)$. Leží-li na $\omega_2$, jsme hotovi: $\omega_2$ je pak kružnice opsaná trojúhelníku $A'DE$, tedy obrazu trojúhelníku $ABC$ v~$h$, čili $h(\omega_1) = \omega_2$, a stejnolehlost zobrazuje střed na střed.

    Jelikož $k < 0$, leží $P$ mezi $A$ a~$A'$ a polopřímky $A'A$, $A'P$ splývají. Stejnolehlost $h$ zobrazuje trojúhelník $APC$ na trojúhelník $A'PE$ (bod $P$ je jejím pevným bodem), a tak
    $$
    |\angle AA'E| = |\angle PA'E| = |\angle PAC| = |\angle ADE|,
    $$
    kde poslední rovnost je druhá podmínka ze zadání.

    Zbývá ověřit polohu bodu $A'$. Bod $P$ je vnitřním bodem úhlopříčky $BD$, leží tedy uvnitř pětiúhelníku, čili v otevřené polorovině s hraniční přímkou $AE$, v níž leží i vrchol $D$. Bod $A'$ leží na polopřímce $AP$ a je různý od $A$, takže leží v téže polorovině. Body $D$ a~$A'$ proto vidí úsečku $AE$ pod stejným úhlem a leží u~ní na téže straně, takže $A$, $E$, $D$, $A'$ leží na jedné kružnici. Tou je $\omega_2$, a tedy $A'$ na $\omega_2$ leží.

    Trojúhelník $A'DE$ je obrazem trojúhelníku $ABC$, jeho vrcholy tedy nejsou kolineární a kružnice $h(\omega_1)$ je jimi jednoznačně určena. Všechny tři leží na $\omega_2$, proto $h(\omega_1) = \omega_2$ a~$h(O_1) = O_2$. Je-li $O_1 = O_2$, je $O_1$ pevným bodem stejnolehlosti s koeficientem $k \neq 1$, čili $O_1 = O_2 = P$ a přímka ze zadání není určena. V opačném případě je $O_1 \neq P$ a body $P$, $O_1$, $O_2 = h(O_1)$ leží z definice stejnolehlosti na jedné přímce.
}

\EnvId{NsM3t1SwEv7CD5wXQgwbU-usecka-pq-rovnobezna-s-bc}
\Problem{0}{IMO 2019, P2}{
    V trojúhelníku~$ABC$ nechť~$B_1$ je bod na straně~$AC$ a~$C_1$ bod na straně~$AB$. Nechť $P$ a~$Q$ jsou body po řadě na úsečkách~$BB_1$ a~$CC_1$ tak, že přímka~$PQ$ je rovnoběžná s~$BC$. Nechť~$P_1$ je bod na přímce~$PC_1$ takový, že~$C_1$ leží ostře mezi~$P$ a~$P_1$ a $|\angle PP_1A|=|\angle CBA|$. Nechť~$Q_1$ je bod na přímce~$QB_1$ takový, že~$B_1$ leží ostře mezi~$Q$ a~$Q_1$ a $|\angle AQ_1Q|=|\angle ACB|$. Dokažte, že $P$, $Q$, $P_1$, $Q_1$ leží na jedné kružnici.
}{
    Cílem je přidat něco, co vyjasní zvláštní úhlovou podmínku pro~$P_1$ a~$Q_1$. Úloha však má příliš mnoho bodů a různé věci, které vypadají slibně, lze přidat bez valné pomoci, například průsečík~$P_1P$ s~$BC$ nebo~$PQ$ s~$AB$ (zkuste to). Potíž s těmito průsečíky je, že sice dají tětivový čtyřúhelník, ale ten neposkytne žádné další užitečné úhly k přenesení. Je proto lepší protnout nějaké přímky s kružnicí tak, aby to do rovnosti~$|\angle PP_1A|=|\angle CBA|$ vneslo ještě jeden úhel.
}{
    Klíčovým kouzelným bodem je druhý průsečík přímky~$CC_1$ s kružnicí opsanou trojúhelníku~$ABC$, nazvěme ho~$C_2$. Díky němu dostaneme $|\angle CBA|=|\angle CC_2A|$, což kouzelně dá tětivový čtyřúhelník $C_1C_2P_1A$. Bod~$B_2$ definujeme analogicky. A věřte nebo ne, teď už je to jen středně těžká úloha na počítání úhlů.
}{
    Označme $\omega$ kružnici opsanou trojúhelníku $ABC$, $\beta = |\angle CBA|$ a $\gamma = |\angle ACB|$.

    Úhel $\beta$ z podmínky pro bod $P_1$ je v obrázku vidět u vrcholu $B$, tedy daleko od bodů $C_1$ a~$P_1$, kterých se podmínka týká. Zároveň je to obvodový úhel nad tětivou $AC$, takže je vidět z každého bodu oblouku $AC$ obsahujícího $B$. Stačí tedy protnout s~$\omega$ vhodnou přímku. Tou pravou je $CC_1$: leží na ní $C_1$ i $Q$. Nechť je tedy $C_2$ její druhý průsečík s~$\omega$ a analogicky nechť $B_2$ je druhý průsečík přímky $BB_1$ s~$\omega$.

    Body $B_1$, $C_1$ jsou vnitřní body stran $AC$, $AB$ (pro $C_1 = A$ by $P_1$ ležel na přímce $PA$ a úhel $|\angle PP_1A|$ by byl nulový nebo přímý, pro $C_1 = B$ by $Q$ ležel na přímce $BC$ a přímka $PQ$ by s ní splynula; pro $B_1$ analogicky), takže leží uvnitř $\omega$. Kdyby $P = B$, byla by přímka $PQ$ rovnoběžkou s~$BC$ vedenou bodem $B$, tedy samotnou přímkou $BC$; proto $P \neq B$ a stejně tak $Q \neq C$. Úsečka $BB_1$ leží kromě koncového bodu $B$ uvnitř $\omega$, takže uvnitř $\omega$ leží $P$, a stejně tak $Q$.

    Přímka $CC_1$ protíná přímku $AB$ jedině v jejím vnitřním bodě $C_1$, který leží uvnitř $\omega$; na přímce $CC_1$ je proto pořadí bodů $C$, $C_1$, $C_2$ a bod $C_2$ leží na otevřeném oblouku $AB$ neobsahujícím $C$. Ten je částí oblouku $AC$ obsahujícího $B$, takže z obvodových úhlů nad tětivou $AC$ dostáváme
    $$
    |\angle AC_2C| = |\angle ABC| = \beta.
    $$

    Jelikož $B_1$ je vnitřní bod strany $AC$, leží celá úsečka $BB_1$ kromě bodu $B$ v otevřené polorovině s hraniční přímkou $AB$ obsahující $C$; speciálně tam leží $P$. Bod $C_1$ leží na přímce $AB$ a mezi $P$ a~$P_1$, takže $P_1$ leží v opačné otevřené polorovině, a v té leží i $C_2$, neboť $C_1$ leží mezi $C$ a~$C_2$. Z toho, že $C_1$ leží na polopřímce $P_1P$ i na polopřímce $C_2C$, dále plyne
    $$
    |\angle AP_1C_1| = |\angle AP_1P| = \beta = |\angle AC_2C| = |\angle AC_2C_1|.
    $$
    Body $P_1$ a~$C_2$ tedy vidí úsečku $AC_1$ pod stejným úhlem a leží na téže straně přímky $AC_1$, takže podle obrácené věty o obvodovém úhlu leží $A$, $C_1$, $C_2$, $P_1$ na jedné kružnici. Analogicky, s úhlem $\gamma$ místo $\beta$, leží na jedné kružnici body $A$, $B_1$, $B_2$, $Q_1$, přičemž $B_2$ leží na otevřeném oblouku $AC$ neobsahujícím $B$.

    Bod $Q$ leží na úsečce $CC_1$, tedy v uzavřené polorovině přímky $AB$ obsahující $C$; totéž platí pro $P$. Bod $P_1$ leží v opačné otevřené polorovině, takže $P_1 \neq P$ i $P_1 \neq Q$, a analogicky $Q_1 \neq P$ i $Q_1 \neq Q$.

    Dále už jen počítáme úhly. Abychom se vyhnuli rozboru vzájemných poloh, přejdeme k orientovaným úhlům modulo $180^\circ$: $\angle(\ell,m)$ je úhel, o který je třeba otočit přímku $\ell$, aby byla rovnoběžná s přímkou $m$. Použijeme dvě standardní vlastnosti: čtyři navzájem různé body $X$, $Y$, $Z$, $W$ leží na jedné kružnici nebo přímce právě tehdy, když $\angle(XZ,XW) = \angle(YZ,YW)$, a rovnoběžné přímky svírají s každou další přímkou týž orientovaný úhel.

    Přímky $BB_1$ a~$CC_1$ protínají přímku $BC$, takže nejsou rovnoběžné s~$PQ$ a s přímkou $PQ$ mají společný jediný bod, po řadě $P$ a~$Q$. Body $B_2$, $C_2$ leží na $\omega$, kdežto $P$ a~$Q$ uvnitř ní, takže ani $B_2$, ani $C_2$ neleží na přímce $PQ$; oblouky, na nichž leží, jsou navíc disjunktní, takže $B_2 \neq C_2$. Označme $\Omega$ kružnici procházející body $P$, $Q$, $C_2$. V následujících krocích je tak vždy možnost \uv{nebo na přímce} vyloučena.

    Nejprve ukážeme, že na $\Omega$ leží i $B_2$. Bod $P$ leží na přímce $BB_2$, bod $Q$ na přímce $CC_2$ a $PQ \parallel BC$, takže
    $$
    \eqalign{
        \angle(B_2P, B_2C_2) &= \angle(B_2B, B_2C_2) = \angle(CB, CC_2) \cr
        &= \angle(QP, QC_2),
    }
    $$
    kde prostřední rovnost jsou obvodové úhly nad tětivou $BC_2$ kružnice $\omega$. Body $P$, $Q$, $C_2$, $B_2$ tedy leží na $\Omega$.

    Nyní bod $P_1$; pro $P_1 = C_2$ není co dokazovat, nechť je tedy $P_1 \neq C_2$. Bod $C_1$ leží na přímce $P_1P$ i na přímce $AB$ a body $A$, $C_1$, $C_2$, $P_1$ leží na jedné kružnici, takže
    $$
    \eqalign{
        \angle(P_1P, P_1C_2) &= \angle(P_1C_1, P_1C_2) = \angle(AC_1, AC_2) \cr
        &= \angle(AB, AC_2) = \angle(CB, CC_2) = \angle(QP, QC_2),
    }
    $$
    kde předposlední rovnost jsou obvodové úhly nad tětivou $BC_2$ kružnice $\omega$ a poslední pochází z předchozího výpočtu. Bod $P_1$ proto leží na kružnici $\Omega$.

    S bodem $Q_1$ naložíme symetricky, jen místo $C_2$ použijeme $B_2$ (pro $Q_1 = B_2$ opět není co dokazovat). Bod $B_1$ leží na přímce $Q_1Q$ i na přímce $AC$ a body $A$, $B_1$, $B_2$, $Q_1$ leží na jedné kružnici, takže
    $$
    \eqalign{
        \angle(Q_1Q, Q_1B_2) &= \angle(Q_1B_1, Q_1B_2) = \angle(AB_1, AB_2) \cr
        &= \angle(AC, AB_2) = \angle(BC, BB_2) = \angle(PQ, PB_2),
    }
    $$
    kde předposlední rovnost jsou obvodové úhly nad tětivou $CB_2$ kružnice $\omega$ a poslední využívá $PQ \parallel BC$ a to, že $P$ leží na přímce $BB_2$. Bod $Q_1$ tak leží na kružnici procházející body $P$, $Q$, $B_2$, což je $\Omega$.

    Na kružnici $\Omega$ tedy leží body $P$, $Q$, $B_2$, $C_2$, $P_1$ i~$Q_1$; speciálně leží na jedné kružnici body $P$, $Q$, $P_1$ a~$Q_1$.
}

\EnvId{sIGVitSRoddubnK736EUu-bod-na-vysce-pravouhleho-trojuhelniku}
\Problem{0}{IMO 2012, P5}{
    V pravoúhlém trojúhelníku~$ABC$ s pravým úhlem u vrcholu~$A$ nechť~$D$ je pata výšky z~$A$. Nechť~$X$ je libovolný vnitřní bod úsečky~$AD$. Nechť~$K$ je bod na úsečce~$BX$ takový, že $|CK| = |CA|$. Podobně nechť~$L$ je bod na úsečce~$CX$ takový, že $|BL| = |BA|$. Označme~$M$ průsečík přímek~$BL$ a~$CK$. Dokažte, že $|MK| = |ML|$.
}{
    V konfiguraci nelze vidět nic smysluplného. Něco se tam skrývá, ale není snadné vidět přesně co. Jednou z možností je zkusit udělat něco smysluplného s přímkami~$BX$ a~$CX$. Body~$K$ a~$L$ na nich leží a nevidíme, jak to využít. Jednou z možností je protnout je s něčím smysluplným.
}{
    Klíčovými body jsou průsečíky~$E$ a~$F$ přímek~$BX$ a~$CX$ s kružnicí nad průměrem~$BC$. Ani ty samotné nestačí, ale další bod je docela přirozený: průsečík~$G$ přímek~$BF$ a~$CE$. Ten zřejmě leží na přímce~$AXD$ a~$X$ je ortocentrum trojúhelníku~$GBC$. Ačkoli se to tak nemusí zdát, úlohu lze nyní dotáhnout, a to kombinací mocnosti bodu, obvodových úhlů a tětivových čtyřúhelníků. I v této fázi to není snadná úloha.
}{
    Z předpokladů o bodech~$K$ a~$L$ využijeme jen tolik: $K$ leží na přímce~$BX$ a $|CK| = |CA|$, tedy leží na kružnici~$\omega_c$ se středem~$C$ a poloměrem $|CA|$; podobně $L$ leží na přímce~$CX$ a na kružnici~$\omega_b$ se středem~$B$ a poloměrem $|BA|$. Obě kružnice procházejí bodem~$A$.

    Jediné, co body~$K$ a~$L$ v obrázku drží, jsou přímky~$BX$ a~$CX$, a ty samy o sobě uchopit nelze. Protněme je proto s kružnicí, která o konfiguraci něco ví: s Thaletovou kružnicí~$\tau$ nad průměrem~$BC$. Leží na ní i bod~$A$, neboť úhel při~$A$ je pravý.

    Jelikož úhly při~$B$ a~$C$ jsou ostré, leží bod~$D$ uvnitř úsečky~$BC$; bod~$X$ tedy neleží na přímce~$BC$ a přímky~$BX$, $CX$ jsou různé. Ani jedna z nich není kolmá na~$BC$: kolmice na~$BC$ vedená bodem~$B$ je rovnoběžná s přímkou~$AD$ a od ní různá (bod~$B$ na~$AD$ neleží), takže bod~$X$ neobsahuje, a stejně pro~$C$. Přímka~$BX$ proto není tečnou~$\tau$ v bodě~$B$ a protíná $\tau$ ještě v bodě $E \neq B$; přitom $E \neq C$, neboť jinak by $X$ ležel na~$BC$. Analogicky nechť~$F$ je druhý průsečík přímky~$CX$ s~$\tau$, opět $F \neq B$ a $F \neq C$. Podle Thaletovy věty je
    $$
    |\angle BEC| = |\angle BFC| = 90^\circ,
    $$
    tedy $CE \perp BX$ a $BF \perp CX$.

    Přímka~$CE$ je kolmice z~$C$ na~$BX$, tedy osa tětivy, kterou kružnice~$\omega_c$ vytíná na přímce~$BX$; podobně $BF$ je osa tětivy kružnice~$\omega_b$ na přímce~$CX$. Jejich průsečík je přirozený kandidát na střed kružnice procházející body~$K$ i~$L$, a~to je bod, který úlohu rozlouskne. Označme jej~$G$. Existuje a je jediný, neboť přímky~$BX$, $CX$ jsou různé a protínají se v~$X$, takže nejsou rovnoběžné, a tedy nejsou rovnoběžné ani kolmice~$CE$, $BF$ na ně.

    Body~$G$, $B$, $C$ tvoří trojúhelník. Bod~$E$ leží na~$\tau$ a je různý od~$B$ i od~$C$, takže neleží na přímce~$BC$ a přímky~$CE$, $BC$ se protínají jedině v~$C$; stejně $F$ neleží na~$BC$, tedy $C$ neleží na přímce~$BF$. Kdyby $G$ ležel na~$BC$, z prvního by bylo $G = C$, což druhé vylučuje.

    Přímka~$BX$ prochází vrcholem~$B$ a je kolmá na přímku $CG = CE$, je to tedy výška trojúhelníku~$GBC$ z vrcholu~$B$; stejně $CX$ je jeho výška z vrcholu~$C$. Bod~$X$ je proto ortocentrum trojúhelníku~$GBC$ a třetí výška dává $GX \perp BC$. Body~$G$ a~$X$ jsou přitom různé: kdyby $G = X$, ležel by $X$ na přímce~$CE$ i na přímce~$BX$, ty jsou na sebe kolmé a protínají se jedině v~$E$, tedy by bylo $X = E$; jenže $A$ leží na~$\tau$ a~$D$ uvnitř ní, takže celá úsečka~$AD$ kromě bodu~$A$ leží uvnitř~$\tau$, kdežto $E$ leží na~$\tau$. Přímky~$GX$ a~$AD$ jsou tak obě kolmé na~$BC$ a obě procházejí bodem~$X$, tedy splývají a bod~$G$ leží na přímce~$AD$.

    Teď už jen počítáme, a to stále týmž nástrojem: je-li přímka~$p$ kolmá na přímku~$q$ v bodě~$T$, pak pro libovolné body~$Y$ na~$p$ a~$Z$ na~$q$ platí $|YZ|^2 = |YT|^2 + |TZ|^2$ (pro $Y = T$ nebo $Z = T$ triviálně). Použijeme to v bodech~$D$, $E$, $F$ a nakonec v~$K$ a~$L$.

    Přímka~$AD$ je kolmá na~$BC$ v bodě~$D$ a leží na ní body~$A$ i~$G$. Rozdíly pro dvojice $(G,B)$, $(G,C)$, resp. $(A,B)$, $(A,C)$ proto dávají
    $$
    \gather
    |GB|^2 - |GC|^2 = |DB|^2 - |DC|^2 = \\
    = |AB|^2 - |AC|^2.
    \endgather
    $$
    To je klíčová identita a~dá se číst dvěma způsoby. Za prvé ji lze přepsat jako
    $$
    |GB|^2 - |BA|^2 = |GC|^2 - |CA|^2,
    $$
    což říká, že mocnost bodu~$G$ ke kružnici~$\omega_b$ se rovná jeho mocnosti ke kružnici~$\omega_c$, tedy že $G$ leží na jejich chordále (kterou je právě přímka~$AD$). Společnou hodnotu označme~$k$. Za druhé, jelikož úhel při~$A$ je pravý a tedy $|AB|^2 + |AC|^2 = |BC|^2$, dá se přepsat i~do dvou tvarů
    $$
    \gather
    |GB|^2 = |GC|^2 + |BC|^2 - 2|CA|^2, \\
    |GC|^2 = |GB|^2 + |BC|^2 - 2|BA|^2.
    \endgather
    $$

    Přímka~$CE$ je kolmá na~$BX$ v bodě~$E$, přičemž $C$ i~$G$ leží na~$CE$ a~$B$ i~$K$ na~$BX$. Rozdíly pro dvojice $(G,K)$, $(C,K)$, resp. $(G,B)$, $(C,B)$ dávají
    $$
    \gather
    |GK|^2 - |CK|^2 = |GE|^2 - |CE|^2 = \\
    = |GB|^2 - |CB|^2,
    \endgather
    $$
    neboť v první dvojici se odečte $|EK|^2$ a ve druhé $|EB|^2$. Dosadíme $|CK| = |CA|$ a první ze dvou tvarů klíčové identity:
    $$
    \gather
    |GK|^2 = |GB|^2 - |CB|^2 + |CA|^2 = \\
    = |GC|^2 - |CA|^2 = k.
    \endgather
    $$
    Bod~$F$ dá totéž se záměnou $B \leftrightarrow C$, $K \leftrightarrow L$, $E \leftrightarrow F$, která celou konfiguraci zachovává: přímka~$BF$ je kolmá na~$CX$ v bodě~$F$, na~$BF$ leží $B$ i~$G$, na~$CX$ leží $C$ i~$L$, a ze vztahu $|BL| = |BA|$ a druhého tvaru klíčové identity vyjde
    $$
    |GL|^2 = |GC|^2 - |BC|^2 + |BA|^2 = k.
    $$
    Bod~$G$ je tedy od~$K$ i od~$L$ vzdálen $\sqrt k$.

    Je-li $k = 0$, je $K = G = L$ a tvrzení platí triviálně. Nechť tedy $k > 0$ a nechť~$\gamma$ je kružnice se středem~$G$ a poloměrem $\sqrt k$; body~$K$, $L$ na ní leží. Z $|GK|^2 = |GC|^2 - |CA|^2$ a $|CK| = |CA|$ dostaneme
    $$
    |GK|^2 + |CK|^2 = |GC|^2,
    $$
    takže podle obrácené Pythagorovy věty je $|\angle GKC| = 90^\circ$ a přímka~$CK$ je tečnou~$\gamma$ v bodě~$K$; stejně je $BL$ tečnou~$\gamma$ v bodě~$L$. Bod~$M$ leží na obou tečnách, takže poslední použití téže Pythagorovy věty, v bodě~$K$ s přímkami $GK \perp CK$ a poté v bodě~$L$ s přímkami $GL \perp BL$, dává
    $$
    |MK|^2 = |MG|^2 - k = |ML|^2.
    $$
}

\secc Moje oblíbené

Pár konstrukcí, které mi připadají obzvlášť rozkošné. Ne nutně nejtěžší úlohy zde \SlightSmile{}

\EnvId{JLEqfNmB2Hw0P5oc4qFwQ-body-na-vyskach-a-obsahy}
\Problem{0}{}{
    Nechť $ABC$ je ostroúhlý trojúhelník. Body $X$, $Y$, $Z$ leží po řadě na výškách z $A$, $B$, $C$, uvnitř trojúhelníku, tak, že $A_{BCX} + A_{CAY} + A_{ABZ} = S_{ABC}$, kde $A_{PQR}$ značí obsah trojúhelníku $PQR$. Dokažte, že kružnice opsaná trojúhelníku $XYZ$ prochází ortocentrem trojúhelníku $ABC$.
}{
    Když máme součet obsahů, je lepší ho pěkně interpretovat. Jenže tyto trojúhelníky se nešikovně překrývají. Tak raději tyto obsahy \uv{přesuneme} jinam.
}{
    Trik je v tom, že je-li $X'$ libovolný bod na přímce procházející $X$ rovnoběžně s $BC$, pak $A_{BCX}=A_{BCX'}$. Na co nás to navádí?
}{
    Zřejmě dává smysl uvážit přímky rovnoběžné se stranami procházející našimi body a protnout je; zvolme přímku procházející $Y$ rovnoběžně s $AC$ a přímku procházející $Z$ rovnoběžně s $AB$ a nechme je protnout v $T$. Co teď říká podmínka na obsahy? Prakticky jsme u cíle.
}{
    Symbol $S_{ABC}$ ze zadání je obsah trojúhelníku $ABC$, tedy $A_{ABC}$; dále píšeme důsledně $A_{PQR}$.

    Pro bod $P$ roviny označme $A_{BCP}$ orientovaný obsah trojúhelníku $BCP$: obyčejný obsah se znaménkem $+$, leží-li $P$ ve stejné polorovině s~hraniční přímkou $BC$ jako bod $A$, se znaménkem $-$, leží-li v~opačné, a~nulu, leží-li $P$ na přímce $BC$. Analogicky definujeme $A_{CAP}$ (znaménko podle $B$) a~$A_{ABP}$ (podle $C$). Body $X$, $Y$, $Z$ leží uvnitř trojúhelníku, takže pro ně orientovaný obsah splývá s~obyčejným a~podmínka ze zadání se nemění.

    Platí $A_{BCP} = \frac12 |BC| \cdot d(P)$, kde $d(P)$ je orientovaná vzdálenost bodu $P$ od přímky $BC$. Pro dva body je proto $A_{BCP} = A_{BCP'}$ právě tehdy, když $P$ a~$P'$ leží na téže rovnoběžce s~přímkou $BC$; analogicky pro $CA$ a~$AB$. Dále pro každý bod $P$ platí
    $$
    A_{BCP} + A_{CAP} + A_{ABP} = A_{ABC}.
    $$
    Orientovaná vzdálenost od přímky je afinní funkce bodu, takže afinní je i levá strana; v~bodech $A$, $B$, $C$ nabývá hodnoty $A_{ABC}$ a~afinní funkce je svými hodnotami ve třech nekolineárních bodech určena jednoznačně.

    Tři trojúhelníky z~podmínky se nepřehledně překrývají, obsah $A_{BCX}$ se však nezmění, posuneme-li $X$ po rovnoběžce s~$BC$. Všechny tři obsahy tedy umíme přesunout tak, aby příslušely jedinému bodu: stačí protnout vhodné rovnoběžky.

    Nechť $T$ je průsečík přímky vedené bodem $Y$ rovnoběžně s~$CA$ a~přímky vedené bodem $Z$ rovnoběžně s~$AB$; přímky $CA$ a~$AB$ nejsou rovnoběžné, takže $T$ existuje a~je jediný. Z~kritéria výše je $A_{CAT} = A_{CAY}$ a~$A_{ABT} = A_{ABZ}$, čili
    $$
    \eqalign{
        A_{BCT} &= A_{ABC} - A_{CAT} - A_{ABT} \cr
        &= A_{ABC} - A_{CAY} - A_{ABZ} = A_{BCX},\cr
    }
    $$
    kde poslední rovnost je právě podmínka ze zadání. Body $X$ a~$T$ tak leží na téže rovnoběžce s~$BC$.

    Označme $H$ ortocentrum trojúhelníku $ABC$. Body $X$, $H$ leží na výšce z~$A$, která je kolmá na $BC$; podobně $Y$, $H$ leží na výšce z~$B$ kolmé na $CA$ a~$Z$, $H$ na výšce z~$C$ kolmé na $AB$.

    Kdyby bylo $T = H$, ležel by bod $H$ s~bodem $X$ na téže rovnoběžce s~$BC$ a~zároveň na výšce z~$A$; tyto dvě přímky jsou na sebe kolmé, takže mají jediný společný bod a~$X = H$. Stejně by bylo $Y = H$ a~$Z = H$, což je spor s~tím, že $XYZ$ je trojúhelník. Proto $T \neq H$ a~kružnice $k$ s~průměrem $HT$ je dobře definovaná.

    Bod $X$ leží na $k$: pro $X = H$ i~$X = T$ je to zřejmé, jinak je přímka $XH$ výškou z~$A$ a~přímka $XT$ rovnoběžkou s~$BC$, ty jsou na sebe kolmé, takže $|\angle HXT| = 90^\circ$ a~$X$ leží na $k$ podle Thaletovy věty. Stejně leží na $k$ body $Y$ (výška z~$B$ je kolmá na $CA$) a~$Z$ (výška z~$C$ je kolmá na $AB$).

    Kružnice $k$ prochází body $X$, $Y$, $Z$, a~jelikož ty tvoří trojúhelník, je právě $k$ jeho opsanou kružnicí. Ta tedy prochází bodem $H$.
}

\EnvId{V8qpCa8IyPgddyS2NJS8Y-trojuhelnik-s-obvodem-ctyri}
\Problem{0}{}{
    Nechť $ABC$ je trojúhelník s obvodem $4$. Body $P$ a $Q$ leží po řadě na polopřímkách $AB$ a $AC$ tak, že $|AP| = |AQ| = 1$ a úsečky $BC$ a $PQ$ se protínají v bodě $X$. Dokažte, že jeden ze dvou trojúhelníků, na které $AX$ dělí trojúhelník $ABC$, má obvod $2$.
}{
    První trik je vzpomenout si na jedno obzvlášť pěkné místo, kde se v geometrii trojúhelníku objevuje obvod (nebo možná jeho polovina)?
}{
    Vzdálenost z $A$ k bodům dotyku kružnice připsané proti vrcholu $A$ je přesně polovina obvodu, tedy 2. A $|AP|=|AQ|=1$. Je to náhoda?
}{
    Není to náhoda (jinak bych se neptal). Přímka $DE$ je obrazem $A$-poláry kružnice připsané proti vrcholu $A$ ve stejnolehlosti se středem $A$ a koeficientem $1/2$. Co o této přímce víme?
}{
    Inu, tato přímka je chordálou $A$ a kružnice připsané proti vrcholu $A$! A $F$ na ní leží. Blížíme se.
}{
    Posledním bodem, který narýsovat, je bod dotyku připsané kružnice a $BC$. S využitím toho, že $F$ leží na této chordále, dostaneme nějaké pěkné stejné délky.
}{
    Obvod je 4, čili polovina obvodu je 2, a~$|AP| = |AQ| = 1$ je přesně její polovina. Polovina obvodu se v~geometrii trojúhelníku objevuje jako vzdálenost z~vrcholu k~bodům dotyku kružnice připsané proti tomuto vrcholu, přidejme si ji tedy do obrázku.

    Označme $a = |BC|$, $b = |CA|$, $c = |AB|$, takže $a + b + c = 4$, a~nechť $\omega$ je kružnice připsaná trojúhelníku $ABC$ proti vrcholu $A$. Ta se dotýká strany $BC$ v~bodě $U$ a~prodloužení stran $AB$, $AC$ za body $B$, $C$ po řadě v~bodech $V$, $W$. Mocnost bodu $M$ vzhledem ke kružnici $\gamma$ označme $p(M,\gamma)$.

    Z~rovnosti délek tečen z~jednoho bodu je $|BV| = |BU|$, $|CW| = |CU|$ a~$|AV| = |AW|$, a~jelikož $U$ leží na úsečce $BC$, dostáváme
    $$
    |AV| + |AW| = (c + |BU|) + (b + |CU|) = a + b + c = 4.
    $$
    Proto $|AV| = |AW| = 2$, a~tedy i~$|BU| = |BV| = 2 - c$ a~$|CU| = |CW| = 2 - b$.

    Body $P$, $V$ leží na téže polopřímce z~$A$ a~$|AP| = 1 = \tfrac12|AV|$, takže $P$ je středem úsečky $AV$; stejně je $Q$ středem úsečky $AW$. Přímka $PQ$ je tedy obrazem přímky $VW$, čili poláry bodu $A$ vzhledem k~$\omega$, ve stejnolehlosti se středem $A$ a~koeficientem $\tfrac12$. To není náhoda a~dá se to říct mnohem užitečněji: přímka $PQ$ je chordálou kružnice $\omega$ a~kružnice s~nulovým poloměrem a~středem $A$.

    Mocnost bodu $M$ vzhledem k~této degenerované kružnici je $|MA|^2$ a~standardní důkaz, že body se stejnou mocností vzhledem ke dvěma nesoustředným kružnicím tvoří přímku kolmou na spojnici jejich středů, funguje i~pro nulový poloměr. Bod $A$ leží vně $\omega$, neboť $p(A,\omega) = |AV|^2 = 4 > 0$, takže není středem $\omega$ a~tato chordála opravdu je přímka. Přímka $AB$ se dotýká $\omega$ v~bodě $V$, proto
    $$
    p(P,\omega) = |PV|^2 = |PA|^2,
    $$
    kde druhá rovnost platí právě proto, že $P$ je středem úsečky $AV$; stejně $p(Q,\omega) = |QW|^2 = |QA|^2$. Body $P$, $Q$ jsou různé, neboť leží ve vzdálenosti 1 na různých polopřímkách z~$A$, a~hledanou chordálou je tak právě přímka $PQ$.

    Bod $X$ leží na přímce $PQ$, čili $|XA|^2 = p(X,\omega)$. Zároveň leží na přímce $BC$, která se dotýká $\omega$ v~bodě $U$, takže $p(X,\omega) = |XU|^2$. Dostáváme
    $$
    |XA| = |XU|.
    $$

    Úsečka $AX$ dělí trojúhelník $ABC$ na trojúhelníky $ABX$ a~$ACX$. Body $U$ i~$X$ leží na úsečce $BC$, takže $X$ leží na úsečce $BU$ nebo na úsečce $UC$. V~prvním případě je $|BX| + |XU| = |BU|$ a~obvod trojúhelníku $ABX$ je
    $$
    |AB| + |BX| + |XA| = c + |BX| + |XU| = c + |BU| = 2.
    $$
    Ve druhém případě je $|CX| + |XU| = |CU|$ a~obvod trojúhelníku $ACX$ je
    $$
    |AC| + |CX| + |XA| = b + |CX| + |XU| = b + |CU| = 2.
    $$
}

\EnvId{hWnbfKi7JKW34d1E7tUxF-kruznice-nad-stranami-ctyruhelniku}
\Problem{0}{}{
    Čtyřúhelník $ABCD$ splňuje $a + b = c + d$, kde $a, b, c, d$ jsou jeho délky stran. Nad každou z jeho stran jako průměrem je sestrojena kružnice. Dokažte, že existuje kružnice dotýkající se všech čtyř z nich.
}{
    Představte si kýženou kouzelnou kružnici a zkuste uhodnout, kde by mohl být její střed. Přemýšlejte, kde jsou středy našich čtyř Thaletových kružnic.
}{
    Střed kružnice bude středem některé úhlopříčky; funguje jen jedna z nich, tak zkuste obě a nevzdávejte to.
}{
    Vzdálenosti z tohoto středu ke středům Thaletových kružnic plynou z věty o střední příčce. Zbývá spočítat kýžený poloměr a ověřit, že naše kouzelná podmínka ho zaručuje.
}{
    Nechť $a = |AB|$, $b = |BC|$, $c = |CD|$, $d = |DA|$ a nechť $P$, $Q$, $R$, $S$ jsou po řadě středy stran $AB$, $BC$, $CD$, $DA$. Thaletovy kružnice ze zadání jsou tedy $\omega_P$ se středem~$P$ a poloměrem $a/2$, $\omega_Q$ se středem~$Q$ a poloměrem $b/2$, $\omega_R$ se středem~$R$ a poloměrem $c/2$ a $\omega_S$ se středem~$S$ a poloměrem $d/2$.

    Střed hledané kružnice musíme z něčeho vyrobit a jediné, co o daných čtyřech kružnicích víme, jsou polohy jejich středů: jsou to středy stran. Když je obrázek plný středů úseček, vyplatí se přidat další střed, aby vznikly střední příčky, a přirozenými kandidáty jsou středy úhlopříček. Úhlopříčka~$AC$ přitom dělí strany na dvojice $\{AB, BC\}$ a $\{CD, DA\}$, tedy přesně tak, jak je seskupuje předpoklad $a + b = c + d$. Zkusme tedy její střed.

    Nechť $O$ je střed úsečky~$AC$. V trojúhelníku~$ABC$ jsou $O$, $P$ středy stran $CA$, $AB$ a $O$, $Q$ středy stran $AC$, $CB$; v trojúhelníku~$ACD$ jsou $O$, $R$ středy stran $CA$, $CD$ a $O$, $S$ středy stran $AC$, $AD$. Z věty o střední příčce tak dostáváme
    $$
    \gather
    |OP| = \tfrac12 |BC| = \tfrac b2, \qquad |OQ| = \tfrac12 |AB| = \tfrac a2, \\
    |OR| = \tfrac12 |DA| = \tfrac d2, \qquad |OS| = \tfrac12 |CD| = \tfrac c2.
    \endgather
    $$
    O tvaru čtyřúhelníku jsme přitom nic nepředpokládali, vše další je už jen práce s těmito čtyřmi vzdálenostmi.

    Podstatné je křížení: vzdálenost bodu~$O$ od středu kružnice nad stranou~$AB$ je polovina délky~$BC$ a naopak. Proto
    $$
    \gather
    |OP| + \tfrac a2 = |OQ| + \tfrac b2 = \tfrac{a+b}2, \\
    |OR| + \tfrac c2 = |OS| + \tfrac d2 = \tfrac{c+d}2,
    \endgather
    $$
    a díky předpokladu $a + b = c + d$ jsou všechny čtyři hodnoty stejné. Označme tuto společnou hodnotu $\rho = (a+b)/2 = (c+d)/2$ (je to čtvrtina obvodu čtyřúhelníku) a nechť $\Omega$ je kružnice se středem~$O$ a poloměrem~$\rho$.

    Zbývá ověřit dotyk. Mají-li dvě kružnice poloměry $\rho > r > 0$ a jejich středy jsou ve vzdálenosti $\rho - r > 0$, dotýkají se vnitřně: menší leží uvnitř větší a jediný společný bod je průsečík polopřímky ze středu větší do středu menší s větší kružnicí. (Vnější dotyk by znamenal vzdálenost středů $\rho + r$, o ten tu nejde.) Předchozí rovnosti říkají právě toto:
    $$
    |OP| = \rho - \tfrac a2 = \tfrac b2 > 0, \qquad |OQ| = \rho - \tfrac b2 = \tfrac a2 > 0,
    $$
    $$
    |OR| = \rho - \tfrac c2 = \tfrac d2 > 0, \qquad |OS| = \rho - \tfrac d2 = \tfrac c2 > 0,
    $$
    přičemž všechny čtyři poloměry $a/2$, $b/2$, $c/2$, $d/2$ jsou kladné a menší než~$\rho$. Kružnice~$\Omega$ se tedy vnitřně dotýká všech čtyř Thaletových kružnic, které leží v jejím vnitřku.

    Na závěr ještě, proč nefunguje druhá úhlopříčka. Je-li $O'$ střed úsečky~$BD$, táž úvaha v trojúhelnících~$ABD$ a~$CBD$ dá $|O'P| = d/2$, $|O'S| = a/2$, $|O'Q| = c/2$ a $|O'R| = b/2$, takže tentokrát
    $$
    \gather
    |O'P| + \tfrac a2 = |O'S| + \tfrac d2 = \tfrac{a+d}2, \\
    |O'Q| + \tfrac b2 = |O'R| + \tfrac c2 = \tfrac{b+c}2.
    \endgather
    $$
    Kružnice se středem~$O'$ dotýkající se vnitřně všech čtyř by tedy vyžadovala $a + d = b + c$, což je jiná podmínka než ta zadaná. Úhlopříčka~$BD$ seskupuje strany do dvojic $\{AB, DA\}$ a $\{BC, CD\}$, a to je seskupení, které nám předpoklad nedává.
}

\EnvId{KWvpK6SAI-wBt88-OpQjm-sestiuhelnik-s-vepsanou-kruznici}
\Problem{0}{}{
    Je-li do šestiúhelníku~$ABCDEF$ vepsána kružnice, pak úhlopříčky $AD$, $BE$, $CF$ procházejí jedním bodem.
}{
    Narýsujte tři kružnice tak, aby úhlopříčky $AD$, $BE$, $CF$ byly chordálami. Je to záludné.
}{
    Jedna z kružnic se dotýká polopřímek opačných k $BA$ a $DE$. Analogicky ostatní. Jediný trik je umístit je tak, aby chordály fungovaly. K tomu rozhodně potřebujeme využít vepsanou kružnici $ABCDEF$.
}{
    Tři přímky procházející jedním bodem volají po potenčním středu, hledáme tedy tři kružnice, jejichž chordálami jsou právě $AD$, $BE$, $CF$; jediné, co v úloze máme, je vepsaná kružnice, a mocnost vrcholu umíme uchopit tehdy, když z něj vede tečna.

    Nechť je $ABCDEF$ konvexní a nechť se jeho vepsaná kružnice~$\omega$ se středem~$O$ a poloměrem~$r$ dotýká stran $AB$, $BC$, $CD$, $DE$, $EF$, $FA$ po řadě v bodech $P$, $Q$, $R$, $S$, $T$, $U$; v tomto pořadí leží tyto body i na~$\omega$, ve stejném směru oběhu jako vrcholy. Délky tečen z vrcholů označme
    $$
    \gather
    a = |AP| = |AU|, \quad b = |BP| = |BQ|, \quad c = |CQ| = |CR|,\\
    d = |DR| = |DS|, \quad e = |ES| = |ET|, \quad f = |FT| = |FU|.
    \endgather
    $$

    Pro bod~$M$ na tečně kružnice~$\omega'$ s bodem dotyku~$Z$ je $p(M,\omega') = |MZ|^2$, kde $p(M,\omega')$ značí mocnost bodu~$M$ ke kružnici~$\omega'$. Nechť se proto kružnice~$\omega_1$ dotýká přímek protilehlých stran $AB$, $DE$ v bodech $P_1$, $S_1$, kružnice~$\omega_2$ přímek $BC$, $EF$ v bodech $Q_1$, $T_1$ a kružnice~$\omega_3$ přímek $CD$, $FA$ v bodech $R_1$, $U_1$. Každý vrchol pak leží na tečně právě dvou z nich, například~$B$ na tečnách $AB$, $BC$ kružnic $\omega_1$, $\omega_2$, a přímka~$BE$ bude chordálou kružnic $\omega_1$, $\omega_2$, jakmile $|BP_1| = |BQ_1|$ a $|ES_1| = |ET_1|$. Chceme tedy, aby v každém vrcholu byly oba jeho body dotyku od něj stejně vzdálené.

    Původní body dotyku to splňují, jenže dávají $\omega_1 = \omega_2 = \omega_3 = \omega$. Posuneme je proto po stranách o tutéž délku, střídavě ve směru oběhu a proti němu: na každé straně leží právě jeden z vrcholů $B$, $D$, $F$ a bod dotyku posuneme směrem k němu, a to až za něj. Zvolme $t > \max(b,d,f)$ a nechť body $P_1$, $Q_1$ leží na polopřímkách opačných k $BA$, $BC$ ve vzdálenosti $t-b$ od~$B$, body $R_1$, $S_1$ na polopřímkách opačných k $DC$, $DE$ ve vzdálenosti $t-d$ od~$D$ a body $T_1$, $U_1$ na polopřímkách opačných k $FE$, $FA$ ve vzdálenosti $t-f$ od~$F$. Pak
    $$
    \gather
    |AP_1| = |AU_1| = a+t, \quad |BP_1| = |BQ_1| = t-b,\\
    |CQ_1| = |CR_1| = c+t, \quad |DR_1| = |DS_1| = t-d,\\
    |ES_1| = |ET_1| = e+t, \quad |FT_1| = |FU_1| = t-f,
    \endgather
    $$
    například $|AP_1| = |AB| + |BP_1| = (a+b)+(t-b)$ a $|CQ_1| = |CB| + |BQ_1| = (b+c)+(t-b)$.

    Zbývá ukázat, že kružnice dotýkající se $AB$ v bodě~$P_1$ a $DE$ v bodě~$S_1$ vůbec existuje; právě tady se střídání směrů uplatní. Nechť~$m_1$ je osa úsečky~$PS$ a~$\sigma$ osová souměrnost podle ní. Z $|OP| = |OS| = r$ prochází~$m_1$ bodem~$O$, takže~$\sigma$ zobrazuje~$\omega$ na sebe, bod~$P$ na bod~$S$, a tedy i tečnu~$AB$ na tečnu~$DE$. Jelikož je šestiúhelník konvexní, leží při oběhu $A$, $B$, $C$, $D$, $E$, $F$ střed~$O$ u každé strany po téže ruce, řekněme vlevo. Souměrnost je nepřímá shodnost a levou stranu vyměňuje za pravou: směr $A \to B$, který má v bodě~$P$ střed~$O$ vlevo, se proto zobrazí na ten směr přímky~$DE$, který má v bodě~$S$ střed~$O$ vpravo, čili na směr $E \to D$. Body $P_1$, $S_1$ přitom vznikly posunem o~$t$ právě v těchto dvou směrech, takže $\sigma(P_1) = S_1$.

    Přímka~$m_1$ není kolmá na~$AB$: jinak by splynula s přímkou~$OP$, jedinou kolmicí na~$AB$ vedenou bodem~$O$, a~$\sigma$ by nechala~$P$ na místě, čili $P = S$. Kolmice na~$AB$ v bodě~$P_1$ proto protne~$m_1$ v jediném bodě~$O_1$; ten~$\sigma$ nechá na místě a tuto kolmici zobrazí na kolmici na~$DE$ v bodě~$S_1$, takže $|O_1P_1| = |O_1S_1|$ a kružnice~$\omega_1$ se středem~$O_1$ a poloměrem~$|O_1P_1|$ je hledaná. Nulový poloměr by znamenal $O_1 = P_1 = S_1$, tedy že~$P_1$ je průsečíkem různých přímek $AB$, $DE$ (různé jsou, protože se~$\omega$ dotýkají v různých bodech); při dané dvojici protilehlých stran to vylučuje nejvýše jednu hodnotu~$t$, dohromady tedy nejvýše tři, a~$t$ volíme mimo ně. Konečně $O_1 \neq O$, protože~$O$ leží na kolmici na~$AB$ v bodě~$P$ a~$O_1$ na rovnoběžné kolmici v bodě $P_1 \neq P$. Stejně sestrojíme~$\omega_2$ se středem~$O_2$ na ose~$m_2$ úsečky~$QT$ a~$\omega_3$ se středem~$O_3$ na ose~$m_3$ úsečky~$RU$.

    Kružnice jsou po dvou nesoustředné. Osy $m_1$, $m_2$, $m_3$ procházejí bodem~$O$ a jsou navzájem různé: body $P$, $Q$, $S$, $T$ leží na~$\omega$ v tomto cyklickém pořadí, takže se tětivy $PS$ a~$QT$ protínají a nejsou rovnoběžné, a tedy nesplývají ani kolmice na ně vedené bodem~$O$, čili~$m_1$ a~$m_2$; analogicky pro dvojice $QT$, $RU$ a $RU$, $PS$. Přímky $m_1$, $m_2$, $m_3$ tak mají jediný společný bod~$O$, a jelikož~$O_i$ leží na~$m_i$ a $O_i \neq O$, jsou středy $O_1$, $O_2$, $O_3$ navzájem různé.

    Vrcholy leží na tečnách příslušných kružnic, takže
    $$
    \gather
    p(B,\omega_1) = (t-b)^2 = p(B,\omega_2),\\
    p(E,\omega_1) = (e+t)^2 = p(E,\omega_2).
    \endgather
    $$
    Různé body $B$, $E$ tedy leží na chordále kružnic $\omega_1$, $\omega_2$, tou je proto přímka~$BE$. Analogicky hodnoty $(c+t)^2$ v~$C$ a $(t-f)^2$ v~$F$ dělají $CF$ chordálou kružnic $\omega_2$, $\omega_3$ a hodnoty $(a+t)^2$ v~$A$ a $(t-d)^2$ v~$D$ dělají $AD$ chordálou kružnic $\omega_1$, $\omega_3$.

    Tři po dvou nesoustředné kružnice mají buď potenční střed, kterým procházejí všechny tři chordály, nebo jejich středy leží na jedné přímce a pak jsou všechny tři chordály na tuto přímku kolmé, tedy rovnoběžné, případně splynou. Druhá možnost odpadá: úhlopříčka~$AD$ dělí konvexní šestiúhelník na $ABCD$ a~$ADEF$, takže $B$ a~$E$ leží v opačných polorovinách určených přímkou~$AD$ a úsečka~$BE$ ji protíná v jediném bodě. Úhlopříčky $AD$, $BE$, $CF$ tedy procházejí jedním bodem, potenčním středem kružnic $\omega_1$, $\omega_2$, $\omega_3$.
}

\DisplayHints

\bye
